\chapter{Coequational Subobjects and Birkhoff Theorems}
\label{ch:coequational-subobjects-birkhoff}

\section{Overview and standing assumptions}
\label{sec:coeq-birkhoff-overview}

\subsection{Orientation}
\label{subsec:coeq-birkhoff-orientation}

The preceding chapter constructed, for a behaviour specification
\[
    k:K\to \mathbf{Beh}_0,
\]
its derivative image
\[
    \operatorname{Der}_0(K)\hookrightarrow \mathbf{Beh}_0
\]
and the corresponding minimal residual Mealy morphism
\[
    \mathsf{Min}(K):\mathbb A\rightsquigarrow\mathbb B.
\]
The present chapter reorganizes residual semantics in coequational form. The
central semantic object is the canonical behaviour Mealy morphism
\[
    \mathbf{Beh}^{\mathrm{can}}_{\mathbb A,\mathbb B}:
    \mathbb A\rightsquigarrow\mathbb B.
\]
By Chapter~\ref{ch:behaviour-categories-myhill-nerode}, this object is
terminal in the ordinary category
\[
    \mathsf{Mly}^{\mathrm{str}}_{\mathbb A,\mathbb B}
\]
of Mealy morphisms and strict semantic homomorphisms. Hence a local
coequational theory is a subobject of this terminal semantic object:
\[
    V\hookrightarrow \mathbf{Beh}^{\mathrm{can}}_{\mathbb A,\mathbb B}.
\]
A Mealy morphism
\[
    P:\mathbb A\rightsquigarrow\mathbb B
\]
models \(V\) when its unique semantic map to the terminal object factors
through \(V\):
\[
\begin{tikzcd}[column sep=large, row sep=large]
    P
        \arrow[rr, "\beta_P"]
        \arrow[dr, dashed, "\bar\beta_P"']
    &&
    \mathbf{Beh}^{\mathrm{can}}_{\mathbb A,\mathbb B}
    \\
    &
    V
        \arrow[ur, hook, "j_V"']
    &
\end{tikzcd}
\]
Equivalently, the carrier map
\[
    \beta_P:P\to\mathbf{Beh}_0
\]
factors through the underlying carrier subobject
\[
    V\hookrightarrow\mathbf{Beh}_0
\]
in the slice over
\[
    Z:=A_0\times B_0.
\]

The basic examples of coequational theories are semantic images. For a Mealy
morphism \(P\), its semantic image is the regular image of its terminal
semantic map:
\[
    P\twoheadrightarrow \operatorname{Beh}(P)
    \hookrightarrow \mathbf{Beh}^{\mathrm{can}}_{\mathbb A,\mathbb B}.
\]
The first main task of the chapter is to prove that regular images of strict
semantic homomorphisms are created by the ambient carrier slice. This makes
\[
    \operatorname{Beh}(P)
\]
again a subobject of the canonical behaviour Mealy morphism.

The chapter proves a coequational Birkhoff theorem. Locally, a replete full
class of Mealy morphisms is definable by a local coequational theory exactly
when it is closed under strict homomorphic preimages, regular-epimorphic
quotients, and small coproducts. The representing coequational theory
associated to a class
\[
    \mathcal C\subseteq\mathsf{Mly}^{\mathrm{str}}_{\mathbb A,\mathbb B}
\]
is the join of the semantic images of its objects:
\[
    V_{\mathcal C}
    :=
    \bigvee_{P\in\mathcal C}\operatorname{Beh}(P)
    \hookrightarrow \mathbf{Beh}_0.
\]
The use of arbitrary joins of subobjects is the point at which this chapter
works in a Grothendieck topos. The resulting statement is in the same lineage
as Birkhoff and Eilenberg type correspondences for recognizable languages and
categorical varieties of languages
\cite{Eilenberg1974,AdamekMiliusMyersUrbat2019TOCL,Salamanca2016Monad,Salamanca2017Unveiling}.

\subsection{Ambient assumptions}
\label{subsec:coeq-birkhoff-ambient-assumptions}

\begin{assumption}
\label{ass:coeq-birkhoff-grothendieck-topos}
    Throughout this chapter, \(\mathcal E\) is a Grothendieck topos with chosen
    pullbacks. Thus \(\mathcal E\) is Barr-exact and locally cartesian closed,
    has stable regular epi--mono factorizations, effective quotients of
    internal equivalence relations, small coproducts, and complete subobject
    lattices in every slice. All coproducts and joins below are small, relative
    to the fixed ambient universe. We use these topos-theoretic properties in
    the form recorded in
    \cite{Johnstone2002,Borceux1994,Jacobs1999CategoricalLogic}.
\end{assumption}

\paragraph{Internal-language convention.}
All arguments below written in generalized-element notation are interpreted
in the Mitchell--Bénabou internal language of \(\mathcal E\). Equivalently,
each such argument may be read after pullback along an arbitrary generalized
element. In particular, assertions of the form
\[
    x\in X,
\]
implications between subobject-membership predicates, and elementwise
equalities are shorthand for the corresponding internal statements in
\(\mathcal E\).

The input and output internal categories are
\[
    \mathbb A=(A_0,A_1,s_A,t_A,e_A,m_A),
    \qquad
    \mathbb B=(B_0,B_1,s_B,t_B,e_B,m_B).
\]
When the input and output categories are fixed, we write
\[
    Z:=A_0\times B_0.
\]
Composition in \(\mathbb A\) is written
\[
    m_A(a,b)=b\circ a
\]
for composable arrows
\[
    a:u\to v,
    \qquad
    b:v\to w.
\]
Thus a Mealy state over \(w\) processes \(b\circ a\) by first processing \(b\),
then processing \(a\).

Small coproducts include the empty coproduct. Thus every coequationally
definable class contains the initial strict Mealy morphism. Dually, the empty
join of local coequational theories is the bottom subobject of the canonical
behaviour object.

\subsection{Strict Mealy morphisms and terminal semantics}
\label{subsec:coeq-standing-notation-terminal-semantics}

Throughout, Mealy morphisms
\[
    P:\mathbb A\rightsquigarrow\mathbb B
\]
are considered as objects of the ordinary category
\[
    \mathsf{Mly}^{\mathrm{str}}_{\mathbb A,\mathbb B}.
\]
A strict semantic homomorphism
\[
    f:P\to Q
\]
is a map in the carrier slice
\[
    \mathcal E/Z
\]
preserving transition and output on the nose. Thus, writing
\[
    P=(P,p_P,q_P,\delta_P,\omega_P),
    \qquad
    Q=(Q,p_Q,q_Q,\delta_Q,\omega_Q),
\]
strictness means
\[
    p_Qf=p_P,
    \qquad
    q_Qf=q_P,
\]
and
\[
    f\delta_P
    =
    \delta_Q(1\times f),
    \qquad
    \omega_P
    =
    \omega_Q(1\times f).
\]
For every Mealy morphism
\[
    P=(P,p_P,q_P,\delta_P,\omega_P),
\]
write
\[
    D_P
    :=
    A_1\times_{t_A,A_0,p_P}P
\]
for its transition domain, and write
\[
    \lambda_P:D_P\to A_1,
    \qquad
    \rho_P:D_P\to P
\]
for the two pullback projections. Thus
\[
    t_A\lambda_P=p_P\rho_P.
\]
The transition and output typing equations may then be written
\[
    p_P\delta_P=s_A\lambda_P,
\]
and
\[
    s_B\omega_P=q_P\delta_P,
    \qquad
    t_B\omega_P=q_P\rho_P.
\]

For a map
\[
    f:P\to Q
\]
over \(A_0\), write
\[
    \widehat f:
    D_P\to D_Q
\]
for the induced map on transition domains, characterized by
\[
    \lambda_Q\widehat f=\lambda_P,
    \qquad
    \rho_Q\widehat f=f\rho_P.
\]
In generalized-element notation,
\[
    \widehat f(a,x)=(a,fx).
\]
Subobjects, coproducts, and full subcategories below are taken in
\[
    \mathsf{Mly}^{\mathrm{str}}_{\mathbb A,\mathbb B}
\]
unless explicitly stated otherwise. Regular quotients are understood in the
carrier-based sense of the following definition.

\begin{definition}[Regular quotient]
\label{def:regular-quotient-strict-mealy}
    A strict semantic homomorphism
    \[
        e:P\to Q
    \]
    is a \textbf{regular quotient} if its underlying carrier map
    \[
        Ue:UP\to UQ
    \]
    is a regular epimorphism in
    \[
        \mathcal E/Z,
        \qquad
        Z=A_0\times B_0.
    \]
    Thus regular-quotient closure throughout this chapter means closure under
    strict semantic homomorphisms whose carriers are regular epimorphisms in
    the carrier slice.
\end{definition}

A replete full subcategory
\[
    \mathcal C\subseteq
    \mathsf{Mly}^{\mathrm{str}}_{\mathbb A,\mathbb B}
\]
is said to be \textbf{closed under strict homomorphic preimages} if, for every
strict semantic homomorphism
\[
    f:P\to Q,
\]
with
\[
    Q\in\mathcal C,
\]
one has
\[
    P\in\mathcal C.
\]
In particular, this includes closure under strict subobjects, since a strict
subobject is a strict semantic homomorphism into its ambient object.

Recall from Chapter~\ref{ch:behaviour-categories-myhill-nerode} that every
Mealy morphism \(P\) has a residual semantics map
\[
    \beta_P:P\to\mathbf{Beh}_0.
\]
For \(x\in P\), the residual behaviour
\[
    \beta_P(x)=H_x:\mathbb A/p_P(x)\to\mathbb B
\]
is given on objects by
\[
    H_x(r)=q_P(\delta_P(r,x)),
\]
and on residual arrows
\[
    w\xrightarrow{a}u\xrightarrow{r}p_P(x)
\]
by
\[
    H_x(a,r)=\omega_P(a,\delta_P(r,x)).
\]
Equivalently,
\[
    \beta_P:P\to\mathbf{Beh}^{\mathrm{can}}_{\mathbb A,\mathbb B}
\]
is the unique strict semantic homomorphism from \(P\) to the terminal object
\[
    \mathbf{Beh}^{\mathrm{can}}_{\mathbb A,\mathbb B}
\]
of \(\mathsf{Mly}^{\mathrm{str}}_{\mathbb A,\mathbb B}\). The use of terminal
semantic objects follows the categorical semantics of realization and
coalgebraic behaviour maps
\cite{Goguen1973,Rutten2000,Jacobs2016Coalgebra}.

\begin{lemma}[Strict maps preserve residual semantics]
\label{lem:coeq-strict-maps-preserve-beta}
    Let
    \[
        f:P\to Q
    \]
    be a strict semantic homomorphism. Then
    \[
        \beta_Qf=\beta_P.
    \]
\end{lemma}

\begin{proof}
    \leavevmode
    \begin{enumerate}
        \item \textbf{Use terminality.}
        Since \(f\) is strict and \(\beta_Q\) is strict, the composite
        \[
            \beta_Qf:P\to\mathbf{Beh}^{\mathrm{can}}_{\mathbb A,\mathbb B}
        \]
        is a strict semantic homomorphism. The map \(\beta_P\) is also a strict
        semantic homomorphism from \(P\) to the terminal object.

        This gives the parallel pair
        \[
        \begin{tikzcd}[column sep=huge, row sep=large]
            P
                \arrow[r, shift left=0.65ex, "\beta_Qf"]
                \arrow[r, shift right=0.65ex, "\beta_P"']
            &
            \mathbf{Beh}^{\mathrm{can}}_{\mathbb A,\mathbb B}.
        \end{tikzcd}
        \]

        \item \textbf{Conclude uniqueness.}
        By terminality of
        \[
            \mathbf{Beh}^{\mathrm{can}}_{\mathbb A,\mathbb B}
        \]
        in \(\mathsf{Mly}^{\mathrm{str}}_{\mathbb A,\mathbb B}\), the two maps
        agree:
        \[
            \beta_Qf=\beta_P.
        \]
    \end{enumerate}
\end{proof}

\begin{lemma}[Pullbacks and monomorphisms are created by carriers]
\label{lem:coeq-carrier-created-pullbacks-monomorphisms}
    The carrier functor
    \[
        U:
        \mathsf{Mly}^{\mathrm{str}}_{\mathbb A,\mathbb B}
        \longrightarrow
        \mathcal E/Z
    \]
    creates pullbacks.

    Consequently, a strict semantic homomorphism
    \[
        f:P\to Q
    \]
    is monic in
    \[
        \mathsf{Mly}^{\mathrm{str}}_{\mathbb A,\mathbb B}
    \]
    if and only if its underlying carrier map is monic in
    \[
        \mathcal E/Z.
    \]
\end{lemma}

\begin{proof}
    Let
    \[
        f:P\to R,
        \qquad
        g:Q\to R
    \]
    be strict semantic homomorphisms. Form their pullback in the carrier
    slice:
    \[
    \begin{tikzcd}[column sep=large, row sep=large]
        X
            \arrow[r, "\pi_Q"]
            \arrow[d, "\pi_P"']
            \arrow[dr, phantom, "\lrcorner", very near start]
        &
        Q
            \arrow[d, "g"]
        \\
        P
            \arrow[r, "f"']
        &
        R .
    \end{tikzcd}
    \]
    Thus
    \[
        X=P\times_RQ
    \]
    in \(\mathcal E/Z\). Put
    \[
        p_X:=p_P\pi_P=p_Q\pi_Q,
        \qquad
        q_X:=q_P\pi_P=q_Q\pi_Q.
    \]

    Write
    \[
        D_X
        :=
        A_1\times_{t_A,A_0,p_X}X
    \]
    and similarly for \(D_P,D_Q,D_R\). The pullback projections induce maps
    \[
        \widehat\pi_P:D_X\to D_P,
        \qquad
        \widehat\pi_Q:D_X\to D_Q.
    \]

    Since \(f\) and \(g\) are strict,
    \[
    \begin{aligned}
        f\delta_P\widehat\pi_P
        &=
        \delta_R\widehat f\,\widehat\pi_P
        \\
        &=
        \delta_R\widehat g\,\widehat\pi_Q
        \\
        &=
        g\delta_Q\widehat\pi_Q.
    \end{aligned}
    \]
    The middle equality follows from
    \[
        f\pi_P=g\pi_Q,
    \]
    which implies
    \[
        \widehat f\,\widehat\pi_P
        =
        \widehat g\,\widehat\pi_Q.
    \]

    Hence the pair
    \[
        \bigl(
            \delta_P\widehat\pi_P,
            \delta_Q\widehat\pi_Q
        \bigr)
    \]
    factors uniquely through the underlying pullback \(X\) in
    \(\mathcal E\). Define
    \[
        \delta_X:D_X\to X
    \]
    by
    \[
        \pi_P\delta_X
        =
        \delta_P\widehat\pi_P,
        \qquad
        \pi_Q\delta_X
        =
        \delta_Q\widehat\pi_Q.
    \]

    The transition has the required input typing:
    \[
    \begin{aligned}
        p_X\delta_X
        &=
        p_P\pi_P\delta_X
        \\
        &=
        p_P\delta_P\widehat\pi_P
        \\
        &=
        s_A\lambda_P\widehat\pi_P
        \\
        &=
        s_A\lambda_X.
    \end{aligned}
    \]

    Strictness of \(f\) and \(g\) also gives
    \[
    \begin{aligned}
        \omega_P\widehat\pi_P
        &=
        \omega_R\widehat f\,\widehat\pi_P
        \\
        &=
        \omega_R\widehat g\,\widehat\pi_Q
        \\
        &=
        \omega_Q\widehat\pi_Q.
    \end{aligned}
    \]
    Define the common map to be
    \[
        \omega_X:D_X\to B_1.
    \]

    The output map has the required source typing:
    \[
    \begin{aligned}
        s_B\omega_X
        &=
        s_B\omega_P\widehat\pi_P
        \\
        &=
        q_P\delta_P\widehat\pi_P
        \\
        &=
        q_P\pi_P\delta_X
        \\
        &=
        q_X\delta_X.
    \end{aligned}
    \]

    It also has the required target typing. Writing
    \[
        \rho_X:D_X\to X
    \]
    for the state projection,
    \[
    \begin{aligned}
        t_B\omega_X
        &=
        t_B\omega_P\widehat\pi_P
        \\
        &=
        q_P\rho_P\widehat\pi_P
        \\
        &=
        q_P\pi_P\rho_X
        \\
        &=
        q_X\rho_X.
    \end{aligned}
    \]

    The transition unit law holds after composing with both pullback
    projections. For example,
    \[
    \begin{aligned}
        \pi_P\delta_X(1_{p_Xx},x)
        &=
        \delta_P(1_{p_P\pi_Px},\pi_Px)
        \\
        &=
        \pi_Px.
    \end{aligned}
    \]
    The analogous equality holds after composing with \(\pi_Q\). Since the
    pullback projections are jointly monic,
    \[
        \delta_X(1_{p_Xx},x)=x.
    \]

    For composable arrows
    \[
        a:u\to v,
        \qquad
        b:v\to w,
    \]
    and \(x\in X\) over \(w\), one has
    \[
    \begin{aligned}
        \pi_P\delta_X(b\circ a,x)
        &=
        \delta_P(b\circ a,\pi_Px)
        \\
        &=
        \delta_P
        \bigl(
            a,
            \delta_P(b,\pi_Px)
        \bigr)
        \\
        &=
        \pi_P\delta_X
        \bigl(
            a,
            \delta_X(b,x)
        \bigr).
    \end{aligned}
    \]
    The same equality holds after composing with \(\pi_Q\). Joint monicity
    therefore gives
    \[
        \delta_X(b\circ a,x)
        =
        \delta_X
        \bigl(
            a,
            \delta_X(b,x)
        \bigr).
    \]

    The output unit law follows by restriction:
    \[
    \begin{aligned}
        \omega_X(1_{p_Xx},x)
        &=
        \omega_P(1_{p_P\pi_Px},\pi_Px)
        \\
        &=
        1_{q_P\pi_Px}
        \\
        &=
        1_{q_Xx}.
    \end{aligned}
    \]

    Likewise,
    \[
    \begin{aligned}
        \omega_X(b\circ a,x)
        &=
        \omega_P(b\circ a,\pi_Px)
        \\
        &=
        \omega_P(b,\pi_Px)
        \circ
        \omega_P
        \bigl(
            a,
            \delta_P(b,\pi_Px)
        \bigr)
        \\
        &=
        \omega_X(b,x)
        \circ
        \omega_X
        \bigl(
            a,
            \delta_X(b,x)
        \bigr).
    \end{aligned}
    \]

    Thus
    \[
        X=(X,p_X,q_X,\delta_X,\omega_X)
    \]
    is a Mealy morphism. By construction, the projections
    \[
        \pi_P:X\to P,
        \qquad
        \pi_Q:X\to Q
    \]
    preserve transition and output strictly.

    Now suppose that
    \[
        h:T\to P,
        \qquad
        k:T\to Q
    \]
    are strict semantic homomorphisms satisfying
    \[
        fh=gk.
    \]
    The carrier pullback gives a unique map
    \[
        \ell:T\to X
    \]
    in \(\mathcal E/Z\) satisfying
    \[
        \pi_P\ell=h,
        \qquad
        \pi_Q\ell=k.
    \]

    To prove transition preservation, compose with the pullback projections:
    \[
    \begin{aligned}
        \pi_P\ell\delta_T
        &=
        h\delta_T
        \\
        &=
        \delta_P\widehat h
        \\
        &=
        \delta_P\widehat\pi_P\,\widehat\ell
        \\
        &=
        \pi_P\delta_X\widehat\ell,
    \end{aligned}
    \]
    and similarly
    \[
        \pi_Q\ell\delta_T
        =
        \pi_Q\delta_X\widehat\ell.
    \]
    Joint monicity gives
    \[
        \ell\delta_T
        =
        \delta_X\widehat\ell.
    \]

    Output preservation follows from
    \[
    \begin{aligned}
        \omega_X\widehat\ell
        &=
        \omega_P\widehat\pi_P\,\widehat\ell
        \\
        &=
        \omega_P\widehat h
        \\
        &=
        \omega_T.
    \end{aligned}
    \]
    Hence \(\ell\) is strict. This proves that \(U\) creates the displayed
    pullback.

    Suppose now that
    \[
        f:P\to Q
    \]
    is a strict semantic homomorphism whose carrier is monic in
    \(\mathcal E/Z\). Since the carrier functor is faithful, \(f\) is monic
    in the strict Mealy category.

    Conversely, suppose that \(f\) is monic in the strict Mealy category.
    Since pullbacks are created by \(U\), the kernel pair of \(f\) is
    \[
    \begin{tikzcd}[column sep=large, row sep=large]
        P\times_QP
            \arrow[r, shift left=0.65ex, "r_1"]
            \arrow[r, shift right=0.65ex, "r_2"']
        &
        P
            \arrow[r, "f"]
        &
        Q .
    \end{tikzcd}
    \]
    Since
    \[
        fr_1=fr_2
    \]
    and \(f\) is monic,
    \[
        r_1=r_2.
    \]
    Therefore the underlying carrier kernel-pair projections agree. Hence
    the carrier of \(f\) is monic in \(\mathcal E/Z\).
\end{proof}

\section{Local coequational theories}
\label{sec:local-coequational-theories}

\subsection{Canonical derivative structure}
\label{subsec:canonical-derivative-structure}

The canonical behaviour object
\[
    \mathbf{Beh}_0
\]
carries the canonical behaviour Mealy morphism
\[
    \mathbf{Beh}^{\mathrm{can}}_{\mathbb A,\mathbb B}:
    \mathbb A\rightsquigarrow\mathbb B.
\]
A state over \(v\in A_0\) is a residual behaviour
\[
    H:\mathbb A/v\to\mathbb B.
\]
For an arrow
\[
    a:u\to v,
\]
the canonical transition is the derivative
\[
    \partial_aH=H\circ a_*,
\]
and the canonical output comparison is
\[
    \omega_{\mathbf{Beh}}(a,H):
    q_{\mathbf{Beh}}(\partial_aH)\to q_{\mathbf{Beh}}(H),
\]
the arrow assigned by \(H\) to the residual morphism
\[
    a\to 1_v
\]
in \(\mathbb A/v\). This derivative notation is compatible with the
Brzozowski derivative viewpoint in the one-object free-monoid specialization
\cite{Brzozowski1964}.

The output comparison is displayed by applying \(H\) to the residual triangle
\[
\begin{tikzcd}[column sep=large, row sep=large]
    u
        \arrow[r, "a"]
        \arrow[rr, "a", bend left=15]
    &
    v
        \arrow[r, "1_v"']
    &
    v .
\end{tikzcd}
\qquad
\leadsto
\qquad
\begin{tikzcd}[column sep=large, row sep=large]
    H(a:u\to v)
        \arrow[r, "{H(a\to 1_v)}"]
    &
    H(1_v).
\end{tikzcd}
\]
Equivalently, the transition-output structure is the map of spans
\[
\begin{tikzcd}[column sep=large, row sep=large]
    A_1\times_{t_A,A_0,p_{\mathbf{Beh}}}\mathbf{Beh}_0
        \arrow[rr, "{\langle \delta_{\mathbf{Beh}},\omega_{\mathbf{Beh}}\rangle}"]
        \arrow[dr, "{s_A\pi_{A_1}}"']
    &&
    \mathbf{Beh}_0\times_{q_{\mathbf{Beh}},B_0,s_B}B_1
        \arrow[dl, "{p_{\mathbf{Beh}}\pi_{\mathbf{Beh}}}"]
        \arrow[dr, "{t_B\pi_{B_1}}"]
    \\
    &
    A_0
    &&
    B_0 .
\end{tikzcd}
\]

Put
\[
    Z=A_0\times B_0
\]
and write
\[
    D_{\mathbf{Beh}}
    :=
    A_1\times_{t_A,A_0,p_{\mathbf{Beh}}}\mathbf{Beh}_0.
\]
There are two maps
\[
    \operatorname{pr}_{\partial}:
    D_{\mathbf{Beh}}\to\mathbf{Beh}_0,
    \qquad
    \partial:
    D_{\mathbf{Beh}}\to\mathbf{Beh}_0,
\]
defined by
\[
    \operatorname{pr}_{\partial}(a,H)=H,
    \qquad
    \partial(a,H)=\partial_aH.
\]

Whenever a derivative domain occurs as a carrier over \(Z\), it is regarded
over \(Z\) by the derivative state. Thus \(D_{\mathbf{Beh}}\) is regarded as
an object of \(\mathcal E/Z\) by the composite
\[
\begin{tikzcd}[column sep=large, row sep=large]
    D_{\mathbf{Beh}}
        \arrow[r, "\partial"]
        \arrow[dr, dashed,
            "{(s_A\pi_{A_1},\,q_{\mathbf{Beh}}\partial)}"']
    &
    \mathbf{Beh}_0
        \arrow[d, "{(p_{\mathbf{Beh}},q_{\mathbf{Beh}})}"]
    \\
    &
    Z .
\end{tikzcd}
\]
Thus, if
\[
    a:u\to v
\]
and \(H\) lies over \(v\), then \((a,H)\) has \(Z\)-coordinate
\[
    \bigl(
        u,
        q_{\mathbf{Beh}}(\partial_aH)
    \bigr).
\]

With this convention, the derivative map
\[
    \partial:D_{\mathbf{Beh}}\to\mathbf{Beh}_0
\]
is a morphism in \(\mathcal E/Z\). The state projection
\[
    \operatorname{pr}_{\partial}:
    D_{\mathbf{Beh}}\to\mathbf{Beh}_0
\]
is not in general a morphism over \(Z\), since it remembers the interface and
anchor coordinates of the original behaviour rather than those of its
derivative.

Pullbacks and universal quantification along
\(\operatorname{pr}_{\partial}\) below are therefore taken in \(\mathcal E\).
Since subobjects of \(\mathbf{Beh}_0\) in \(\mathcal E\) are canonically the
same as subobjects of \(\mathbf{Beh}_0\) in \(\mathcal E/Z\), the resulting
predicates are again regarded over \(Z\).

\subsection{Behavioural predicates and residual invariance}
\label{subsec:behavioural-predicates-residual-invariance}

\begin{definition}[Behavioural predicate]
\label{def:behavioural-predicate}
    A \textbf{behavioural predicate} over \((\mathbb A,\mathbb B)\) is a
    subobject
    \[
        j_C:C\hookrightarrow\mathbf{Beh}_0
    \]
    in \(\mathcal E/Z\). Write
    \[
        \mathsf{Pred}_{\mathbb A,\mathbb B}
        :=
        \operatorname{Sub}_{\mathcal E/Z}(\mathbf{Beh}_0)
    \]
    for the poset of behavioural predicates, ordered by inclusion.
\end{definition}

A behavioural predicate specifies a class of residual behaviours without yet
requiring stability under future input. The residual stability condition is
exactly derivative-closedness.

\begin{definition}[Residual invariance]
\label{def:residual-invariance}
    A behavioural predicate
    \[
        j_V:V\hookrightarrow\mathbf{Beh}_0
    \]
    is \textbf{residually invariant}, or \textbf{derivative-closed}, if the
    canonical derivative map restricts to a morphism
    \[
        \delta_V:
        A_1\times_{t_A,A_0,p_V}V\to V
    \]
    in \(\mathcal E/Z\), equivalently if
    \[
        j_V\delta_V
        =
        \partial(1\times j_V).
    \]
    In generalized-element notation, this says
    \[
        H\in V,
        \quad
        a:u\to p_V(H)
        \quad\Longrightarrow\quad
        \partial_aH\in V.
    \]
\end{definition}

Write
\[
    \mathsf{Coeq}_{\mathbb A,\mathbb B}
    \subseteq
    \mathsf{Pred}_{\mathbb A,\mathbb B}
\]
for the full subposet of residually invariant predicates, and write
\[
    I_\partial:
    \mathsf{Coeq}_{\mathbb A,\mathbb B}
    \hookrightarrow
    \mathsf{Pred}_{\mathbb A,\mathbb B}
\]
for the inclusion.

\subsection{Generated closure, safety core, and the adjoint triple}
\label{subsec:generation-inclusion-safety}

There are two canonical ways to replace an arbitrary behavioural predicate by
a residually invariant one. The first freely adjoins all residual derivatives;
the second retains only those behaviours whose every residual derivative
already satisfies the predicate.

\begin{definition}[Generated residual closure]
\label{def:generated-residual-closure}
    Let
    \[
        j_C:C\hookrightarrow\mathbf{Beh}_0
    \]
    be a behavioural predicate. Its derivative domain is
    \[
        D(C)
        :=
        A_1\times_{t_A,A_0,p_C}C,
    \]
    regarded over \(Z\) by the derivative state. Factor the derivative map in
    \(\mathcal E/Z\) as
    \[
    \begin{tikzcd}[column sep=large, row sep=large]
        D(C)
            \arrow[rr, "{\partial(1\times j_C)}"]
            \arrow[dr, two heads, "e_C"']
        &&
        \mathbf{Beh}_0
        \\
        &
        \operatorname{Gen}_\partial(C)
            \arrow[ur, hook, "g_C"']
        &
    \end{tikzcd}
    \]
    with \(e_C\) regular epimorphic and \(g_C\) monic. The subobject
    \[
        g_C:
        \operatorname{Gen}_\partial(C)
        \hookrightarrow
        \mathbf{Beh}_0
    \]
    is the \textbf{generated residual closure} of \(C\).
\end{definition}

For the specification given by the inclusion
\[
    j_C:C\to\mathbf{Beh}_0,
\]
this is precisely the derivative image constructed in
Chapter~\ref{ch:behaviour-categories-myhill-nerode}. The derivative-closure
theorem,
Theorem~\ref{thm:derivative-closure}, therefore gives the following universal
property.

\begin{proposition}[Universal property of generated residual closure]
\label{prop:generated-residual-closure}
    For every behavioural predicate \(C\), the subobject
    \[
        \operatorname{Gen}_\partial(C)
        \hookrightarrow
        \mathbf{Beh}_0
    \]
    is residually invariant and is the least residually invariant predicate
    containing \(C\). Consequently, for every
    \(V\in\mathsf{Coeq}_{\mathbb A,\mathbb B}\),
    \[
        \operatorname{Gen}_\partial(C)\leq V
        \quad\Longleftrightarrow\quad
        C\leq I_\partial(V).
    \]
\end{proposition}
\begin{proof}
    Apply Theorem~\ref{thm:derivative-closure} to the behaviour specification
    \[
        j_C:C\to\mathbf{Beh}_0.
    \]
    Its formal derivative orbit is
    \[
        A_1\times_{t_A,A_0,p_C}C,
    \]
    and its residual evaluation map is
    \[
        \partial(1\times j_C):
        A_1\times_{t_A,A_0,p_C}C
        \to
        \mathbf{Beh}_0.
    \]
    By definition, the regular image of this map is
    \[
        \operatorname{Gen}_\partial(C)
        \hookrightarrow
        \mathbf{Beh}_0.
    \]

    Theorem~\ref{thm:derivative-closure} states that this regular image is
    residually invariant and is the smallest residually invariant subobject
    containing the image of \(j_C\). Since \(j_C\) is itself a monomorphism,
    its image is exactly the predicate \(C\). Hence
    \[
        \operatorname{Gen}_\partial(C)
    \]
    is the least residually invariant predicate containing \(C\).

    Therefore, for every residually invariant predicate \(V\),
    \[
        \operatorname{Gen}_\partial(C)\leq V
    \]
    if and only if
    \[
        C\leq V.
    \]
    Equivalently,
    \[
        \operatorname{Gen}_\partial(C)\leq V
        \quad\Longleftrightarrow\quad
        C\leq I_\partial(V).
    \]
\end{proof}

The second construction uses universal quantification over all derivatives of
a behaviour. Since \(\mathcal E\) is locally cartesian closed, pullback along
the state projection
\[
    \operatorname{pr}_{\partial}:
    D_{\mathbf{Beh}}\to\mathbf{Beh}_0
\]
has a right adjoint on subobject posets,
\[
    \forall_{\operatorname{pr}_{\partial}}:
    \operatorname{Sub}_{\mathcal E}(D_{\mathbf{Beh}})
    \to
    \operatorname{Sub}_{\mathcal E}(\mathbf{Beh}_0).
\]

\begin{definition}[Residual safety core]
\label{def:residual-safety-core}
    Let
    \[
        j_C:C\hookrightarrow\mathbf{Beh}_0
    \]
    be a behavioural predicate. Its \textbf{residual safety core} is
    \[
        \operatorname{Safe}_\partial(C)
        :=
        \forall_{\operatorname{pr}_{\partial}}
        \bigl(
            \partial^*C
        \bigr)
        \hookrightarrow
        \mathbf{Beh}_0,
    \]
    where
    \[
        \partial^*C\hookrightarrow D_{\mathbf{Beh}}
    \]
    is the pullback of
    \[
        C\hookrightarrow\mathbf{Beh}_0
    \]
    along the derivative map \(\partial\).
\end{definition}

Internally, the defining predicate is
\[
    H\in\operatorname{Safe}_\partial(C)
    \quad\Longleftrightarrow\quad
    \forall a:u\to p_{\mathbf{Beh}}(H),\;
    \partial_aH\in C.
\]
Thus \(\operatorname{Safe}_\partial(C)\) consists of those behaviours for
which every admissible residual continuation satisfies \(C\).

\begin{proposition}[Universal property of the residual safety core]
\label{prop:residual-safety-core}
    For every behavioural predicate \(C\), the residual safety core
    \[
        \operatorname{Safe}_\partial(C)
        \hookrightarrow
        \mathbf{Beh}_0
    \]
    is residually invariant and is the greatest residually invariant
    predicate contained in \(C\).

    Consequently, for every
    \[
        V\in\mathsf{Coeq}_{\mathbb A,\mathbb B},
    \]
    \[
        I_\partial(V)\leq C
        \quad\Longleftrightarrow\quad
        V\leq\operatorname{Safe}_\partial(C).
    \]
\end{proposition}

\begin{proof}
    \leavevmode
    \begin{enumerate}
        \item \textbf{The safety core lies in the original predicate.}

        The identity map of the input category gives a section
        \[
            \iota_{\mathbf{Beh}}:
            \mathbf{Beh}_0\to D_{\mathbf{Beh}},
            \qquad
            H\longmapsto
            \bigl(
                1_{p_{\mathbf{Beh}}(H)},
                H
            \bigr),
        \]
        satisfying
        \[
            \operatorname{pr}_{\partial}
            \iota_{\mathbf{Beh}}
            =
            1_{\mathbf{Beh}_0},
        \]
        and
        \[
            \partial\iota_{\mathbf{Beh}}
            =
            1_{\mathbf{Beh}_0}.
        \]

        Hence, if every derivative of \(H\) lies in \(C\), then in particular
        the identity derivative of \(H\) lies in \(C\). Therefore
        \[
            \operatorname{Safe}_\partial(C)\leq C.
        \]

        \item \textbf{The safety core is derivative-closed.}

        Let
        \[
            H\in\operatorname{Safe}_\partial(C),
            \qquad
            a:u\to p_{\mathbf{Beh}}(H).
        \]
        To show that
        \[
            \partial_aH
            \in
            \operatorname{Safe}_\partial(C),
        \]
        let
        \[
            c:w\to u.
        \]
        The derivative multiplication law gives
        \[
            \partial_c(\partial_aH)
            =
            \partial_{a\circ c}H.
        \]
        Since \(H\) lies in the safety core, every derivative of \(H\),
        including \(\partial_{a\circ c}H\), lies in \(C\). Thus every
        derivative of \(\partial_aH\) lies in \(C\), and therefore
        \[
            \partial_aH
            \in
            \operatorname{Safe}_\partial(C).
        \]

        Hence \(\operatorname{Safe}_\partial(C)\) is residually invariant.

        \item \textbf{Prove maximality.}

        Let \(V\) be residually invariant and suppose
        \[
            V\leq C.
        \]
        If \(H\in V\), residual invariance gives
        \[
            \partial_aH\in V
        \]
        for every admissible input arrow \(a\). Since \(V\leq C\), every such
        derivative lies in \(C\). Hence
        \[
            H\in\operatorname{Safe}_\partial(C).
        \]
        Therefore
        \[
            V\leq\operatorname{Safe}_\partial(C).
        \]

        Conversely, if
        \[
            V\leq\operatorname{Safe}_\partial(C),
        \]
        then
        \[
            V\leq C
        \]
        by the first part.

        This proves
        \[
            I_\partial(V)\leq C
            \quad\Longleftrightarrow\quad
            V\leq\operatorname{Safe}_\partial(C),
        \]
        and establishes the claimed maximality.
    \end{enumerate}
\end{proof}

\begin{theorem}[Generation--inclusion--safety adjoint triple]
\label{thm:generation-inclusion-safety-adjoint-triple}
    The generated residual closure and residual safety core define monotone
    maps
    \[
        \operatorname{Gen}_\partial,
        \operatorname{Safe}_\partial:
        \mathsf{Pred}_{\mathbb A,\mathbb B}
        \to
        \mathsf{Coeq}_{\mathbb A,\mathbb B},
    \]
    and there is an adjoint triple of posets
    \[
        \operatorname{Gen}_\partial
        \dashv
        I_\partial
        \dashv
        \operatorname{Safe}_\partial.
    \]
    Explicitly, for every behavioural predicate \(C\) and every residually
    invariant predicate \(V\),
    \[
        \operatorname{Gen}_\partial(C)\leq V
        \quad\Longleftrightarrow\quad
        C\leq I_\partial(V),
    \]
    and
    \[
        I_\partial(V)\leq C
        \quad\Longleftrightarrow\quad
        V\leq\operatorname{Safe}_\partial(C).
    \]
\end{theorem}

\begin{proof}
    The first adjunction is
    Proposition~\ref{prop:generated-residual-closure}, and the second is
    Proposition~\ref{prop:residual-safety-core}.
\end{proof}

\begin{corollary}[Closure and interior]
\label{cor:residual-closure-and-interior}
    For every behavioural predicate \(C\),
    \[
        \operatorname{Safe}_\partial(C)
        \leq
        C
        \leq
        \operatorname{Gen}_\partial(C).
    \]
    Moreover, the following are equivalent:
    \begin{enumerate}
        \item \(C\) is residually invariant;
        \item \(\operatorname{Gen}_\partial(C)=C\);
        \item \(\operatorname{Safe}_\partial(C)=C\).
    \end{enumerate}
    Thus residually invariant predicates are simultaneously the fixed points
    of the generated residual closure and of the residual safety interior.
\end{corollary}

\begin{proof}
    The two inequalities are the unit of
    \(\operatorname{Gen}_\partial\dashv I_\partial\) and the counit of
    \(I_\partial\dashv\operatorname{Safe}_\partial\). If \(C\) is residually
    invariant, both universal properties force equality. Conversely, either
    equality exhibits \(C\) as an object of
    \(\mathsf{Coeq}_{\mathbb A,\mathbb B}\).
\end{proof}

\subsection{Coequational subobjects}
\label{subsec:coequational-subobjects}

\begin{definition}[Local coequational theory and residual behavioural contract]
\label{def:local-coequational-theory}
    A \textbf{local coequational theory} over
    \((\mathbb A,\mathbb B)\), equivalently a \textbf{residual behavioural
    contract}, is a residually invariant behavioural predicate
    \[
        j_V:V\hookrightarrow\mathbf{Beh}_0.
    \]
\end{definition}

The terminology records two aspects of the same structure. The phrase
``coequational theory'' emphasizes the Birkhoff and Eilenberg developments of
this chapter and the next. The phrase ``residual behavioural contract''
emphasizes that \(V\) is a behavioural requirement stable under every future
typed input. The following proposition identifies these predicates with
strict subobjects of the canonical behaviour machine.

\begin{proposition}[Residual contracts as subobjects of the canonical behaviour machine]
\label{prop:subobjects-canonical-behaviour-machine}
    To give a local coequational theory
    \[
        j_V:V\hookrightarrow\mathbf{Beh}_0
    \]
    is equivalently to give a subobject
    \[
        V\hookrightarrow
        \mathbf{Beh}^{\mathrm{can}}_{\mathbb A,\mathbb B}
    \]
    in
    \[
        \mathsf{Mly}^{\mathrm{str}}_{\mathbb A,\mathbb B}.
    \]

    Under this equivalence, the canonical derivative and output maps restrict
    to
    \[
        \delta_V:
        A_1\times_{t_A,A_0,p_V}V\to V
    \]
    and
    \[
        \omega_V:
        A_1\times_{t_A,A_0,p_V}V\to B_1.
    \]
\end{proposition}

\begin{proof}
    Suppose first that
    \[
        j_V:
        V\hookrightarrow
        \mathbf{Beh}^{\mathrm{can}}_{\mathbb A,\mathbb B}
    \]
    is a subobject in
    \[
        \mathsf{Mly}^{\mathrm{str}}_{\mathbb A,\mathbb B}.
    \]
    Choose a monic representative \(j_V\). By
    Lemma~\ref{lem:coeq-carrier-created-pullbacks-monomorphisms}, its
    underlying carrier map
    \[
        j_V:V\hookrightarrow\mathbf{Beh}_0
    \]
    is monic in \(\mathcal E/Z\).

    Since \(j_V\) is strict, its transition equation is
    \[
        j_V\delta_V
        =
        \delta_{\mathbf{Beh}}(1\times j_V)
        =
        \partial(1\times j_V).
    \]
    Hence the underlying carrier subobject is derivative-closed. Therefore it
    is a local coequational theory.

    Conversely, suppose
    \[
        j_V:V\hookrightarrow\mathbf{Beh}_0
    \]
    is a derivative-closed subobject in \(\mathcal E/Z\). Write
    \[
        p_V:=p_{\mathbf{Beh}}j_V,
        \qquad
        q_V:=q_{\mathbf{Beh}}j_V.
    \]

    Derivative-closedness supplies a unique map
    \[
        \delta_V:
        A_1\times_{t_A,A_0,p_V}V
        \to
        V
    \]
    satisfying
    \[
        j_V\delta_V
        =
        \delta_{\mathbf{Beh}}(1\times j_V).
    \]
    Define
    \[
        \omega_V
        :=
        \omega_{\mathbf{Beh}}(1\times j_V).
    \]

    The output typing equations are inherited from the canonical behaviour
    machine:
    \[
    \begin{aligned}
        s_B\omega_V
        &=
        s_B\omega_{\mathbf{Beh}}\widehat j_V
        \\
        &=
        q_{\mathbf{Beh}}
        \delta_{\mathbf{Beh}}
        \widehat j_V
        \\
        &=
        q_{\mathbf{Beh}}j_V\delta_V
        \\
        &=
        q_V\delta_V.
    \end{aligned}
    \]
    and
    \[
    \begin{aligned}
        t_B\omega_V
        &=
        t_B\omega_{\mathbf{Beh}}\widehat j_V
        \\
        &=
        q_{\mathbf{Beh}}
        \rho_{\mathbf{Beh}}
        \widehat j_V
        \\
        &=
        q_{\mathbf{Beh}}j_V\rho_V
        \\
        &=
        q_V\rho_V.
    \end{aligned}
    \]
    
    The transition unit law follows from
    \[
    \begin{aligned}
        j_V\delta_V(1_{p_VH},H)
        &=
        \delta_{\mathbf{Beh}}
        (1_{p_{\mathbf{Beh}}j_VH},j_VH)
        \\
        &=
        j_VH.
    \end{aligned}
    \]
    Since \(j_V\) is monic,
    \[
        \delta_V(1_{p_VH},H)=H.
    \]

    For composable arrows
    \[
        a:u\to v,
        \qquad
        b:v\to w,
    \]
    and \(H\in V\) over \(w\),
    \[
    \begin{aligned}
        j_V\delta_V(b\circ a,H)
        &=
        \delta_{\mathbf{Beh}}(b\circ a,j_VH)
        \\
        &=
        \delta_{\mathbf{Beh}}
        \bigl(
            a,
            \delta_{\mathbf{Beh}}(b,j_VH)
        \bigr)
        \\
        &=
        \delta_{\mathbf{Beh}}
        \bigl(
            a,
            j_V\delta_V(b,H)
        \bigr)
        \\
        &=
        j_V\delta_V
        \bigl(
            a,
            \delta_V(b,H)
        \bigr).
    \end{aligned}
    \]
    Monicity gives
    \[
        \delta_V(b\circ a,H)
        =
        \delta_V
        \bigl(
            a,
            \delta_V(b,H)
        \bigr).
    \]

    The output unit law is obtained by direct restriction:
    \[
    \begin{aligned}
        \omega_V(1_{p_VH},H)
        &=
        \omega_{\mathbf{Beh}}
        (1_{p_{\mathbf{Beh}}j_VH},j_VH)
        \\
        &=
        1_{q_{\mathbf{Beh}}j_VH}
        \\
        &=
        1_{q_VH}.
    \end{aligned}
    \]

    The output multiplication law is likewise the restriction of the
    canonical output law:
    \[
    \begin{aligned}
        \omega_V(b\circ a,H)
        &=
        \omega_{\mathbf{Beh}}(b\circ a,j_VH)
        \\
        &=
        \omega_{\mathbf{Beh}}(b,j_VH)
        \circ
        \omega_{\mathbf{Beh}}
        \bigl(
            a,
            \delta_{\mathbf{Beh}}(b,j_VH)
        \bigr)
        \\
        &=
        \omega_V(b,H)
        \circ
        \omega_{\mathbf{Beh}}
        \bigl(
            a,
            j_V\delta_V(b,H)
        \bigr)
        \\
        &=
        \omega_V(b,H)
        \circ
        \omega_V
        \bigl(
            a,
            \delta_V(b,H)
        \bigr).
    \end{aligned}
    \]

    Thus \(V\) carries a Mealy structure. By construction,
    \[
        j_V:
        V\to
        \mathbf{Beh}^{\mathrm{can}}_{\mathbb A,\mathbb B}
    \]
    preserves transition and output strictly. Its carrier is monic, so
    Lemma~\ref{lem:coeq-carrier-created-pullbacks-monomorphisms} implies that
    it is monic in the strict Mealy category. Hence it is a subobject of the
    canonical behaviour machine.
\end{proof}

The derivative-closedness condition is the factorization
\[
\begin{tikzcd}[column sep=large, row sep=large]
    A_1\times_{t_A,A_0,p_V}V
        \arrow[rr, "{\partial(1\times j_V)}"]
        \arrow[dr, "\delta_V"']
    &&
    \mathbf{Beh}_0
    \\
    &
    V
        \arrow[ur, hook, "j_V"']
    &
\end{tikzcd}
\]
in \(\mathcal E/Z\). Equivalently, the transition-output cell of the canonical
machine restricts along \(j_V\):
\[
\begin{tikzcd}[column sep=huge, row sep=large]
    A_1\times_{A_0}V
        \arrow[r, "{1\times j_V}"]
        \arrow[d, "{\langle\delta_V,\omega_V\rangle}"']
    &
    A_1\times_{A_0}\mathbf{Beh}_0
        \arrow[d, "{\langle\delta_{\mathbf{Beh}},\omega_{\mathbf{Beh}}\rangle}"]
    \\
    V\times_{B_0}B_1
        \arrow[r, hook, "{j_V\times 1}"']
    &
    \mathbf{Beh}_0\times_{B_0}B_1 .
\end{tikzcd}
\]

When \(V\hookrightarrow\mathbf{Beh}_0\) is a local coequational theory, we
write
\[
    \delta_V(a,H)=\partial_aH
\]
and
\[
    \omega_V(a,H)=\omega_{\mathbf{Beh}}(a,H).
\]

\subsection{Models of local coequational theories}
\label{subsec:models-local-coequational-theories}

\begin{definition}[Models and contract satisfaction]
\label{def:models-coequational-theory}
    Let
    \[
        j_V:
        V\hookrightarrow
        \mathbf{Beh}^{\mathrm{can}}_{\mathbb A,\mathbb B}
    \]
    be a local coequational theory, and regard its carrier as the monomorphism
    \[
        j_V:V\hookrightarrow\mathbf{Beh}_0
    \]
    in
    \[
        \mathcal E/Z,
        \qquad
        Z=A_0\times B_0.
    \]

    A Mealy morphism
    \[
        P:\mathbb A\rightsquigarrow\mathbb B
    \]
    \textbf{models} \(V\), or \textbf{satisfies the residual behavioural
    contract} \(V\), if its carrier semantic map
    \[
        \beta_P:P\to\mathbf{Beh}_0
    \]
    factors through \(j_V\) in the carrier slice \(\mathcal E/Z\).
    Equivalently, there is a necessarily unique morphism
    \[
        \bar\beta_P:P\to V
    \]
    in \(\mathcal E/Z\) such that
    \[
        j_V\bar\beta_P=\beta_P.
    \]

    We write
    \[
        P\models V
    \]
    and denote by
    \[
        \mathsf{Mod}(V)
    \]
    the full subcategory of
    \[
        \mathsf{Mly}^{\mathrm{str}}_{\mathbb A,\mathbb B}
    \]
    on the objects that model \(V\).
\end{definition}

\begin{lemma}[Model factorizations are strict]
\label{lem:model-factorizations-strict}
    If \(P\models V\), then the unique factorization
    \[
        \bar\beta_P:P\to V
    \]
    satisfying
    \[
        j_V\bar\beta_P=\beta_P
    \]
    is a strict semantic homomorphism.
\end{lemma}

\begin{proof}
    \leavevmode
    \begin{enumerate}
        \item \textbf{Check the map over \(Z\).}
        The equation
        \[
            j_V\bar\beta_P=\beta_P
        \]
        and the fact that both \(j_V\) and \(\beta_P\) are maps over \(Z\)
        imply that \(\bar\beta_P\) is a map over \(Z\). This is the lower
        triangle in
        \[
        \begin{tikzcd}[column sep=large, row sep=large]
            P
                \arrow[rr, "\beta_P"]
                \arrow[dr, "\bar\beta_P"']
                \arrow[ddr, "{(p_P,q_P)}"', bend right=18]
            &&
            \mathbf{Beh}_0
                \arrow[ddl, "{(p_{\mathbf{Beh}},q_{\mathbf{Beh}})}", bend left=18]
            \\
            &
            V
                \arrow[ur, hook, "j_V"']
                \arrow[d, "{(p_V,q_V)}"]
            &
            \\
            &
            Z .
        \end{tikzcd}
        \]

        \item \textbf{Check transition preservation.}
        Both composites
        \[
            j_V\bar\beta_P\delta_P
            \qquad\text{and}\qquad
            j_V\delta_V(1\times\bar\beta_P)
        \]
        are equal to the transition-preservation composite for the strict map
        \[
            \beta_P:P\to\mathbf{Beh}^{\mathrm{can}}_{\mathbb A,\mathbb B}.
        \]
        The comparison is the outer rectangle
        \[
        \begin{tikzcd}[column sep=huge, row sep=large]
            A_1\times_{A_0}P
                \arrow[r, "{1\times \bar\beta_P}"]
                \arrow[d, "\delta_P"']
                \arrow[rr, bend left=16, "{1\times \beta_P}"]
            &
            A_1\times_{A_0}V
                \arrow[r, "{1\times j_V}"]
                \arrow[d, "\delta_V"]
            &
            A_1\times_{A_0}\mathbf{Beh}_0
                \arrow[d, "\delta_{\mathbf{Beh}}"]
            \\
            P
                \arrow[r, "\bar\beta_P"']
                \arrow[rr, bend right=16, "\beta_P"']
            &
            V
                \arrow[r, hook, "j_V"']
            &
            \mathbf{Beh}_0 .
        \end{tikzcd}
        \]
        Since \(j_V\) is monic, the transition equation holds in \(V\).

        \item \textbf{Check output preservation.}
        The output equation for \(\bar\beta_P\) follows from the output equation
        for \(\beta_P\), since \(\omega_V\) is the restriction of
        \(\omega_{\mathbf{Beh}}\). It is represented by
        \[
        \begin{tikzcd}[column sep=huge, row sep=large]
            A_1\times_{A_0}P
                \arrow[r, "{1\times \bar\beta_P}"]
                \arrow[dr, "\omega_P"']
            &
            A_1\times_{A_0}V
                \arrow[d, "\omega_V"]
                \arrow[r, "{1\times j_V}"]
            &
            A_1\times_{A_0}\mathbf{Beh}_0
                \arrow[dl, "\omega_{\mathbf{Beh}}"]
            \\
            &
            B_1 .
        \end{tikzcd}
        \]
    \end{enumerate}
\end{proof}

\subsection{Restricted canonical subobjects}
\label{subsec:restricted-canonical-subobjects}

\begin{lemma}[Semantic map of a restricted canonical subobject]
\label{lem:semantic-map-of-a-restricted-canonical-subobject}
    Let
    \[
        V\hookrightarrow\mathbf{Beh}^{\mathrm{can}}_{\mathbb A,\mathbb B}
    \]
    be a local coequational theory. Then the residual semantics map
    \[
        \beta_V:V\to\mathbf{Beh}_0
    \]
    is the inclusion
    \[
        j_V:V\hookrightarrow\mathbf{Beh}_0.
    \]
\end{lemma}

\begin{proof}
    \leavevmode
    \begin{enumerate}
        \item \textbf{Compute residual objects.}
        Let
        \[
            H\in V_v
        \]
        and let
        \[
            r:u\to v
        \]
        be a residual object. In the restricted canonical object, execution of
        \(r\) from \(H\) gives
        \[
            \partial_rH.
        \]
        Its anchor output is
        \[
            q_V(\partial_rH)=H(r).
        \]
        This is the object part of the equality displayed by
        \[
        \begin{tikzcd}[column sep=large, row sep=large]
            H
                \arrow[r, maps to]
                \arrow[d, "{\partial_r}"']
            &
            H(1_v)
            \\
            \partial_rH
                \arrow[r, maps to]
            &
            H(r).
        \end{tikzcd}
        \]

        \item \textbf{Compute residual arrows.}
        For a residual arrow
        \[
            w\xrightarrow{a}u\xrightarrow{r}v,
        \]
        the output produced by the restricted canonical object at
        \(\partial_rH\) along \(a\) is
        \[
            \omega_V(a,\partial_rH).
        \]
        This is the arrow assigned by \(H\) to the residual morphism
        \[
            r\circ a\to r.
        \]
        The comparison is
        \[
        \begin{tikzcd}[column sep=large, row sep=large]
            w
                \arrow[r, "a"]
                \arrow[rr, "r\circ a", bend left=15]
            &
            u
                \arrow[r, "r"']
            &
            v
        \end{tikzcd}
        \qquad
        \leadsto
        \qquad
        \begin{tikzcd}[column sep=large, row sep=large]
            H(r\circ a)
                \arrow[r, "{H(r\circ a\to r)}"]
            &
            H(r).
        \end{tikzcd}
        \]

        \item \textbf{Conclude equality of represented behaviours.}
        Thus the residual behaviour computed from \(H\) in the restricted
        canonical object agrees with \(H\) on residual objects and residual
        arrows. Hence
        \[
            \beta_V(H)=H.
        \]
        Therefore
        \[
            \beta_V=j_V.
        \]
        Equivalently, the terminal semantic triangle is
        \[
        \begin{tikzcd}[column sep=large, row sep=large]
            V
                \arrow[rr, "\beta_V"]
                \arrow[dr, equal]
            &&
            \mathbf{Beh}^{\mathrm{can}}
            \\
            &
            V
                \arrow[ur, hook, "j_V"']
            &
        \end{tikzcd}
        \]
        and the diagonal identity is forced by the previous computation.
    \end{enumerate}
\end{proof}

\section{Semantic images and regular-image structure}
\label{sec:regular-semantic-images}

\subsection{Regular-epimorphic descent for subobjects}
\label{subsec:regular-epi-descent-subobjects}

The following elementary descent fact will be used repeatedly.

\begin{lemma}[Regular-epimorphic descent for subobject membership]
\label{lem:regular-epi-descent-for-subobjects-membership}
    Let \(\mathcal X\) be either \(\mathcal E\) or a slice
    \(\mathcal E/S\). Let
    \[
        e:X\twoheadrightarrow Y
    \]
    be a regular epimorphism in \(\mathcal X\), and let
    \[
        j:V\hookrightarrow W
    \]
    be a monomorphism in \(\mathcal X\). If
    \[
        h:Y\to W
    \]
    is a morphism in \(\mathcal X\) such that
    \[
        he
    \]
    factors through \(j\), then \(h\) factors uniquely through \(j\).
\end{lemma}

\begin{proof}
    \leavevmode
    \begin{enumerate}
        \item \textbf{Use the kernel pair of the regular epimorphism.}
        Let
        \[
            R\rightrightarrows X
        \]
        be the kernel pair of \(e\) in \(\mathcal X\). Since \(e\) is a
        regular epimorphism, it is the coequalizer of this kernel pair:
        \[
        \begin{tikzcd}[column sep=large, row sep=large]
            R
                \arrow[r, shift left=0.65ex, "r_1"]
                \arrow[r, shift right=0.65ex, "r_2"']
            &
            X
                \arrow[r, two heads, "e"]
            &
            Y .
        \end{tikzcd}
        \]

        \item \textbf{Descend the factorization.}
        Suppose
        \[
            he=jg
        \]
        for some \(g:X\to V\). Since
        \[
            he r_1=he r_2,
        \]
        we have
        \[
            jg r_1=jg r_2.
        \]
        Since \(j\) is monic,
        \[
            g r_1=g r_2.
        \]
        Therefore \(g\) descends uniquely through \(e\) to a map
        \[
            \bar g:Y\to V
        \]
        with
        \[
            \bar g e=g.
        \]
        The descent diagram is
        \[
        \begin{tikzcd}[column sep=large, row sep=large]
            R
                \arrow[r, shift left=0.65ex, "r_1"]
                \arrow[r, shift right=0.65ex, "r_2"']
            &
            X
                \arrow[r, two heads, "e"]
                \arrow[dr, "g"']
            &
            Y
                \arrow[d, dashed, "\bar g"]
            \\
            &&
            V .
        \end{tikzcd}
        \]

        \item \textbf{Compare with \(h\).}
        Now
        \[
            j\bar g e=jg=he.
        \]
        Since \(e\) is epic,
        \[
            j\bar g=h.
        \]
        Uniqueness follows from monicity of \(j\). The final comparison is
        \[
        \begin{tikzcd}[column sep=large, row sep=large]
            X
                \arrow[r, two heads, "e"]
                \arrow[dr, "g"']
                \arrow[ddr, "he"', bend right=18]
            &
            Y
                \arrow[d, dashed, "\bar g"]
                \arrow[dd, "h", bend left=18]
            \\
            &
            V
                \arrow[d, hook, "j"]
            \\
            &
            W .
        \end{tikzcd}
        \]
    \end{enumerate}
\end{proof}

\subsection{Regular images of strict homomorphisms}
\label{subsec:regular-images-strict-homomorphisms}

This section proves that regular images of strict semantic homomorphisms are
created by the carrier slice.

The following proposition should be read as a carrier-creation statement for
regular images. In particular, the regular image factorization of a strict
semantic homomorphism is computed in the carrier slice
\[
    \mathcal E/Z
\]
and then carries the unique Mealy structure making the two factorization maps
strict. Thus the phrase ``regular image'' below always refers to the regular
image in the carrier slice, equipped with this created strict Mealy structure.

\begin{proposition}[Regular images are created by the carrier slice]
\label{prop:regular-image-factorization-mealy}
    Let
    \[
        f:P\to Q
    \]
    be a strict semantic homomorphism. Factor its underlying carrier map in
    \(\mathcal E/Z\) as
    \[
        P\xrightarrow{e}I\xrightarrow{m}Q,
    \]
    where \(e\) is a regular epimorphism and \(m\) is a monomorphism.

    Then \(I\) carries a unique Mealy structure such that
    \[
        e:P\twoheadrightarrow I
    \]
    and
    \[
        m:I\hookrightarrow Q
    \]
    are strict semantic homomorphisms. Moreover:

    \begin{enumerate}
        \item \(e\) is a regular epimorphism in
        \[
            \mathsf{Mly}^{\mathrm{str}}_{\mathbb A,\mathbb B};
        \]

        \item \(m\) is a monomorphism in
        \[
            \mathsf{Mly}^{\mathrm{str}}_{\mathbb A,\mathbb B};
        \]

        \item the factorization
        \[
            f=me
        \]
        is the regular image factorization of \(f\) in the strict Mealy
        category.
    \end{enumerate}

    Consequently the carrier functor
    \[
        U:
        \mathsf{Mly}^{\mathrm{str}}_{\mathbb A,\mathbb B}
        \to
        \mathcal E/Z
    \]
    creates regular image factorizations.
\end{proposition}

\begin{proof}
    Put
    \[
        p_I:=p_Qm,
        \qquad
        q_I:=q_Qm,
    \]
    and write
    \[
        D_I
        :=
        A_1\times_{t_A,A_0,p_I}I.
    \]

    Since regular epimorphisms are stable under pullback, the induced map
    \[
        \widehat e:D_P\twoheadrightarrow D_I
    \]
    is a regular epimorphism. In generalized-element notation,
    \[
        \widehat e(a,x)=(a,ex).
    \]

    Consider the underlying map in \(\mathcal E\)
    \[
        h_\delta
        :=
        \delta_Q\widehat m:
        D_I\to Q.
    \]
    Since \(f=me\) and \(f\) is strict,
    \[
    \begin{aligned}
        h_\delta\widehat e
        &=
        \delta_Q\widehat m\,\widehat e
        \\
        &=
        \delta_Q\widehat f
        \\
        &=
        f\delta_P
        \\
        &=
        me\delta_P.
    \end{aligned}
    \]
    Thus \(h_\delta\widehat e\) factors through the underlying monomorphism
    \[
        m:I\hookrightarrow Q.
    \]

    At this point the transition map on \(I\) has not yet been constructed,
    so \(D_I\) is not yet being regarded over \(Z\) by a derivative-state
    coordinate. We therefore apply
    Lemma~\ref{lem:regular-epi-descent-for-subobjects-membership}
    in the underlying category \(\mathcal E\).

    Since \(\widehat e\) is regular epimorphic in \(\mathcal E\), \(m\) is
    monic in \(\mathcal E\), and
    \[
        h_\delta\widehat e
        =
        m(e\delta_P),
    \]
    there is a unique map
    \[
        \delta_I:D_I\to I
    \]
    satisfying
    \[
        m\delta_I
        =
        \delta_Q\widehat m.
    \]

    Define
    \[
        \omega_I
        :=
        \omega_Q\widehat m:
        D_I\to B_1.
    \]

    The transition has the required input typing:
    \[
    \begin{aligned}
        p_I\delta_I
        &=
        p_Qm\delta_I
        \\
        &=
        p_Q\delta_Q\widehat m
        \\
        &=
        s_A\lambda_Q\widehat m
        \\
        &=
        s_A\lambda_I.
    \end{aligned}
    \]

    The output map has the required source typing:
    \[
    \begin{aligned}
        s_B\omega_I
        &=
        s_B\omega_Q\widehat m
        \\
        &=
        q_Q\delta_Q\widehat m
        \\
        &=
        q_Qm\delta_I
        \\
        &=
        q_I\delta_I.
    \end{aligned}
    \]

    It also has the required target typing:
    \[
    \begin{aligned}
        t_B\omega_I
        &=
        t_B\omega_Q\widehat m
        \\
        &=
        q_Q\rho_Q\widehat m
        \\
        &=
        q_Qm\rho_I
        \\
        &=
        q_I\rho_I.
    \end{aligned}
    \]

    These equations show that \(\delta_I\) and \(\omega_I\) are correctly
    typed. In particular, \(D_I\) may now be regarded over \(Z\) by
    \[
        \left\langle
            s_A\lambda_I,\,
            q_I\delta_I
        \right\rangle:
        D_I\to A_0\times B_0.
    \]

    The transition unit law holds after composing with \(m\):
    \[
    \begin{aligned}
        m\delta_I(1_{p_Iz},z)
        &=
        \delta_Q(1_{p_Qmz},mz)
        \\
        &=
        mz.
    \end{aligned}
    \]
    Since \(m\) is monic,
    \[
        \delta_I(1_{p_Iz},z)=z.
    \]

    For composable arrows
    \[
        a:u\to v,
        \qquad
        b:v\to w,
    \]
    and \(z\in I\) over \(w\),
    \[
    \begin{aligned}
        m\delta_I(b\circ a,z)
        &=
        \delta_Q(b\circ a,mz)
        \\
        &=
        \delta_Q
        \bigl(
            a,
            \delta_Q(b,mz)
        \bigr)
        \\
        &=
        \delta_Q
        \bigl(
            a,
            m\delta_I(b,z)
        \bigr)
        \\
        &=
        m\delta_I
        \bigl(
            a,
            \delta_I(b,z)
        \bigr).
    \end{aligned}
    \]
    Monicity of \(m\) gives
    \[
        \delta_I(b\circ a,z)
        =
        \delta_I
        \bigl(
            a,
            \delta_I(b,z)
        \bigr).
    \]

    The output unit law follows by restriction:
    \[
    \begin{aligned}
        \omega_I(1_{p_Iz},z)
        &=
        \omega_Q(1_{p_Qmz},mz)
        \\
        &=
        1_{q_Qmz}
        \\
        &=
        1_{q_Iz}.
    \end{aligned}
    \]

    Similarly,
    \[
    \begin{aligned}
        \omega_I(b\circ a,z)
        &=
        \omega_Q(b\circ a,mz)
        \\
        &=
        \omega_Q(b,mz)
        \circ
        \omega_Q
        \bigl(
            a,
            \delta_Q(b,mz)
        \bigr)
        \\
        &=
        \omega_I(b,z)
        \circ
        \omega_Q
        \bigl(
            a,
            m\delta_I(b,z)
        \bigr)
        \\
        &=
        \omega_I(b,z)
        \circ
        \omega_I
        \bigl(
            a,
            \delta_I(b,z)
        \bigr).
    \end{aligned}
    \]

    Hence
    \[
        I=(I,p_I,q_I,\delta_I,\omega_I)
    \]
    is a Mealy morphism.

    The equations
    \[
        m\delta_I
        =
        \delta_Q\widehat m
    \]
    and
    \[
        \omega_I
        =
        \omega_Q\widehat m
    \]
    show that \(m\) is strict.

    To show that \(e\) is strict, compute
    \[
    \begin{aligned}
        me\delta_P
        &=
        f\delta_P
        \\
        &=
        \delta_Q\widehat f
        \\
        &=
        \delta_Q\widehat m\,\widehat e
        \\
        &=
        m\delta_I\widehat e.
    \end{aligned}
    \]
    Since \(m\) is monic,
    \[
        e\delta_P
        =
        \delta_I\widehat e.
    \]

    For the output,
    \[
    \begin{aligned}
        \omega_I\widehat e
        &=
        \omega_Q\widehat m\,\widehat e
        \\
        &=
        \omega_Q\widehat f
        \\
        &=
        \omega_P.
    \end{aligned}
    \]
    Hence \(e\) is strict.

    Since the carrier of \(m\) is monic,
    Lemma~\ref{lem:coeq-carrier-created-pullbacks-monomorphisms}
    implies that \(m\) is monic in the strict Mealy category.

    We next prove that \(e\) is a regular epimorphism in the strict Mealy
    category. Let
    \[
        R_e\rightrightarrows P
    \]
    be its kernel pair there. By
    Lemma~\ref{lem:coeq-carrier-created-pullbacks-monomorphisms}, this kernel
    pair is created by the carrier functor.

    Let
    \[
        h:P\to X
    \]
    be a strict semantic homomorphism coequalizing the kernel pair of \(e\).
    Since the carrier map \(e\) is the coequalizer of its carrier kernel pair,
    the underlying carrier map of \(h\) descends uniquely to a map
    \[
        \bar h:I\to X
    \]
    in \(\mathcal E/Z\) satisfying
    \[
        \bar h e=h.
    \]

    To prove transition preservation, precompose with the regular epimorphism
    \(\widehat e\):
    \[
    \begin{aligned}
        \bar h\delta_I\widehat e
        &=
        \bar h e\delta_P
        \\
        &=
        h\delta_P
        \\
        &=
        \delta_X\widehat h
        \\
        &=
        \delta_X\widehat{\bar h}\,\widehat e.
    \end{aligned}
    \]
    Since \(\widehat e\) is epic,
    \[
        \bar h\delta_I
        =
        \delta_X\widehat{\bar h}.
    \]

    Likewise,
    \[
    \begin{aligned}
        \omega_X\widehat{\bar h}\,\widehat e
        &=
        \omega_X\widehat h
        \\
        &=
        \omega_P
        \\
        &=
        \omega_I\widehat e.
    \end{aligned}
    \]
    Since \(\widehat e\) is epic,
    \[
        \omega_X\widehat{\bar h}
        =
        \omega_I.
    \]
    Thus \(\bar h\) is strict.

    Therefore \(e\) is the coequalizer of its kernel pair in
    \[
        \mathsf{Mly}^{\mathrm{str}}_{\mathbb A,\mathbb B},
    \]
    and hence is a regular epimorphism there.

    Finally, suppose that \(f\) factors through a strict monomorphism
    \[
        n:J\hookrightarrow Q:
    \]
    \[
        f=nh
    \]
    for some strict semantic homomorphism
    \[
        h:P\to J.
    \]
    Since
    \[
        P\xrightarrow{e}I\xrightarrow{m}Q
    \]
    is the carrier image factorization of \(f\), there is a unique carrier map
    \[
        u:I\to J
    \]
    in \(\mathcal E/Z\) satisfying
    \[
        ue=h,
        \qquad
        nu=m.
    \]

    The map \(u\) preserves transition:
    \[
    \begin{aligned}
        nu\delta_I
        &=
        m\delta_I
        \\
        &=
        \delta_Q\widehat m
        \\
        &=
        \delta_Q\widehat n\,\widehat u
        \\
        &=
        n\delta_J\widehat u.
    \end{aligned}
    \]
    Since the carrier of \(n\) is monic,
    \[
        u\delta_I
        =
        \delta_J\widehat u.
    \]

    It also preserves output:
    \[
    \begin{aligned}
        \omega_J\widehat u
        &=
        \omega_Q\widehat n\,\widehat u
        \\
        &=
        \omega_Q\widehat m
        \\
        &=
        \omega_I.
    \end{aligned}
    \]
    Hence \(u\) is strict.

    Its uniqueness as a strict map follows from its uniqueness as a carrier
    map. Therefore \(m:I\hookrightarrow Q\) has the universal property of the
    image of \(f\), and
    \[
        P\xrightarrow{e}I\xrightarrow{m}Q
    \]
    is the regular image factorization in the strict Mealy category.

    Finally, the Mealy structure on \(I\) is unique. Indeed, strictness of
    \(m\) forces
    \[
        m\delta_I=\delta_Q\widehat m
    \]
    and
    \[
        \omega_I=\omega_Q\widehat m,
    \]
    and monicity of \(m\) determines \(\delta_I\) uniquely.
\end{proof}

\begin{corollary}[Regular epimorphisms are detected by carriers]
\label{cor:carrier-regular-quotients-categorical}
    Let
    \[
        e:P\to Q
    \]
    be a strict semantic homomorphism. Then the following are equivalent:
    \begin{enumerate}
        \item the underlying carrier map
        \[
            Ue:UP\to UQ
        \]
        is a regular epimorphism in
        \[
            \mathcal E/Z;
        \]

        \item the strict semantic homomorphism
        \[
            e:P\to Q
        \]
        is a regular epimorphism in
        \[
            \mathsf{Mly}^{\mathrm{str}}_{\mathbb A,\mathbb B}.
        \]
    \end{enumerate}

    Consequently, the regular quotients of
    Definition~\ref{def:regular-quotient-strict-mealy} are exactly the
    categorical regular epimorphisms in the strict Mealy category.
\end{corollary}

\begin{proof}
    \leavevmode
    \begin{enumerate}
        \item \textbf{Carrier regular epimorphisms are categorical regular
        epimorphisms.}

        Suppose that
        \[
            Ue:UP\to UQ
        \]
        is a regular epimorphism in \(\mathcal E/Z\).

        Apply
        Proposition~\ref{prop:regular-image-factorization-mealy} to \(e\).
        Since the carrier image of a regular epimorphism is its entire
        codomain, its carrier regular-image factorization is, up to the unique
        canonical isomorphism of image representatives,
        \[
            UP\xrightarrow{Ue}UQ\xrightarrow{1_{UQ}}UQ.
        \]
        The proposition creates this factorization in the strict Mealy
        category and identifies its first map as a categorical regular
        epimorphism. Hence \(e\) is a regular epimorphism in
        \[
            \mathsf{Mly}^{\mathrm{str}}_{\mathbb A,\mathbb B}.
        \]

        \item \textbf{Categorical regular epimorphisms have regular-epimorphic
        carriers.}

        Conversely, suppose
        \[
            e:P\to Q
        \]
        is a regular epimorphism in the strict Mealy category.

        Factor its carrier in \(\mathcal E/Z\) as
        \[
            UP
            \xrightarrow{Ue'}
            UI
            \xrightarrow{Um}
            UQ,
        \]
        where \(Ue'\) is regular epimorphic and \(Um\) is monic.

        By
        Proposition~\ref{prop:regular-image-factorization-mealy}, the object
        \(I\) carries a unique Mealy structure for which
        \[
            P\xrightarrow{e'}I\xrightarrow{m}Q
        \]
        is the regular image factorization of \(e\) in the strict Mealy
        category. Thus
        \[
            e=me',
        \]
        with \(e'\) a regular epimorphism and \(m\) a monomorphism.

        Regular epimorphisms are extremal. Since the regular epimorphism \(e\)
        factors through the monomorphism \(m\), the map \(m\) is an
        isomorphism. Therefore its carrier
        \[
            Um:UI\to UQ
        \]
        is an isomorphism in \(\mathcal E/Z\).

        Since
        \[
            Ue=(Um)(Ue')
        \]
        and \(Ue'\) is regular epimorphic, it follows that \(Ue\) is a regular
        epimorphism in \(\mathcal E/Z\).
    \end{enumerate}
\end{proof}

\begin{remark}
\label{rem:regular-quotient-terminology}
    By the preceding corollary, there is no distinction in this chapter
    between:

    \begin{enumerate}
        \item a strict semantic homomorphism whose carrier is a regular
        epimorphism in \(\mathcal E/Z\);

        \item a regular epimorphism in
        \[
            \mathsf{Mly}^{\mathrm{str}}_{\mathbb A,\mathbb B}.
        \]
    \end{enumerate}

    Accordingly, the terms
    \[
        \text{regular quotient}
        \qquad\text{and}\qquad
        \text{regular-epimorphic quotient}
    \]
    may be used interchangeably for strict semantic homomorphisms.
\end{remark}

\subsection{Semantic images}
\label{subsec:semantic-images}

\begin{definition}[Semantic image]
\label{def:semantic-image}
    For a Mealy morphism \(P\), its \textbf{semantic image} is the regular
    image factorization of the unique strict map to the terminal object:
    \[
        P\xrightarrow{e_P}\operatorname{Beh}(P)
        \xrightarrow{m_P}
        \mathbf{Beh}^{\mathrm{can}}_{\mathbb A,\mathbb B}.
    \]
    On carriers, this is the regular image factorization of
    \[
        \beta_P:P\to\mathbf{Beh}_0
    \]
    in \(\mathcal E/Z\).
\end{definition}

\begin{remark}[Image subobjects and chosen representatives]
\label{rem:semantic-image-subobject-convention}
    The notation
    \[
        \operatorname{Beh}(P)
    \]
    is used primarily for the regular-image subobject
    \[
        \operatorname{Im}(\beta_P)
        \hookrightarrow
        \mathbf{Beh}^{\mathrm{can}}_{\mathbb A,\mathbb B},
    \]
    regarded as an element of
    \[
        \operatorname{Sub}_{\mathcal E/Z}(\mathbf{Beh}_0).
    \]

    Accordingly, an equality such as
    \[
        \operatorname{Beh}(P)=V
    \]
    means equality of subobject classes in this subobject poset. If particular
    regular epi--mono factorizations have been chosen, the corresponding
    carrier objects need not be literally equal; instead, there is a unique
    canonical isomorphism between the chosen image representatives commuting
    with their monomorphisms into \(\mathbf{Beh}_0\).

    The same convention applies to chosen representatives of derivative
    images, joins, meets, and other regular-image subobjects appearing below.
\end{remark}

The semantic image is displayed by
\[
\begin{tikzcd}[column sep=large, row sep=large]
    P
        \arrow[rr, "\beta_P"]
        \arrow[dr, two heads, "e_P"']
    &&
    \mathbf{Beh}_0
    \\
    &
    \operatorname{Beh}(P)
        \arrow[ur, hook, "m_P"']
    &
\end{tikzcd}
\qquad
\text{in }\mathcal E/Z.
\]
The same factorization, with the Mealy structure created by
Proposition~\ref{prop:regular-image-factorization-mealy}, is
\[
\begin{tikzcd}[column sep=large, row sep=large]
    P
        \arrow[rr, "\beta_P"]
        \arrow[dr, two heads, "e_P"']
    &&
    \mathbf{Beh}^{\mathrm{can}}
    \\
    &
    \operatorname{Beh}(P)
        \arrow[ur, hook, "m_P"']
    &
\end{tikzcd}
\qquad
\text{in }\mathsf{Mly}^{\mathrm{str}}_{\mathbb A,\mathbb B}.
\]

\begin{corollary}[Semantic images are coequational theories]
\label{cor:semantic-images-are-coeq-theories}
    For every Mealy morphism \(P\), the semantic image
    \[
        \operatorname{Beh}(P)\hookrightarrow\mathbf{Beh}^{\mathrm{can}}
    \]
    is a local coequational theory. Moreover,
    \[
        P\twoheadrightarrow\operatorname{Beh}(P)
    \]
    is a regular quotient in
    \[
        \mathsf{Mly}^{\mathrm{str}}_{\mathbb A,\mathbb B}.
    \]
\end{corollary}

\begin{proof}
    \leavevmode
    \begin{enumerate}
        \item \textbf{Apply image creation.}
        The semantic image is the regular image of the strict semantic
        homomorphism
        \[
            \beta_P:P\to\mathbf{Beh}^{\mathrm{can}}.
        \]
        Proposition~\ref{prop:regular-image-factorization-mealy} creates this
        image in the strict Mealy category.

        \item \textbf{Identify the image as a subobject of the terminal object.}
        The monic part
        \[
            \operatorname{Beh}(P)\hookrightarrow\mathbf{Beh}^{\mathrm{can}}
        \]
        is a subobject, hence a local coequational theory.

        \item \textbf{Identify the quotient map.}
        The regular-epimorphic part
        \[
            P\twoheadrightarrow\operatorname{Beh}(P)
        \]
        is strict by the same proposition.
    \end{enumerate}
\end{proof}

\subsection{Meets and coproducts}
\label{subsec:meets-and-coproducts-local-theories}

\begin{lemma}[Meets of local theories]
\label{lem:meets-local-theories}
    If
    \[
        S,T\hookrightarrow
        \mathbf{Beh}^{\mathrm{can}}_{\mathbb A,\mathbb B}
    \]
    are local coequational theories, then their meet
    \[
        S\wedge T
        \hookrightarrow
        \mathbf{Beh}^{\mathrm{can}}_{\mathbb A,\mathbb B}
    \]
    is a local coequational theory.

    Its carrier is the pullback subobject
    \[
        S\wedge T
        =
        S\times_{\mathbf{Beh}_0}T
    \]
    in \(\mathcal E/Z\).
\end{lemma}

\begin{proof}
    Regard \(S\) and \(T\) as strict subobjects of the canonical behaviour
    machine by
    Proposition~\ref{prop:subobjects-canonical-behaviour-machine}. Form the
    pullback
    \[
    \begin{tikzcd}[column sep=large, row sep=large]
        S\wedge T
            \arrow[r]
            \arrow[d]
            \arrow[dr, phantom, "\lrcorner", very near start]
        &
        T
            \arrow[d, hook]
        \\
        S
            \arrow[r, hook]
        &
        \mathbf{Beh}^{\mathrm{can}}_{\mathbb A,\mathbb B}
    \end{tikzcd}
    \]
    in
    \[
        \mathsf{Mly}^{\mathrm{str}}_{\mathbb A,\mathbb B}.
    \]

    By
    Lemma~\ref{lem:coeq-carrier-created-pullbacks-monomorphisms}, this
    pullback is created by the carrier functor. Its carrier is therefore
    \[
        S\times_{\mathbf{Beh}_0}T
    \]
    in \(\mathcal E/Z\), which is the meet of the two carrier subobjects.

    The induced map
    \[
        S\wedge T
        \longrightarrow
        \mathbf{Beh}^{\mathrm{can}}_{\mathbb A,\mathbb B}
    \]
    is strict and has monic carrier. Hence it is a subobject in the strict
    Mealy category. By
    Proposition~\ref{prop:subobjects-canonical-behaviour-machine}, its carrier
    is a derivative-closed behavioural predicate. Therefore it is a local
    coequational theory.

    Explicitly, its transition and output are the restrictions
    \[
        \delta_{S\wedge T}(a,H)=\partial_aH
    \]
    and
    \[
        \omega_{S\wedge T}(a,H)
        =
        \omega_{\mathbf{Beh}}(a,H).
    \]
\end{proof}

\begin{proposition}[Small coproducts are created by the carrier slice]
\label{prop:small-coproducts-mealy}
    A small family
    \[
        (P_i)_{i\in I}
    \]
    of Mealy morphisms over \((\mathbb A,\mathbb B)\) has a coproduct in
    \[
        \mathsf{Mly}^{\mathrm{str}}_{\mathbb A,\mathbb B}.
    \]
    Its carrier is the coproduct
    \[
        \coprod_{i\in I}P_i
    \]
    in the slice
    \[
        \mathcal E/Z,
    \]
    and its transition and output maps are induced componentwise.
\end{proposition}

\begin{proof}
    \leavevmode
    \begin{enumerate}
        \item \textbf{Use preservation of coproducts by pullback.}
        In a Grothendieck topos, pullback preserves small coproducts in every
        slice. Hence
        \[
            A_1\times_{t_A,A_0,p_{\coprod_iP_i}}
            \coprod_iP_i
            \cong
            \coprod_i
            \left(
                A_1\times_{t_A,A_0,p_{P_i}}P_i
            \right).
        \]
        The coproduct of transition domains is displayed by
        \[
        \begin{tikzcd}[column sep=large, row sep=large]
            A_1\times_{A_0}P_i
                \arrow[r]
                \arrow[d]
            &
            A_1\times_{A_0}\coprod_iP_i
                \arrow[d]
            \\
            P_i
                \arrow[r]
            &
            \coprod_iP_i .
        \end{tikzcd}
        \]

        \item \textbf{Define transition and output componentwise.}
        The transition and output maps are induced by the component maps
        \[
            \delta_{P_i},
            \qquad
            \omega_{P_i}.
        \]
        Equivalently, for each \(i\) the following diagrams commute:
        \[
        \begin{tikzcd}[column sep=huge, row sep=large]
            A_1\times_{A_0}P_i
                \arrow[r]
                \arrow[d, "\delta_{P_i}"']
            &
            A_1\times_{A_0}\coprod_iP_i
                \arrow[d, "\delta_{\coprod_iP_i}"]
            \\
            P_i
                \arrow[r]
            &
            \coprod_iP_i
        \end{tikzcd}
        \qquad
        \begin{tikzcd}[column sep=huge, row sep=large]
            A_1\times_{A_0}P_i
                \arrow[r]
                \arrow[dr, "\omega_{P_i}"']
            &
            A_1\times_{A_0}\coprod_iP_i
                \arrow[d, "\omega_{\coprod_iP_i}"]
            \\
            &
            B_1 .
        \end{tikzcd}
        \]

        \item \textbf{Verify the laws.}
        The Mealy laws hold after restriction to each coproduct summand. Since
        the summand inclusions are jointly epic, the laws hold on the
        coproduct.

        \item \textbf{Check the universal property.}
        A strict semantic homomorphism
        \[
            \coprod_iP_i\to Q
        \]
        is equivalently a compatible family of strict semantic homomorphisms
        \[
            P_i\to Q.
        \]
        This is the universal property of the coproduct in \(\mathcal E/Z\)
        together with the componentwise Mealy structure:
        \[
        \begin{tikzcd}[column sep=large, row sep=large]
            P_i
                \arrow[r]
                \arrow[dr, "f_i"']
            &
            \coprod_iP_i
                \arrow[d, dashed, "f"]
            \\
            &
            Q .
        \end{tikzcd}
        \]
    \end{enumerate}
\end{proof}

\section{Universal realization}
\label{sec:universal-realization}

\subsection{The behaviour--realization correspondence}
\label{subsec:universal-realization-correspondence}

The semantic-image construction is the local form of the classical
behaviour--realization adjunction. In the terminology of Goguen and Naudé,
behaviour and realization form an adjoint situation: a system has a universal
behaviour quotient, and a behaviour object has a canonical realization
\cite{Goguen1973,Naude1977,Naude1979}. In the present setting, behaviour
objects are local coequational subobjects of the terminal residual semantic
Mealy morphism.

Let
\[
    \mathsf{Coeq}_{\mathbb A,\mathbb B}
\]
denote the poset category of local coequational theories
\[
    V\hookrightarrow
    \mathbf{Beh}^{\mathrm{can}}_{\mathbb A,\mathbb B},
\]
ordered by inclusion as subobjects in
\[
    \mathcal E/Z,
    \qquad
    Z=A_0\times B_0.
\]
Thus a morphism
\[
    V\to W
\]
exists precisely when
\[
    V\leq W.
\]

If
\[
    j_V:V\hookrightarrow
    \mathbf{Beh}^{\mathrm{can}}_{\mathbb A,\mathbb B}
\]
is a local coequational theory, then by
Proposition~\ref{prop:subobjects-canonical-behaviour-machine} the carrier
\(V\) already carries the restricted canonical Mealy structure. Explicitly,
its carrier span is
\[
    A_0
        \xleftarrow{p_V}
    V
        \xrightarrow{q_V}
    B_0,
\]
its transition is
\[
    \delta_V(a,H)=\partial_aH,
\]
and its output is
\[
    \omega_V(a,H)=\omega_{\mathbf{Beh}}(a,H).
\]
Equivalently, the structure is obtained by restriction along
\[
    j_V:V\hookrightarrow\mathbf{Beh}_0:
\]
\[
\begin{tikzcd}[column sep=large, row sep=large]
    A_1\times_{t_A,A_0,p_V}V
        \arrow[r, "{1\times j_V}"]
        \arrow[d, "\delta_V"']
        \arrow[dr, phantom, "\circlearrowleft", very near start]
    &
    A_1\times_{t_A,A_0,p_{\mathbf{Beh}}}\mathbf{Beh}_0
        \arrow[d, "\delta_{\mathbf{Beh}}"]
        \arrow[r, "\omega_{\mathbf{Beh}}"]
    &
    B_1
    \\
    V
        \arrow[r, hook, "j_V"']
    &
    \mathbf{Beh}_0 .
\end{tikzcd}
\]
The output map is
\[
    \omega_V
    =
    \omega_{\mathbf{Beh}}(1\times j_V).
\]

If
\[
    V\leq W
\]
are local coequational theories, then the inclusion
\[
    V\hookrightarrow W
\]
is a strict semantic homomorphism, since both transition and output are
inherited from the canonical behaviour machine. Thus local coequational
theories embed into
\[
    \mathsf{Mly}^{\mathrm{str}}_{\mathbb A,\mathbb B}
\]
by sending each subobject to its restricted canonical Mealy morphism.

On the other side, the assignment
\[
    P\longmapsto \operatorname{Beh}(P)
\]
is functorial as a map
\[
    \mathsf{Mly}^{\mathrm{str}}_{\mathbb A,\mathbb B}
    \longrightarrow
    \mathsf{Coeq}_{\mathbb A,\mathbb B}.
\]
Indeed, if
\[
    f:P\to Q
\]
is a strict semantic homomorphism, then by
Lemma~\ref{lem:coeq-strict-maps-preserve-beta},
\[
    \beta_Qf=\beta_P.
\]
Taking regular images in the carrier slice gives
\[
    \operatorname{Beh}(P)\leq\operatorname{Beh}(Q),
\]
because the regular image of a composite is contained in the regular image of
its second factor. The image-containment diagram is
\[
\begin{tikzcd}[column sep=large, row sep=large]
    P
        \arrow[r, "f"]
        \arrow[d, two heads, "e_P"']
        \arrow[rr, bend left=18, "\beta_P"]
    &
    Q
        \arrow[r, "\beta_Q"]
        \arrow[d, two heads, "e_Q"']
    &
    \mathbf{Beh}^{\mathrm{can}}
    \\
    \operatorname{Beh}(P)
        \arrow[r, dashed]
    &
    \operatorname{Beh}(Q)
        \arrow[ur, hook, "m_Q"']
    &
\end{tikzcd}
\]
where the dashed arrow exists because
\[
    m_Qe_Qf=\beta_Qf=\beta_P.
\]
Since
\[
    \mathsf{Coeq}_{\mathbb A,\mathbb B}
\]
is posetal, these inclusions are automatically functorial.

\subsection{The universal property}
\label{subsec:universal-realization-universal-property}

\begin{theorem}[Universal realization adjunction]
\label{thm:universal-realization-adjunction}
    The semantic-image assignment
    \[
        \operatorname{Beh}:
        \mathsf{Mly}^{\mathrm{str}}_{\mathbb A,\mathbb B}
        \longrightarrow
        \mathsf{Coeq}_{\mathbb A,\mathbb B}
    \]
    is left adjoint to the inclusion of local coequational theories as
    restricted canonical Mealy morphisms
    \[
        \iota:
        \mathsf{Coeq}_{\mathbb A,\mathbb B}
        \longrightarrow
        \mathsf{Mly}^{\mathrm{str}}_{\mathbb A,\mathbb B}.
    \]
    Thus there is an adjunction
    \[
        \operatorname{Beh}
        \dashv
        \iota .
    \]
    Explicitly, for every Mealy morphism
    \[
        P:\mathbb A\rightsquigarrow\mathbb B
    \]
    and every local coequational theory
    \[
        V\hookrightarrow
        \mathbf{Beh}^{\mathrm{can}}_{\mathbb A,\mathbb B},
    \]
    there is a natural bijection
    \[
        \mathsf{Mly}^{\mathrm{str}}_{\mathbb A,\mathbb B}
        (P,\iota V)
        \cong
        \mathsf{Coeq}_{\mathbb A,\mathbb B}
        \bigl(\operatorname{Beh}(P),V\bigr).
    \]
    Equivalently,
    \[
        P\to V
        \text{ strictly}
        \qquad\Longleftrightarrow\qquad
        \operatorname{Beh}(P)\leq V.
    \]
\end{theorem}

\begin{proof}
    Let
    \[
        j_V:V\hookrightarrow
        \mathbf{Beh}^{\mathrm{can}}
    \]
    be a local coequational theory.

    \begin{enumerate}
        \item \textbf{From a strict map into \(V\) to semantic containment.}
        Suppose first that there is a strict semantic homomorphism
        \[
            f:P\to V.
        \]
        Then
        \[
            j_Vf:P\to\mathbf{Beh}^{\mathrm{can}}
        \]
        is a strict semantic homomorphism. By terminality of
        \[
            \mathbf{Beh}^{\mathrm{can}}
        \]
        in
        \[
            \mathsf{Mly}^{\mathrm{str}}_{\mathbb A,\mathbb B},
        \]
        this composite is the terminal semantic map:
        \[
            j_Vf=\beta_P.
        \]
        Hence \(\beta_P\) factors through \(V\). Since
        \[
            \operatorname{Beh}(P)
        \]
        is the regular image of \(\beta_P\), it follows that
        \[
            \operatorname{Beh}(P)\leq V.
        \]
        This implication is displayed by
        \[
        \begin{tikzcd}[column sep=large, row sep=large]
            P
                \arrow[r, "f"]
                \arrow[rr, bend left=18, "\beta_P"]
                \arrow[d, two heads, "e_P"']
            &
            V
                \arrow[r, hook, "j_V"]
            &
            \mathbf{Beh}^{\mathrm{can}}
            \\
            \operatorname{Beh}(P)
                \arrow[ur, dashed]
                \arrow[urr, hook, "m_P"', bend right=18]
            &
            &
        \end{tikzcd}
        \]

        \item \textbf{From semantic containment to a strict map into \(V\).}
        Conversely, suppose
        \[
            \operatorname{Beh}(P)\leq V.
        \]
        The semantic image factorization of \(P\) is
        \[
            P
                \xrightarrow{e_P}
            \operatorname{Beh}(P)
                \xrightarrow{m_P}
            \mathbf{Beh}^{\mathrm{can}}.
        \]
        Since
        \[
            \operatorname{Beh}(P)\leq V,
        \]
        the monomorphism \(m_P\) factors through \(j_V\):
        \[
        \begin{tikzcd}[column sep=large, row sep=large]
            \operatorname{Beh}(P)
                \arrow[rr, hook, "m_P"]
                \arrow[dr, dashed, "\bar m_P"']
            &&
            \mathbf{Beh}^{\mathrm{can}}
            \\
            &
            V
                \arrow[ur, hook, "j_V"']
            &
        \end{tikzcd}
        \]
        Therefore \(\beta_P\) factors through \(V\):
        \[
        \begin{tikzcd}[column sep=large, row sep=large]
            P
                \arrow[r, two heads, "e_P"]
                \arrow[rrr, bend left=18, "\beta_P"]
            &
            \operatorname{Beh}(P)
                \arrow[r, hook, "\bar m_P"]
            &
            V
                \arrow[r, hook, "j_V"]
            &
            \mathbf{Beh}^{\mathrm{can}} .
        \end{tikzcd}
        \]
        By Lemma~\ref{lem:model-factorizations-strict}, the induced map
        \[
            P\to V
        \]
        is a strict semantic homomorphism.

        \item \textbf{Prove uniqueness.}
        Let
        \[
            f,g:P\rightrightarrows V
        \]
        be strict semantic homomorphisms. Then
        \[
            j_Vf
            \qquad\text{and}\qquad
            j_Vg
        \]
        are strict semantic homomorphisms from \(P\) to
        \[
            \mathbf{Beh}^{\mathrm{can}}.
        \]
        Hence, by terminality,
        \[
            j_Vf=\beta_P=j_Vg.
        \]
        Since \(j_V\) is monic,
        \[
            f=g.
        \]

        \item \textbf{Conclude the bijection.}
        The left-hand hom-set
        \[
            \mathsf{Mly}^{\mathrm{str}}_{\mathbb A,\mathbb B}(P,V)
        \]
        is empty or singleton according as
        \[
            \operatorname{Beh}(P)\leq V
        \]
        fails or holds. This is exactly the right-hand hom-set in the poset
        \[
            \mathsf{Coeq}_{\mathbb A,\mathbb B}.
        \]
        Therefore the required natural bijection holds.
    \end{enumerate}
\end{proof}

\subsection{Unit, counit, and immediate consequences}
\label{subsec:universal-realization-consequences}

\begin{proposition}[Unit and counit]
\label{prop:universal-realization-unit-counit}
    Under the adjunction of
    Theorem~\ref{thm:universal-realization-adjunction}, the unit at a Mealy
    morphism \(P\) is the semantic quotient
    \[
        \eta_P:
        P\twoheadrightarrow
        \operatorname{Beh}(P).
    \]
    The counit at a local coequational theory \(V\) is the canonical
    isomorphism
    \[
        \operatorname{Beh}(V)
        \xrightarrow{\cong}
        V.
    \]
\end{proposition}

\begin{proof}
    \leavevmode
    \begin{enumerate}
        \item \textbf{Identify the unit.}
        The unit at \(P\) is the unique strict semantic homomorphism
        \[
            P\to\operatorname{Beh}(P)
        \]
        corresponding to the identity inclusion
        \[
            \operatorname{Beh}(P)\leq\operatorname{Beh}(P).
        \]
        By construction, this is the regular-epimorphic part of the semantic
        image factorization:
        \[
        \begin{tikzcd}[column sep=large, row sep=large]
            P
                \arrow[r, two heads, "\eta_P"]
                \arrow[rr, bend left=18, "\beta_P"]
            &
            \operatorname{Beh}(P)
                \arrow[r, hook, "m_P"]
            &
            \mathbf{Beh}^{\mathrm{can}} .
        \end{tikzcd}
        \]

        \item \textbf{Identify the counit.}
        Let
        \[
            j_V:V\hookrightarrow\mathbf{Beh}^{\mathrm{can}}
        \]
        be a local coequational theory. By
        Lemma~\ref{lem:semantic-map-of-a-restricted-canonical-subobject}, the
        semantic map of the restricted canonical object \(V\) is its inclusion:
        \[
            \beta_V=j_V.
        \]
        Hence the semantic image of \(V\) is exactly \(V\):
        \[
            \operatorname{Beh}(V)\cong V.
        \]
        Diagrammatically:
        \[
        \begin{tikzcd}[column sep=large, row sep=large]
            V
                \arrow[r, equal]
                \arrow[rr, bend left=18, "\beta_V=j_V"]
            &
            \operatorname{Beh}(V)
                \arrow[r, hook]
            &
            \mathbf{Beh}^{\mathrm{can}} .
        \end{tikzcd}
        \]
    \end{enumerate}
\end{proof}

\begin{corollary}[Full faithfulness of restricted canonical realization]
\label{cor:universal-realization-fully-faithful}
    The inclusion of
    \[
        \mathsf{Coeq}_{\mathbb A,\mathbb B}
    \]
    into
    \[
        \mathsf{Mly}^{\mathrm{str}}_{\mathbb A,\mathbb B}
    \]
    by restricted canonical Mealy morphisms is fully faithful.
\end{corollary}

\begin{proof}
    \leavevmode
    \begin{enumerate}
        \item \textbf{Use the standard adjunction criterion.}
        In an adjunction
        \[
            L\dashv R,
        \]
        the right adjoint \(R\) is fully faithful precisely when the counit is
        an isomorphism.

        \item \textbf{Apply the counit computation.}
        Here the counit is an isomorphism by
        Proposition~\ref{prop:universal-realization-unit-counit}. Therefore
        the inclusion of local coequational theories as restricted canonical
        Mealy morphisms is fully faithful.
    \end{enumerate}
\end{proof}

\begin{corollary}[Terminal realization of a coequational theory]
\label{cor:terminal-realization-coequational-theory}
    Let
    \[
        V\hookrightarrow\mathbf{Beh}^{\mathrm{can}}
    \]
    be a local coequational theory. Then \(V\), with its restricted canonical
    Mealy structure, is terminal in
    \[
        \mathsf{Mod}(V).
    \]
\end{corollary}

\begin{proof}
    \leavevmode
    \begin{enumerate}
        \item \textbf{Use the model condition.}
        If
        \[
            P\models V,
        \]
        then
        \[
            \operatorname{Beh}(P)\leq V.
        \]

        \item \textbf{Apply the universal property.}
        By Theorem~\ref{thm:universal-realization-adjunction}, the containment
        \[
            \operatorname{Beh}(P)\leq V
        \]
        is equivalent to the existence of a unique strict semantic homomorphism
        \[
            P\to V.
        \]

        \item \textbf{Identify the unique map.}
        Hence \(V\) is terminal in its model category. The unique map is the
        factorization of the terminal semantic map through \(V\):
        \[
        \begin{tikzcd}[column sep=large, row sep=large]
            P
                \arrow[rr, "\beta_P"]
                \arrow[dr, dashed, "\bar\beta_P"']
            &&
            \mathbf{Beh}^{\mathrm{can}}
            \\
            &
            V
                \arrow[ur, hook, "j_V"']
            &
        \end{tikzcd}
        \]
    \end{enumerate}
\end{proof}

\begin{corollary}[Behaviourally separated reflection]
\label{cor:behaviourally-separated-reflection}
    The semantic quotient
    \[
        \eta_P:
        P\twoheadrightarrow\operatorname{Beh}(P)
    \]
    is an isomorphism if and only if \(P\) is behaviourally separated.
\end{corollary}
\begin{proof}
    The semantic image factorization is
    \[
    \begin{tikzcd}[column sep=large, row sep=large]
        P
            \arrow[r, two heads, "\eta_P"]
            \arrow[rr, bend left=18, "\beta_P"]
        &
        \operatorname{Beh}(P)
            \arrow[r, hook, "m_P"]
        &
        \mathbf{Beh}^{\mathrm{can}} .
    \end{tikzcd}
    \]

    \begin{enumerate}
        \item \textbf{Behavioural separatedness implies reflection
        isomorphism.}

        Suppose \(P\) is behaviourally separated. Then
        \[
            \beta_P=m_P\eta_P
        \]
        is monic.

        We claim that \(\eta_P\) is monic. Let
        \[
            f,g:X\rightrightarrows P
        \]
        satisfy
        \[
            \eta_Pf=\eta_Pg.
        \]
        Then
        \[
        \begin{aligned}
            \beta_Pf
            &=
            m_P\eta_Pf
            \\
            &=
            m_P\eta_Pg
            \\
            &=
            \beta_Pg.
        \end{aligned}
        \]
        Since \(\beta_P\) is monic,
        \[
            f=g.
        \]
        Hence \(\eta_P\) is monic.

        By construction, \(\eta_P\) is also a regular epimorphism in the
        strict Mealy category. A morphism that is both a regular epimorphism
        and a monomorphism is an isomorphism. Therefore
        \[
            \eta_P:
            P\overset{\cong}{\longrightarrow}\operatorname{Beh}(P).
        \]

        \item \textbf{Reflection isomorphism implies behavioural
        separatedness.}

        Conversely, suppose
        \[
            \eta_P:
            P\overset{\cong}{\longrightarrow}\operatorname{Beh}(P)
        \]
        is an isomorphism. Since
        \[
            m_P:
            \operatorname{Beh}(P)\hookrightarrow
            \mathbf{Beh}^{\mathrm{can}}
        \]
        is monic, the composite
        \[
            \beta_P=m_P\eta_P
        \]
        is monic. Therefore \(P\) is behaviourally separated.
    \end{enumerate}
\end{proof} 

\subsection{Relation with minimal residual realization}
\label{subsec:universal-realization-minimal-residual-realization}

\begin{remark}[Relation with minimal residual realization]
\label{rem:universal-realization-minimal-residual-realization}
    Let
    \[
        k:K\to\mathbf{Beh}_0
    \]
    be a behaviour specification. By
    Theorem~\ref{thm:derivative-closure}, its derivative image
    \[
        \operatorname{Der}_0(K)\hookrightarrow\mathbf{Beh}_0
    \]
    is derivative-closed, hence is a local coequational theory. The minimal
    residual realization of
    Chapter~\ref{ch:behaviour-categories-myhill-nerode} is precisely the
    restricted canonical Mealy morphism on this derivative image:
    \[
        \mathsf{Min}(K)
        =
        \operatorname{Der}_0(K)
    \]
    as a restricted canonical Mealy morphism.

    If
    \[
        i:K\to P
    \]
    is a reachable realization of \(k\), then
    \[
        \operatorname{Beh}(P)
        =
        \operatorname{Der}_0(K).
    \]
    Under this identification, the unit at \(P\) is the regular semantic
    quotient
    \[
        P\twoheadrightarrow\mathsf{Min}(K)
    \]
    from the categorical Myhill--Nerode theorem. This is the diagram
    \[
    \begin{tikzcd}[column sep=large, row sep=large]
        P
            \arrow[r, two heads]
            \arrow[rr, bend left=18, "\beta_P"]
        &
        \operatorname{Beh}(P)
            \arrow[r, hook]
            \arrow[d, "\cong"']
        &
        \mathbf{Beh}^{\mathrm{can}}
        \\
        &
        \operatorname{Der}_0(K)
            \arrow[ur, hook, "j_K"']
        &
    \end{tikzcd}
    \]
    and the restricted canonical Mealy morphism on
    \[
        \operatorname{Der}_0(K)
    \]
    is
    \[
        \mathsf{Min}(K).
    \]
\end{remark}

\section{The local coequational Birkhoff theorem}
\label{sec:local-coequational-birkhoff}

\subsection{Closure properties of model classes}
\label{subsec:closure-properties-model-classes}

\begin{theorem}[Local coequational closure]
\label{thm:local-coequational-closure}
    Let
    \[
        V\hookrightarrow\mathbf{Beh}^{\mathrm{can}}
    \]
    be a local coequational theory. Then:
    \begin{enumerate}
        \item \(\mathsf{Mod}(V)\) is closed under strict homomorphic preimages.
        \item \(\mathsf{Mod}(V)\) is closed under regular-epimorphic quotients.
        \item \(\mathsf{Mod}(V)\) is closed under small coproducts.
        \item \(V\), regarded as the restricted canonical object, is itself a
        model of \(V\), and
        \[
            \operatorname{Beh}(V)=V.
        \]
    \end{enumerate}
    Consequently the assignment
    \[
        V\longmapsto\mathsf{Mod}(V)
    \]
    is monotone and order-reflecting.
\end{theorem}

\begin{proof}
    Let
    \[
        j_V:V\hookrightarrow\mathbf{Beh}^{\mathrm{can}}
    \]
    denote the inclusion.

    \begin{enumerate}
        \item \textbf{Closure under strict homomorphic preimages.}
        Let
        \[
            f:P\to Q
        \]
        be strict and suppose \(Q\models V\). Choose
        \[
            \bar\beta_Q:Q\to V
        \]
        with
        \[
            j_V\bar\beta_Q=\beta_Q.
        \]
        Then, by Lemma~\ref{lem:coeq-strict-maps-preserve-beta},
        \[
            \beta_P=\beta_Qf=j_V\bar\beta_Qf,
        \]
        so \(P\models V\). The factorization is
        \[
        \begin{tikzcd}[column sep=large, row sep=large]
            P
                \arrow[r, "f"]
                \arrow[rr, bend left=16, "\beta_P"]
                \arrow[dr, dashed, "\bar\beta_Qf"']
            &
            Q
                \arrow[r, "\beta_Q"]
                \arrow[d, "\bar\beta_Q"]
            &
            \mathbf{Beh}^{\mathrm{can}}
            \\
            &
            V
                \arrow[ur, hook, "j_V"']
            &
        \end{tikzcd}
        \]

        \item \textbf{Closure under regular-epimorphic quotients.}
        Let
        \[
            e:P\twoheadrightarrow Q
        \]
        be a regular-epimorphic strict semantic homomorphism and suppose
        \(P\models V\). Since
        \[
            \beta_P=\beta_Qe,
        \]
        and \(\beta_P\) factors through
        \[
            j_V:V\hookrightarrow\mathbf{Beh}^{\mathrm{can}},
        \]
        Lemma~\ref{lem:regular-epi-descent-for-subobjects-membership} implies that
        \(\beta_Q\) factors through \(j_V\). Hence \(Q\models V\).
        Diagrammatically,
        \[
        \begin{tikzcd}[column sep=large, row sep=large]
            P
                \arrow[r, two heads, "e"]
                \arrow[dr, "\bar\beta_P"']
                \arrow[ddr, "\beta_P"', bend right=18]
            &
            Q
                \arrow[d, dashed, "\bar\beta_Q"]
                \arrow[dd, "\beta_Q", bend left=18]
            \\
            &
            V
                \arrow[d, hook, "j_V"]
            \\
            &
            \mathbf{Beh}^{\mathrm{can}} .
        \end{tikzcd}
        \]

        \item \textbf{Closure under small coproducts.}
        Let
        \[
            (P_i)_{i\in I}
        \]
        be a small family of models. For each \(i\), the factorization
        \[
            \bar\beta_i:P_i\to V
        \]
        is strict by Lemma~\ref{lem:model-factorizations-strict}. By the
        coproduct universal property in
        \(\mathsf{Mly}^{\mathrm{str}}_{\mathbb A,\mathbb B}\), these strict
        maps induce a strict map
        \[
            \coprod_iP_i\to V.
        \]
        Therefore the coproduct models \(V\). The construction is
        \[
        \begin{tikzcd}[column sep=large, row sep=large]
            P_i
                \arrow[r]
                \arrow[dr, "\bar\beta_i"']
            &
            \coprod_iP_i
                \arrow[d, dashed, "\coprod_i\bar\beta_i"]
                \arrow[dd, bend left=18, "\beta_{\coprod_iP_i}"]
            \\
            &
            V
                \arrow[d, hook, "j_V"]
            \\
            &
            \mathbf{Beh}^{\mathrm{can}} .
        \end{tikzcd}
        \]

        \item \textbf{The theory models itself.}
        The identity map
        \[
            V\to V
        \]
        exhibits \(V\) as a model of itself. By
        Lemma~\ref{lem:semantic-map-of-a-restricted-canonical-subobject}, the
        semantic map of a restricted canonical subobject is its inclusion
        \[
            j_V:V\hookrightarrow\mathbf{Beh}^{\mathrm{can}}.
        \]
        Hence
        \[
            \operatorname{Beh}(V)=V.
        \]

        \item \textbf{Prove monotonicity and order-reflection.}
        If \(V\leq W\), then every factorization through \(V\) is a
        factorization through \(W\), so
        \[
            \mathsf{Mod}(V)\subseteq\mathsf{Mod}(W).
        \]
        Conversely, if
        \[
            \mathsf{Mod}(V)\subseteq\mathsf{Mod}(W),
        \]
        then \(V\in\mathsf{Mod}(V)\), hence \(V\in\mathsf{Mod}(W)\). Therefore
        the inclusion
        \[
            V\hookrightarrow\mathbf{Beh}^{\mathrm{can}}
        \]
        factors through
        \[
            W\hookrightarrow\mathbf{Beh}^{\mathrm{can}},
        \]
        so \(V\leq W\). This is the triangle
        \[
        \begin{tikzcd}[column sep=large, row sep=large]
            V
                \arrow[rr, hook, "j_V"]
                \arrow[dr, dashed]
            &&
            \mathbf{Beh}^{\mathrm{can}}
            \\
            &
            W
                \arrow[ur, hook, "j_W"']
            &
        \end{tikzcd}
        \]
        obtained from the fact that \(V\) is a model of \(W\).
    \end{enumerate}
\end{proof}

Closure under small coproducts includes the empty coproduct. Hence the empty
replete full subcategory is not coequationally definable.

\subsection{Joins of local theories}
\label{subsec:joins-local-theories}

Let
\[
    \mathcal C\subseteq\mathsf{Mly}^{\mathrm{str}}_{\mathbb A,\mathbb B}
\]
be a replete full subcategory. Define
\[
    \mathcal S_{\mathcal C}
    :=
    \{
        S\hookrightarrow\mathbf{Beh}^{\mathrm{can}}
        \mid
        S=\operatorname{Beh}(P)
        \text{ for some }P\in\mathcal C
    \}.
\]
Although \(\mathcal C\) may be large, the collection
\[
    \mathcal S_{\mathcal C}
\]
is a subclass of the set
\[
    \operatorname{Sub}_{\mathcal E/Z}(\mathbf{Beh}_0).
\]
We therefore choose a small set of representatives for the subobjects in
\(\mathcal S_{\mathcal C}\) and form their join. Define
\[
    V_{\mathcal C}
    :=
    \bigvee_{S\in\mathcal S_{\mathcal C}}S
    \hookrightarrow\mathbf{Beh}^{\mathrm{can}},
\]
where the displayed join is taken over the chosen small representative set.

\begin{lemma}[Joins of local coequational theories]
\label{lem:joins-derivative-closed}
    Let
    \[
        (S_i\hookrightarrow\mathbf{Beh}^{\mathrm{can}})_{i\in I}
    \]
    be a small family of local coequational theories. Then their join
    \[
        S:=\bigvee_{i\in I}S_i
    \]
    in
    \[
        \operatorname{Sub}_{\mathcal E/Z}(\mathbf{Beh}_0)
    \]
    is again a local coequational theory.
\end{lemma}

\begin{proof}
    \leavevmode
    \begin{enumerate}
        \item \textbf{Represent the join.}
        In the Grothendieck topos \(\mathcal E/Z\), the join is represented by
        the regular image of the coproduct map
        \[
            \coprod_iS_i\to\mathbf{Beh}_0.
        \]
        Write this image factorization as
        \[
        \begin{tikzcd}[column sep=large, row sep=large]
            \coprod_iS_i
                \arrow[rr]
                \arrow[dr, two heads]
            &&
            \mathbf{Beh}_0
            \\
            &
            S
                \arrow[ur, hook]
            &
        \end{tikzcd}
        \]
        in \(\mathcal E/Z\).

        \item \textbf{Regard the coproduct as a Mealy object.}
        By Proposition~\ref{prop:small-coproducts-mealy}, the coproduct
        \[
            \coprod_iS_i
        \]
        carries the componentwise Mealy structure. Since each
        \[
            S_i\hookrightarrow\mathbf{Beh}^{\mathrm{can}}
        \]
        is strict, the induced coproduct map to
        \(\mathbf{Beh}^{\mathrm{can}}\) is strict. Strictness is tested after
        precomposition with the jointly epic coproduct injections:
        \[
        \begin{tikzcd}[column sep=large, row sep=large]
            S_i
                \arrow[r]
                \arrow[dr, hook]
            &
            \coprod_iS_i
                \arrow[d, dashed]
            \\
            &
            \mathbf{Beh}^{\mathrm{can}} .
        \end{tikzcd}
        \]

        \item \textbf{Take the regular image in the strict Mealy category.}
        By Proposition~\ref{prop:regular-image-factorization-mealy}, the
        regular image of this strict map is created in \(\mathcal E/Z\) and
        carries a Mealy structure. Hence
        \[
            S\hookrightarrow\mathbf{Beh}^{\mathrm{can}}
        \]
        is a subobject in the strict Mealy category. The created image is
        \[
        \begin{tikzcd}[column sep=large, row sep=large]
            \coprod_iS_i
                \arrow[rr]
                \arrow[dr, two heads]
            &&
            \mathbf{Beh}^{\mathrm{can}}
            \\
            &
            S
                \arrow[ur, hook]
            &
        \end{tikzcd}
        \]
        in \(\mathsf{Mly}^{\mathrm{str}}_{\mathbb A,\mathbb B}\).

        \item \textbf{Conclude local coequationality.}
        By Proposition~\ref{prop:subobjects-canonical-behaviour-machine},
        being a subobject of the canonical behaviour object is precisely being
        a local coequational theory.
    \end{enumerate}
\end{proof}

\begin{lemma}[The representing theory belongs to the class]
\label{lem:representing-theory-belongs-to-class}
    Let
    \[
        \mathcal C
        \subseteq
        \mathsf{Mly}^{\mathrm{str}}_{\mathbb A,\mathbb B}
    \]
    be a replete full subcategory closed under regular-epimorphic quotients
    and small coproducts. Let
    \[
        \mathcal S_{\mathcal C}
        :=
        \{
            S\hookrightarrow\mathbf{Beh}^{\mathrm{can}}
            \mid
            S=\operatorname{Beh}(P)
            \text{ for some }P\in\mathcal C
        \},
    \]
    and let
    \[
        V_{\mathcal C}
        :=
        \bigvee_{S\in\mathcal S_{\mathcal C}}S
        \hookrightarrow
        \mathbf{Beh}^{\mathrm{can}}
    \]
    be formed using a chosen small representative set of the distinct
    semantic-image subobjects.

    Then the restricted canonical Mealy morphism
    \[
        V_{\mathcal C}
    \]
    belongs to \(\mathcal C\).
\end{lemma}

\begin{proof}
    For every chosen representative
    \[
        S\in\mathcal S_{\mathcal C},
    \]
    choose an object
    \[
        P_S\in\mathcal C
    \]
    together with an identification
    \[
        S=\operatorname{Beh}(P_S)
    \]
    as a subobject of
    \[
        \mathbf{Beh}^{\mathrm{can}}_{\mathbb A,\mathbb B}.
    \]

    The semantic-image factorization of \(P_S\) gives a regular quotient
    \[
        P_S\twoheadrightarrow\operatorname{Beh}(P_S).
    \]
    Since \(\mathcal C\) is closed under regular-epimorphic quotients,
    \[
        \operatorname{Beh}(P_S)\in\mathcal C.
    \]
    The chosen representative \(S\) is canonically isomorphic to this
    semantic-image object. Repleteness therefore gives
    \[
        S\in\mathcal C.
    \]

    Form the small coproduct
    \[
        C_{\mathcal C}
        :=
        \coprod_{S\in\mathcal S_{\mathcal C}}S.
    \]
    Since every summand belongs to \(\mathcal C\) and \(\mathcal C\) is
    closed under small coproducts,
    \[
        C_{\mathcal C}\in\mathcal C.
    \]

    The inclusions
    \[
        S\hookrightarrow
        \mathbf{Beh}^{\mathrm{can}}_{\mathbb A,\mathbb B}
    \]
    induce a strict semantic homomorphism
    \[
        c_{\mathcal C}:
        C_{\mathcal C}
        \longrightarrow
        \mathbf{Beh}^{\mathrm{can}}_{\mathbb A,\mathbb B}.
    \]
    Its regular image in the carrier slice is the join of the component
    subobjects:
    \[
        \operatorname{Im}(c_{\mathcal C})
        =
        \bigvee_{S\in\mathcal S_{\mathcal C}}S
        =
        V_{\mathcal C}.
    \]
    By
    Proposition~\ref{prop:regular-image-factorization-mealy}, this carrier
    image is created in the strict Mealy category. Thus there is a regular
    quotient
    \[
        C_{\mathcal C}
        \twoheadrightarrow
        V_{\mathcal C}
    \]
    fitting into
    \[
    \begin{tikzcd}[column sep=large, row sep=large]
        \displaystyle
        \coprod_{S\in\mathcal S_{\mathcal C}}S
            \arrow[rr, "c_{\mathcal C}"]
            \arrow[dr, two heads]
        &&
        \mathbf{Beh}^{\mathrm{can}}_{\mathbb A,\mathbb B}
        \\
        &
        V_{\mathcal C}
            \arrow[ur, hook, "j_{V_{\mathcal C}}"']
        &
    \end{tikzcd}
    \]
    in
    \[
        \mathsf{Mly}^{\mathrm{str}}_{\mathbb A,\mathbb B}.
    \]

    Since \(C_{\mathcal C}\in\mathcal C\) and \(\mathcal C\) is closed under
    regular-epimorphic quotients, the chosen image representative belongs to
    \(\mathcal C\) up to canonical isomorphism. Repleteness therefore gives
    \[
        V_{\mathcal C}\in\mathcal C.
    \]
\end{proof}

\subsection{The theorem}
\label{subsec:local-coequational-birkhoff-theorem}

\begin{theorem}[Local coequational Birkhoff theorem]
\label{thm:local-coequational-birkhoff}
    Let
    \[
        \mathcal C\subseteq\mathsf{Mly}^{\mathrm{str}}_{\mathbb A,\mathbb B}
    \]
    be a replete full subcategory. Then \(\mathcal C\) is of the form
    \[
        \mathcal C=\mathsf{Mod}(V)
    \]
    for a local coequational theory
    \[
        V\hookrightarrow\mathbf{Beh}^{\mathrm{can}}_{\mathbb A,\mathbb B}
    \]
    if and only if \(\mathcal C\) is closed under:
    \begin{enumerate}
        \item strict homomorphic preimages;
        \item regular-epimorphic quotients;
        \item small coproducts.
    \end{enumerate}
    When these conditions hold, the defining theory is
    \[
        V_{\mathcal C}
        =
        \bigvee_{P\in\mathcal C}\operatorname{Beh}(P)
        \hookrightarrow\mathbf{Beh}^{\mathrm{can}},
    \]
    where the join is taken over a chosen small set of representatives for the
    resulting semantic-image subobjects.
\end{theorem}

\begin{proof}
    \leavevmode
    \begin{enumerate}
        \item \textbf{Prove necessity.}

        Suppose
        \[
            \mathcal C=\mathsf{Mod}(V)
        \]
        for a local coequational theory \(V\). Closure under strict
        homomorphic preimages, regular-epimorphic quotients, and small
        coproducts is
        Theorem~\ref{thm:local-coequational-closure}.

        \item \textbf{Construct the representing theory.}

        Conversely, suppose that \(\mathcal C\) is closed under strict
        homomorphic preimages, regular-epimorphic quotients, and small
        coproducts.

        Define
        \[
            \mathcal S_{\mathcal C}
            :=
            \{
                S\hookrightarrow\mathbf{Beh}^{\mathrm{can}}
                \mid
                S=\operatorname{Beh}(P)
                \text{ for some }P\in\mathcal C
            \}.
        \]
        Choose a small representative set of the distinct subobjects in
        \(\mathcal S_{\mathcal C}\), and put
        \[
            V_{\mathcal C}
            :=
            \bigvee_{S\in\mathcal S_{\mathcal C}}S
            \hookrightarrow
            \mathbf{Beh}^{\mathrm{can}}_{\mathbb A,\mathbb B}.
        \]

        Every semantic image
        \[
            \operatorname{Beh}(P)
            \hookrightarrow
            \mathbf{Beh}^{\mathrm{can}}_{\mathbb A,\mathbb B}
        \]
        is a local coequational theory by
        Corollary~\ref{cor:semantic-images-are-coeq-theories}. Hence
        Lemma~\ref{lem:joins-derivative-closed} implies that
        \[
            V_{\mathcal C}
        \]
        is a local coequational theory.

        \item \textbf{Place the representing theory inside the class.}

        By
        Lemma~\ref{lem:representing-theory-belongs-to-class},
        \[
            V_{\mathcal C}\in\mathcal C.
        \]

        \item \textbf{Show that every object of \(\mathcal C\) models the
        representing theory.}

        Let
        \[
            P\in\mathcal C.
        \]
        By construction,
        \[
            \operatorname{Beh}(P)\leq V_{\mathcal C}.
        \]
        The semantic map factors through its regular image:
        \[
            P
            \twoheadrightarrow
            \operatorname{Beh}(P)
            \hookrightarrow
            \mathbf{Beh}^{\mathrm{can}}.
        \]
        Composing the semantic-image inclusion with
        \[
            \operatorname{Beh}(P)\hookrightarrow V_{\mathcal C}
        \]
        gives a factorization of \(\beta_P\) through \(V_{\mathcal C}\):
        \[
        \begin{tikzcd}[column sep=large, row sep=large]
            P
                \arrow[r, two heads]
                \arrow[rrr, bend left=18, "\beta_P"]
            &
            \operatorname{Beh}(P)
                \arrow[r, hook]
            &
            V_{\mathcal C}
                \arrow[r, hook, "j_{V_{\mathcal C}}"]
            &
            \mathbf{Beh}^{\mathrm{can}} .
        \end{tikzcd}
        \]
        Hence
        \[
            P\models V_{\mathcal C}.
        \]
        Therefore
        \[
            \mathcal C
            \subseteq
            \mathsf{Mod}(V_{\mathcal C}).
        \]

        \item \textbf{Use the representing theory to prove the converse
        inclusion.}

        Let
        \[
            X\in\mathsf{Mod}(V_{\mathcal C}).
        \]
        By definition, the semantic map of \(X\) admits a unique carrier
        factorization
        \[
            \bar\beta_X:
            X\to V_{\mathcal C}
        \]
        satisfying
        \[
            j_{V_{\mathcal C}}\bar\beta_X
            =
            \beta_X.
        \]
        This is the triangle
        \[
        \begin{tikzcd}[column sep=large, row sep=large]
            X
                \arrow[rr, "\beta_X"]
                \arrow[dr, "\bar\beta_X"']
            &&
            \mathbf{Beh}^{\mathrm{can}}_{\mathbb A,\mathbb B}
            \\
            &
            V_{\mathcal C}
                \arrow[ur, hook, "j_{V_{\mathcal C}}"']
            &
        \end{tikzcd}
        \]

        By Lemma~\ref{lem:model-factorizations-strict}, the factorization
        \[
            \bar\beta_X:X\to V_{\mathcal C}
        \]
        is a strict semantic homomorphism.

        Since
        \[
            V_{\mathcal C}\in\mathcal C
        \]
        and \(\mathcal C\) is closed under strict homomorphic preimages, it
        follows immediately that
        \[
            X\in\mathcal C.
        \]
        Therefore
        \[
            \mathsf{Mod}(V_{\mathcal C})
            \subseteq
            \mathcal C.
        \]

        \item \textbf{Conclude.}

        Combining the two inclusions gives
        \[
            \mathcal C
            =
            \mathsf{Mod}(V_{\mathcal C}).
        \]
    \end{enumerate}
\end{proof}

\begin{corollary}[Uniqueness of the defining local theory]
\label{cor:uniqueness-defining-local-coequational-theory}
    Let
    \[
        V\hookrightarrow
        \mathbf{Beh}^{\mathrm{can}}_{\mathbb A,\mathbb B}
    \]
    be a local coequational theory. Then
    \[
        V
        =
        \bigvee_{P\in\mathsf{Mod}(V)}
        \operatorname{Beh}(P)
    \]
    as a subobject of the canonical behaviour machine.

    Consequently, if a replete full subcategory
    \[
        \mathcal C
        \subseteq
        \mathsf{Mly}^{\mathrm{str}}_{\mathbb A,\mathbb B}
    \]
    is coequationally definable, then its defining local coequational theory
    is unique and is
    \[
        V_{\mathcal C}
        =
        \bigvee_{P\in\mathcal C}
        \operatorname{Beh}(P).
    \]
\end{corollary}

\begin{proof}
    If
    \[
        P\in\mathsf{Mod}(V),
    \]
    then the semantic map of \(P\) factors through \(V\). Therefore its
    semantic image satisfies
    \[
        \operatorname{Beh}(P)\leq V.
    \]
    Taking the join over all models gives
    \[
        \bigvee_{P\in\mathsf{Mod}(V)}
        \operatorname{Beh}(P)
        \leq
        V.
    \]

    Conversely, the restricted canonical object \(V\) is itself a model of
    \(V\), and
    \[
        \operatorname{Beh}(V)=V
    \]
    by
    Lemma~\ref{lem:semantic-map-of-a-restricted-canonical-subobject}.
    Hence \(V\) occurs among the semantic images entering the join, so
    \[
        V
        \leq
        \bigvee_{P\in\mathsf{Mod}(V)}
        \operatorname{Beh}(P).
    \]
    The two inequalities give
    \[
        V
        =
        \bigvee_{P\in\mathsf{Mod}(V)}
        \operatorname{Beh}(P).
    \]

    Now suppose
    \[
        \mathcal C=\mathsf{Mod}(V).
    \]
    Then
    \[
        V
        =
        \bigvee_{P\in\mathcal C}
        \operatorname{Beh}(P)
        =
        V_{\mathcal C}.
    \]
    Thus the defining local coequational theory is uniquely determined by
    \(\mathcal C\).
\end{proof}

\section{Structured output and algebraic closure}
\label{sec:structured-output-algebraic-closure}

\subsection{Structured output categories}
\label{subsec:structured-output-categories}

This section explains how additional closure operations on behaviours arise
from algebraic structure on the output category. The categorical treatment of
varieties of languages uses such algebraic closure conditions as part of the
bridge between syntactic structure and semantic classes
\cite{Eilenberg1974,AdamekMiliusMyersUrbat2019TOCL,Salamanca2016Monad}.

\begin{definition}[Structured output category]
\label{def:structured-output-category}
    Let \(\mathbb T\) be a single-sorted finite-product theory. A
    \textbf{\(\mathbb T\)-structured output category} is a \(\mathbb T\)-algebra
    object in internal categories, equivalently a finite-product-preserving
    functor
    \[
        \mathbb T\longrightarrow \mathbf{Cat}_{\mathrm{int}}(\mathcal E).
    \]
    We write
    \[
        \mathbb B
    \]
    for the internal category assigned to the generic object of \(\mathbb T\).
    Thus each operation of arity \(n\) in \(\mathbb T\) determines an internal
    functor
    \[
        \theta:\mathbb B^n\to\mathbb B,
    \]
    and the equations of \(\mathbb T\) hold internally.
\end{definition}

\subsection{Pointwise operations on behaviours}
\label{subsec:pointwise-operations-behaviours}

For \(n>0\), put
\[
    \mathbf{Beh}_0^{[n]}
    :=
    \underbrace{
    \mathbf{Beh}_0
    \times_{A_0}
    \cdots
    \times_{A_0}
    \mathbf{Beh}_0
    }_{n}.
\]
A generalized element is a tuple
\[
    (H_1,\dots,H_n)
\]
of residual behaviours
\[
    H_i:\mathbb A/v\to\mathbb B
\]
over the same input object \(v\). For \(n=0\), put
\[
    \mathbf{Beh}_0^{[0]}:=A_0.
\]

Let
\[
    \theta:\mathbb B^n\to\mathbb B
\]
be the internal functor corresponding to an \(n\)-ary operation. For \(n>0\),
write
\[
    \mathbf{Beh}^{(n)}_{\theta,0}
\]
for the underlying object
\[
    \mathbf{Beh}_0^{[n]}
\]
equipped with the structure map
\[
    z_\theta:
    \mathbf{Beh}_0^{[n]}
    \to
    Z=A_0\times B_0
\]
defined by
\[
    z_\theta
    =
    \left\langle
        p_{\mathbf{Beh}}\pi_1,\,
        \theta_0
        \left\langle
            q_{\mathbf{Beh}}\pi_1,
            \dots,
            q_{\mathbf{Beh}}\pi_n
        \right\rangle
    \right\rangle.
\]
Thus
\[
    z_\theta(H_1,\dots,H_n)
    =
    \Bigl(
        p_{\mathbf{Beh}}(H_1),\,
        \theta_0
        \bigl(
            q_{\mathbf{Beh}}H_1,
            \dots,
            q_{\mathbf{Beh}}H_n
        \bigr)
    \Bigr).
\]

The structure map is characterized by the product cone
\[
\begin{tikzcd}[column sep=huge, row sep=large]
    &
    \mathbf{Beh}^{(n)}_{\theta,0}
        \arrow[dl, "{p_{\mathbf{Beh}}\pi_1}"']
        \arrow[d, dashed, "z_\theta"]
        \arrow[dr,
            "{\theta_0
            \langle
                q_{\mathbf{Beh}}\pi_1,
                \dots,
                q_{\mathbf{Beh}}\pi_n
            \rangle}"]
    &
    \\
    A_0
    &
    Z=A_0\times B_0
        \arrow[l, "\pi_{A_0}"]
        \arrow[r, "\pi_{B_0}"']
    &
    B_0 .
\end{tikzcd}
\]

For \(n=0\), if
\[
    \theta:1\to\mathbb B
\]
selects the object
\[
    b_\theta:1\to B_0,
\]
then
\[
    \mathbf{Beh}^{(0)}_{\theta,0}=A_0
\]
is regarded as an object of \(\mathcal E/Z\) by the structure map
\[
    \left\langle
        1_{A_0},
        b_\theta!
    \right\rangle:
    A_0\to A_0\times B_0.
\]
In generalized-element notation,
\[
    v\longmapsto(v,b_\theta).
\]

Notice that
\[
    \mathbf{Beh}^{(n)}_{\theta,0}
\]
is not, in general, the \(n\)-fold product of \(\mathbf{Beh}_0\) in the slice
\[
    \mathcal E/Z.
\]
Its underlying object is first formed as the fibre product
\[
    \mathbf{Beh}_0
    \times_{A_0}
    \cdots
    \times_{A_0}
    \mathbf{Beh}_0
\]
over the common input coordinate. It is then regarded as an object of
\(\mathcal E/Z\) using the output coordinate induced by
\[
    \theta_0:B_0^n\to B_0.
\]
This distinction is essential when closure under \(\mathbb T\)-operations is
expressed as a factorization in \(\mathcal E/Z\).

\begin{definition}[Pointwise structured behaviour]
\label{def:pointwise-structured-behaviour}
    Let
    \[
        \theta:\mathbb B^n\to\mathbb B
    \]
    be an operation of \(\mathbb T\). The induced map
    \[
        \theta_{\mathbf{Beh}}:
        \mathbf{Beh}^{(n)}_{\theta,0}\to\mathbf{Beh}_0
    \]
    in \(\mathcal E/Z\) is defined by
    \[
        \theta_{\mathbf{Beh}}(H_1,\dots,H_n)
        =
        \theta\circ\langle H_1,\dots,H_n\rangle
    \]
    for \(n>0\). For \(n=0\), the nullary operation
    \[
        \theta:1\to\mathbb B
    \]
    induces the constant residual behaviour
    \[
        \mathbb A/v\to 1\xrightarrow{\theta}\mathbb B
    \]
    over each \(v\in A_0\).
\end{definition}

The construction is summarized by
\[
\begin{tikzcd}[column sep=huge, row sep=large]
    \mathbb A/v
        \arrow[r, "{\langle H_1,\dots,H_n\rangle}"]
        \arrow[rr, "{\theta_{\mathbf{Beh}}(H_i)_i}"', bend right=14]
    &
    \mathbb B^n
        \arrow[r, "\theta"]
    &
    \mathbb B .
\end{tikzcd}
\]
On anchor objects and residual arrows this is the pair of diagrams
\[
\begin{tikzcd}[column sep=large, row sep=large]
    1_v
        \arrow[r, maps to]
        \arrow[d, maps to]
    &
    (H_1(1_v),\dots,H_n(1_v))
        \arrow[d, "\theta_0"]
    \\
    \theta_{\mathbf{Beh}}(H_i)_i(1_v)
        \arrow[r, equal]
    &
    \theta_0(H_1(1_v),\dots,H_n(1_v))
\end{tikzcd}
\]
and
\[
\begin{tikzcd}[column sep=large, row sep=large]
    r\circ a
        \arrow[r]
        \arrow[d, maps to]
    &
    r
        \arrow[d, maps to]
    \\
    (H_i(r\circ a))_i
        \arrow[r, "{(H_i(r\circ a\to r))_i}"']
        \arrow[d, "\theta_0"']
    &
    (H_i(r))_i
        \arrow[d, "\theta_0"]
    \\
    \theta(H_i(r\circ a))_i
        \arrow[r, "{\theta_1(H_i(r\circ a\to r))_i}"']
    &
    \theta(H_i(r))_i .
\end{tikzcd}
\]

For a nullary operation \(\theta:1\to\mathbb B\), write
\[
    \theta()
\]
for the Mealy morphism with carrier \(A_0\), input projection \(1_{A_0}\),
constant output \(b_\theta\), transition
\[
    \delta_{\theta()}(a,v)=s_A(a),
\]
and identity output arrow at \(b_\theta\).

\begin{lemma}[Structured operations on strict Mealy morphisms]
\label{lem:structured-operations-semantic}
    For every operation
    \[
        \theta:\mathbb B^n\to\mathbb B,
    \]
    there is an induced \(n\)-ary functor
    \[
        \theta:
        \bigl(\mathsf{Mly}^{\mathrm{str}}_{\mathbb A,\mathbb B}\bigr)^n
        \to
        \mathsf{Mly}^{\mathrm{str}}_{\mathbb A,\mathbb B}
    \]
    for \(n>0\), with the evident nullary object for \(n=0\). The residual
    semantics is compatible with this operation:
    \[
        \beta_{\theta(P_i)_i}
        =
        \theta_{\mathbf{Beh}}(\beta_{P_i})_i.
    \]
\end{lemma}

\begin{proof}
    \leavevmode
    \begin{enumerate}
        \item \textbf{Construct the carrier.}
        For \(n>0\), the carrier is
        \[
            P_\theta:=P_1\times_{A_0}\cdots\times_{A_0}P_n.
        \]
        Its input projection is the common projection to \(A_0\), and its
        output projection is
        \[
            q_\theta(x_1,\dots,x_n)
            =
            \theta_0(q_1x_1,\dots,q_nx_n).
        \]
        This is the pullback-and-output diagram
        \[
        \begin{tikzcd}[column sep=large, row sep=large]
            P_\theta
                \arrow[r, "\pi_i"]
                \arrow[d, "p_\theta"']
                \arrow[dr, phantom, "\lrcorner", very near start]
                \arrow[rr, bend left=18, "{(q_i\pi_i)_i}"]
            &
            P_i
                \arrow[d, "p_i"]
                \arrow[r, "q_i"]
            &
            B_0
            \\
            A_0
                \arrow[r, equal]
            &
            A_0
                \arrow[ur, phantom, "\cdots" description]
            &
        \end{tikzcd}
        \]
        together with
        \[
            q_\theta=\theta_0(q_i\pi_i)_i.
        \]
        For \(n=0\), the carrier is \(A_0\), with input projection
        \(1_{A_0}\) and output projection
        \[
            A_0\to 1\xrightarrow{b_\theta}B_0.
        \]

        \item \textbf{Identify the transition domain.}
        For \(n>0\), there is a canonical pullback identification
        \[
        \begin{aligned}
            A_1\times_{A_0}
            \left(P_1\times_{A_0}\cdots\times_{A_0}P_n\right)
            \cong
            \left(A_1\times_{A_0}P_1\right)
            \times_{A_1}
            \cdots
            \times_{A_1}
            \left(A_1\times_{A_0}P_n\right).
        \end{aligned}
        \]
        Under this identification, a transition datum consists of one input
        arrow \(a\) and compatible states \((x_i)_i\). The compatibility is
        expressed by
        \[
        \begin{tikzcd}[column sep=large, row sep=large]
            A_1\times_{A_0}P_\theta
                \arrow[r]
                \arrow[d]
                \arrow[dr, phantom, "\lrcorner", very near start]
            &
            P_\theta
                \arrow[d]
            \\
            A_1
                \arrow[r, "t_A"']
            &
            A_0 .
        \end{tikzcd}
        \]

        \item \textbf{Define transition and output.}
        For \(n>0\), define
        \[
            \delta_\theta(a,(x_i)_i)
            =
            (\delta_i(a,x_i))_i.
        \]
        The output arrows
        \[
            \omega_i(a,x_i):
            q_i(\delta_i(a,x_i))\to q_i(x_i)
        \]
        form an arrow in \(\mathbb B^n\). Applying the internal functor
        \[
            \theta:\mathbb B^n\to\mathbb B
        \]
        gives a well-typed arrow
        \[
            \theta_1(\omega_i(a,x_i))_i:
            \theta_0(q_i\delta_i(a,x_i))_i
            \to
            \theta_0(q_ix_i)_i.
        \]
        Define
        \[
            \omega_\theta(a,(x_i)_i)
            =
            \theta_1(\omega_i(a,x_i))_i.
        \]
        The construction is summarized by
        \[
        \begin{tikzcd}[column sep=large, row sep=large]
            \theta_0(q_i\delta_i(a,x_i))_i
                \arrow[r, "{\theta_1(\omega_i(a,x_i))_i}"]
                \arrow[d, equal]
            &
            \theta_0(q_ix_i)_i
                \arrow[d, equal]
            \\
            q_\theta\delta_\theta(a,(x_i)_i)
                \arrow[r, "{\omega_\theta(a,(x_i)_i)}"']
            &
            q_\theta(x_i)_i .
        \end{tikzcd}
        \]

        For \(n=0\), the transition is
        \[
            \delta_\theta(a,v)=s_A(a),
        \]
        and the output map is the composite
        \[
            A_1
            \xrightarrow{!}
            1
            \xrightarrow{b_\theta}
            B_0
            \xrightarrow{e_B}
            B_1.
        \]
        In generalized-element notation,
        \[
            \omega_\theta(a,v)=e_B(b_\theta).
        \]

        \item \textbf{Define the operation on strict maps.}
        If
        \[
            f_i:P_i\to Q_i
        \]
        are strict semantic homomorphisms, define
        \[
            \theta(f_i)_i:
            P_1\times_{A_0}\cdots\times_{A_0}P_n
            \to
            Q_1\times_{A_0}\cdots\times_{A_0}Q_n
        \]
        componentwise. Strictness follows from the componentwise transition
        equations and from functoriality of
        \[
            \theta:\mathbb B^n\to\mathbb B.
        \]
        The output square is
        \[
        \begin{tikzcd}[column sep=large, row sep=large]
            A_1\times_{A_0}P_\theta
                \arrow[r, "{1\times\theta(f_i)}"]
                \arrow[d, "\omega_\theta^P"']
            &
            A_1\times_{A_0}Q_\theta
                \arrow[d, "\omega_\theta^Q"]
            \\
            B_1
                \arrow[r, equal]
            &
            B_1 .
        \end{tikzcd}
        \]
        Identities and composition are inherited from the product in the
        carrier slice.

        \item \textbf{Verify the Mealy laws structurally.}
        The transition laws are inherited componentwise. The output laws follow
        from functoriality of
        \[
            \theta:\mathbb B^n\to\mathbb B.
        \]
        In the nullary case the result follows from the identity laws in
        \(\mathbb B\).

        \item \textbf{Compare residual behaviours.}
        For a residual object
        \[
            r:u\to v,
        \]
        one has
        \[
        \begin{aligned}
            q_\theta(\delta_\theta(r,(x_i)_i))
            &=
            \theta_0(q_i(\delta_i(r,x_i)))_i,
        \end{aligned}
        \]
        which is the object value of
        \[
            \theta_{\mathbf{Beh}}(\beta_{P_i}(x_i))_i.
        \]
        The same calculation on residual arrows gives
        \[
            \omega_\theta(a,\delta_\theta(r,(x_i)_i))
            =
            \theta_1(\omega_i(a,\delta_i(r,x_i)))_i.
        \]
        Therefore
        \[
            \beta_{\theta(P_i)_i}
            =
            \theta_{\mathbf{Beh}}(\beta_{P_i})_i.
        \]
    \end{enumerate}
\end{proof}

\begin{lemma}[Structured operations preserve derivatives and canonical output]
\label{lem:operations-preserve-derivatives-output}
    For every operation
    \[
        \theta:\mathbb B^n\to\mathbb B
    \]
    and every arrow
    \[
        a:u\to v
    \]
    of \(\mathbb A\),
    \[
        \partial_a\bigl(\theta_{\mathbf{Beh}}(H_i)_i\bigr)
        =
        \theta_{\mathbf{Beh}}(\partial_aH_i)_i
    \]
    and
    \[
        \omega_{\mathbf{Beh}}
        \bigl(a,\theta_{\mathbf{Beh}}(H_i)_i\bigr)
        =
        \theta_1\bigl(\omega_{\mathbf{Beh}}(a,H_i)\bigr)_i.
    \]
\end{lemma}

\begin{proof}
    \leavevmode
    \begin{enumerate}
        \item \textbf{Use precomposition for derivatives.}
        Since
        \[
            \partial_aH=H\circ a_*,
        \]
        one has
        \[
        \begin{aligned}
            \partial_a\bigl(\theta\circ\langle H_i\rangle_i\bigr)
            &=
            \bigl(\theta\circ\langle H_i\rangle_i\bigr)\circ a_* \\
            &=
            \theta\circ\langle H_i\circ a_*\rangle_i \\
            &=
            \theta_{\mathbf{Beh}}(\partial_aH_i)_i.
        \end{aligned}
        \]
        The commutation is
        \[
        \begin{tikzcd}[column sep=huge, row sep=large]
            \mathbb A/u
                \arrow[r, "a_*"]
                \arrow[dr, "{\langle \partial_aH_i\rangle_i}"']
            &
            \mathbb A/v
                \arrow[d, "{\langle H_i\rangle_i}"]
            \\
            &
            \mathbb B^n
                \arrow[d, "\theta"]
            \\
            &
            \mathbb B .
        \end{tikzcd}
        \]

        \item \textbf{Use functoriality for canonical output.}
        The canonical output is the value on the residual morphism
        \[
            a\to 1_v.
        \]
        Applying
        \[
            \theta\circ\langle H_i\rangle_i
        \]
        gives
        \[
            \theta_1(H_i(a,1_v))_i
            =
            \theta_1(\omega_{\mathbf{Beh}}(a,H_i))_i.
        \]
        Diagrammatically,
        \[
        \begin{tikzcd}[column sep=large, row sep=large]
            (H_i(a))_i
                \arrow[r, "{(H_i(a\to 1_v))_i}"]
                \arrow[d, "\theta_0"']
            &
            (H_i(1_v))_i
                \arrow[d, "\theta_0"]
            \\
            \theta(H_i(a))_i
                \arrow[r, "{\theta_1(H_i(a\to 1_v))_i}"']
            &
            \theta(H_i(1_v))_i .
        \end{tikzcd}
        \]

        \item \textbf{Handle nullary operations.}
        In the nullary case both claims say that a constant functor remains
        constant under residual reindexing and sends residual morphisms to the
        identity arrow selected by the nullary functor.
    \end{enumerate}
\end{proof}

\subsection{\texorpdfstring{\(\mathbb T\)}{T}-coequational theories}
\label{subsec:T-coequational-theories}

\begin{definition}[\(\mathbb T\)-coequational theory]
\label{def:T-coequational-theory}
    Let \(\mathbb T\) be a single-sorted finite-product theory, and let
    \(\mathbb B\) be a \(\mathbb T\)-structured output category.

    A \textbf{local \(\mathbb T\)-coequational theory} is a local
    coequational theory
    \[
        j_V:
        V\hookrightarrow
        \mathbf{Beh}^{\mathrm{can}}_{\mathbb A,\mathbb B}
    \]
    that is closed under the pointwise operation induced by every morphism
    \[
        \theta:n\to1
    \]
    of \(\mathbb T\).

    For \(n>0\), put
    \[
        V^{[n]}
        :=
        \underbrace{
        V\times_{A_0}\cdots\times_{A_0}V
        }_{n}.
    \]
    This is the \(n\)-fold product of
    \[
        p_V:V\to A_0
    \]
    in the slice \(\mathcal E/A_0\).

    For an operation
    \[
        \theta:n\to1
    \]
    whose interpretation is the internal functor
    \[
        \theta:\mathbb B^n\to\mathbb B,
    \]
    regard \(V^{[n]}\) as an object of
    \[
        \mathcal E/Z,
        \qquad
        Z=A_0\times B_0,
    \]
    by the structure map
    \[
        z_{\theta,V}
        :=
        \left\langle
            p_V\pi_1,\,
            \theta_0
            \left\langle
                q_V\pi_1,
                \dots,
                q_V\pi_n
            \right\rangle
        \right\rangle.
    \]
    The resulting object of \(\mathcal E/Z\) is denoted
    \[
        V^{(n)}_\theta.
    \]

    Thus, in generalized-element notation,
    \[
        z_{\theta,V}(H_1,\dots,H_n)
        =
        \Bigl(
            p_V(H_1),\,
            \theta_0
            \bigl(
                q_V(H_1),
                \dots,
                q_V(H_n)
            \bigr)
        \Bigr).
    \]

    For a nullary theory operation
    \[
        \theta:0\to1,
    \]
    we use the same symbol
    \[
        \theta:1\to\mathbb B
    \]
    for its interpretation in the \(\mathbb T\)-structured internal category.
    Let
    \[
        b_\theta:1\to B_0
    \]
    be its selected output object. Put
    \[
        V^{(0)}_\theta:=A_0,
    \]
    regarded over \(Z\) by
    \[
        \left\langle
            1_{A_0},
            b_\theta!
        \right\rangle:
        A_0\to A_0\times B_0.
    \]

    Closure under a positive-arity operation
    \[
        \theta:n\to1
    \]
    means that the pointwise behaviour operation restricts to \(V\):
    there is a unique map
    \[
        \theta_V:
        V^{(n)}_\theta\to V
    \]
    in \(\mathcal E/Z\) satisfying
    \[
        j_V\theta_V
        =
        \theta_{\mathbf{Beh}}j_V^{[n]},
    \]
    where
    \[
        j_V^{[n]}
        :=
        j_V\times_{A_0}\cdots\times_{A_0}j_V:
        V^{[n]}
        \to
        \mathbf{Beh}_0^{[n]}.
    \]

    Equivalently, after forgetting the \(B_0\)-coordinate,
    \[
        \theta_V:V^{[n]}\to V
    \]
    is a map in \(\mathcal E/A_0\) satisfying
    \[
        q_V\theta_V
        =
        \theta_0
        \left\langle
            q_V\pi_1,
            \dots,
            q_V\pi_n
        \right\rangle.
    \]

    In the nullary case, closure means that there is a unique map
    \[
        \theta_V:A_0\to V
    \]
    satisfying
    \[
        j_V\theta_V
        =
        \theta_{\mathbf{Beh}},
    \]
    equivalently
    \[
        p_V\theta_V=1_{A_0},
        \qquad
        q_V\theta_V=b_\theta!.
    \]

    The proposition below shows that the resulting operations
    \[
        \theta_V:V^{[n]}\to V
    \]
    form a \(\mathbb T\)-algebra structure on
    \[
        p_V:V\to A_0
    \]
    in the slice \(\mathcal E/A_0\), and that the output map preserves this
    structure.
\end{definition}

The closure condition is displayed by
\[
\begin{tikzcd}[column sep=large, row sep=large]
    V^{(n)}_\theta
        \arrow[rr, "{\theta_{\mathbf{Beh}}}"]
        \arrow[dr, dashed, "\theta_V"']
    &&
    \mathbf{Beh}_0
    \\
    &
    V
        \arrow[ur, hook, "j_V"']
    &
\end{tikzcd}
\qquad
\text{in }\mathcal E/Z.
\]
The compatibility between algebraic closure and derivative closure is
summarized by the cube-like diagram
\[
\begin{tikzcd}[column sep=huge, row sep=large]
    A_1\times_{A_0}V^{(n)}_\theta
        \arrow[r, "{1\times \theta_V}"]
        \arrow[d, "\delta_{V^{(n)}_\theta}"']
        \arrow[dr, phantom, "\circlearrowleft", very near start]
    &
    A_1\times_{A_0}V
        \arrow[d, "\delta_V"]
    \\
    V^{(n)}_\theta
        \arrow[r, "\theta_V"']
        \arrow[d, hook]
    &
    V
        \arrow[d, hook, "j_V"]
    \\
    \mathbf{Beh}^{(n)}_{\theta,0}
        \arrow[r, "\theta_{\mathbf{Beh}}"']
    &
    \mathbf{Beh}_0 .
\end{tikzcd}
\]

\begin{proposition}[Inherited structured operations]
\label{prop:inherited-structured-operations}
    Let
    \[
        V\hookrightarrow
        \mathbf{Beh}^{\mathrm{can}}_{\mathbb A,\mathbb B}
    \]
    be a local \(\mathbb T\)-coequational theory, and let
    \[
        \theta:n\to1
    \]
    be a morphism of the finite-product theory \(\mathbb T\), interpreted as
    an internal functor
    \[
        \theta:\mathbb B^n\to\mathbb B.
    \]

    For \(n>0\), the object
    \[
        V^{(n)}_\theta
    \]
    carries the Mealy structure of
    \[
        \theta(V,\dots,V).
    \]
    Explicitly,
    \[
        \delta_{V^{(n)}_\theta}
        \bigl(
            a,(H_1,\dots,H_n)
        \bigr)
        =
        \bigl(
            \partial_aH_1,
            \dots,
            \partial_aH_n
        \bigr),
    \]
    and
    \[
        \omega_{V^{(n)}_\theta}
        \bigl(
            a,(H_1,\dots,H_n)
        \bigr)
        =
        \theta_1
        \bigl(
            \omega_V(a,H_1),
            \dots,
            \omega_V(a,H_n)
        \bigr).
    \]

    The closure factorization
    \[
        \theta_V:
        V^{(n)}_\theta\to V
    \]
    is a strict semantic homomorphism. The analogous assertion holds for
    nullary operations.

    After forgetting the \(B_0\)-coordinate, the family
    \[
        \theta_V:V^{[n]}\to V
    \]
    preserves projections and substitution. Consequently it defines a
    \(\mathbb T\)-algebra structure on
    \[
        p_V:V\to A_0
    \]
    in the slice \(\mathcal E/A_0\).

    Moreover,
    \[
        \langle p_V,q_V\rangle:
        V\to A_0\times B_0
    \]
    is a homomorphism in \(\mathcal E/A_0\) from this
    \(\mathbb T\)-algebra to the fibrewise \(\mathbb T\)-algebra
    \[
        A_0\times B_0\to A_0
    \]
    induced by the \(\mathbb T\)-algebra structure on \(B_0\).
\end{proposition}

\begin{proof}
    Suppose first that \(n>0\).

    By Lemma~\ref{lem:structured-operations-semantic}, the fibre product
    \[
        V^{[n]}
        =
        V\times_{A_0}\cdots\times_{A_0}V
    \]
    carries the Mealy structure
    \[
        \theta(V,\dots,V)
    \]
    when equipped with the \(\theta\)-induced output coordinate. Thus its
    carrier over \(Z\) is \(V^{(n)}_\theta\), its transition is
    \[
        \delta_{V^{(n)}_\theta}
        \bigl(
            a,(H_1,\dots,H_n)
        \bigr)
        =
        \bigl(
            \partial_aH_1,
            \dots,
            \partial_aH_n
        \bigr),
    \]
    and its output is
    \[
        \omega_{V^{(n)}_\theta}
        \bigl(
            a,(H_1,\dots,H_n)
        \bigr)
        =
        \theta_1
        \bigl(
            \omega_V(a,H_1),
            \dots,
            \omega_V(a,H_n)
        \bigr).
    \]

    Since \(V\) is closed under \(\theta\), there is a unique map
    \[
        \theta_V:
        V^{(n)}_\theta\to V
    \]
    in \(\mathcal E/Z\) satisfying
    \[
        j_V\theta_V
        =
        \theta_{\mathbf{Beh}}j_V^{[n]}.
    \]

    We prove that \(\theta_V\) preserves transition. Let
    \[
        a:u\to v
    \]
    and let
    \[
        (H_1,\dots,H_n)\in V^{[n]}
    \]
    lie over \(v\). After composing with \(j_V\),
    \[
    \begin{aligned}
        &
        j_V\theta_V
        \delta_{V^{(n)}_\theta}
        \bigl(
            a,(H_1,\dots,H_n)
        \bigr)
        \\
        &\qquad=
        \theta_{\mathbf{Beh}}
        \bigl(
            j_V\partial_aH_1,
            \dots,
            j_V\partial_aH_n
        \bigr)
        \\
        &\qquad=
        \theta_{\mathbf{Beh}}
        \bigl(
            \partial_a(j_VH_1),
            \dots,
            \partial_a(j_VH_n)
        \bigr)
        \\
        &\qquad=
        \partial_a
        \bigl(
            \theta_{\mathbf{Beh}}
            (j_VH_1,\dots,j_VH_n)
        \bigr)
        \\
        &\qquad=
        \partial_a
        \bigl(
            j_V\theta_V(H_1,\dots,H_n)
        \bigr)
        \\
        &\qquad=
        j_V\delta_V
        \bigl(
            a,
            \theta_V(H_1,\dots,H_n)
        \bigr).
    \end{aligned}
    \]
    Here the third equality is
    Lemma~\ref{lem:operations-preserve-derivatives-output}. Since \(j_V\) is
    monic,
    \[
        \theta_V\delta_{V^{(n)}_\theta}
        =
        \delta_V\widehat{\theta_V}.
    \]

    The map \(\theta_V\) also preserves output. Indeed,
    \[
    \begin{aligned}
        &
        \omega_V
        \bigl(
            a,
            \theta_V(H_1,\dots,H_n)
        \bigr)
        \\
        &\qquad=
        \omega_{\mathbf{Beh}}
        \bigl(
            a,
            j_V\theta_V(H_1,\dots,H_n)
        \bigr)
        \\
        &\qquad=
        \omega_{\mathbf{Beh}}
        \bigl(
            a,
            \theta_{\mathbf{Beh}}
            (j_VH_1,\dots,j_VH_n)
        \bigr)
        \\
        &\qquad=
        \theta_1
        \bigl(
            \omega_{\mathbf{Beh}}(a,j_VH_1),
            \dots,
            \omega_{\mathbf{Beh}}(a,j_VH_n)
        \bigr)
        \\
        &\qquad=
        \theta_1
        \bigl(
            \omega_V(a,H_1),
            \dots,
            \omega_V(a,H_n)
        \bigr)
        \\
        &\qquad=
        \omega_{V^{(n)}_\theta}
        \bigl(
            a,(H_1,\dots,H_n)
        \bigr).
    \end{aligned}
    \]
    Thus \(\theta_V\) is a strict semantic homomorphism.

    Now let
    \[
        \theta:0\to1
    \]
    be nullary. Then
    \[
        V^{(0)}_\theta=A_0
    \]
    carries the nullary Mealy morphism \(\theta()\). Closure gives a map
    \[
        \theta_V:A_0\to V
    \]
    satisfying
    \[
        j_V\theta_V=\theta_{\mathbf{Beh}}.
    \]
    The nullary cases of
    Lemma~\ref{lem:operations-preserve-derivatives-output}
    show, by the same monicity argument, that \(\theta_V\) preserves
    transition and output. Hence it is strict.

    We next verify that the underlying maps in \(\mathcal E/A_0\)
    preserve the finite-product-theory structure.

    Let
    \[
        \pi_i:n\to1
    \]
    be the \(i\)-th product projection in \(\mathbb T\). The corresponding
    pointwise operation on behaviours is the ordinary projection
    \[
        \pi_{i,\mathbf{Beh}}:
        \mathbf{Beh}_0^{[n]}\to\mathbf{Beh}_0.
    \]
    The ordinary projection
    \[
        \pi_i:V^{[n]}\to V
    \]
    satisfies
    \[
        j_V\pi_i
        =
        \pi_{i,\mathbf{Beh}}j_V^{[n]}.
    \]
    By uniqueness of factorization through the monomorphism \(j_V\),
    \[
        \pi_{i,V}=\pi_i.
    \]
    Hence the restricted operations preserve the distinguished projections.

    Let
    \[
        \theta:k\to1
    \]
    and
    \[
        \varphi_i:n\to1
        \qquad
        (1\leq i\leq k)
    \]
    be morphisms in \(\mathbb T\), and put
    \[
        \psi
        :=
        \theta
        \circ
        \left\langle
            \varphi_1,\dots,\varphi_k
        \right\rangle:
        n\to1.
    \]

    In \(\mathcal E/A_0\), the restricted operations determine a pairing
    \[
        \left\langle
            \varphi_{1,V},
            \dots,
            \varphi_{k,V}
        \right\rangle:
        V^{[n]}\to V^{[k]}.
    \]
    Composing with \(j_V\), one obtains
    \[
    \begin{aligned}
        &
        j_V
        \theta_V
        \left\langle
            \varphi_{1,V},
            \dots,
            \varphi_{k,V}
        \right\rangle
        \\
        &\qquad=
        \theta_{\mathbf{Beh}}
        j_V^{[k]}
        \left\langle
            \varphi_{1,V},
            \dots,
            \varphi_{k,V}
        \right\rangle
        \\
        &\qquad=
        \theta_{\mathbf{Beh}}
        \left\langle
            j_V\varphi_{1,V},
            \dots,
            j_V\varphi_{k,V}
        \right\rangle
        \\
        &\qquad=
        \theta_{\mathbf{Beh}}
        \left\langle
            \varphi_{1,\mathbf{Beh}}j_V^{[n]},
            \dots,
            \varphi_{k,\mathbf{Beh}}j_V^{[n]}
        \right\rangle
        \\
        &\qquad=
        \theta_{\mathbf{Beh}}
        \left\langle
            \varphi_{1,\mathbf{Beh}},
            \dots,
            \varphi_{k,\mathbf{Beh}}
        \right\rangle
        j_V^{[n]}
        \\
        &\qquad=
        \psi_{\mathbf{Beh}}j_V^{[n]}
        \\
        &\qquad=
        j_V\psi_V.
    \end{aligned}
    \]
    Since \(j_V\) is monic,
    \[
        \theta_V
        \left\langle
            \varphi_{1,V},
            \dots,
            \varphi_{k,V}
        \right\rangle
        =
        \psi_V.
    \]

    Thus the restricted operations preserve projections and substitution.
    These are precisely the structural requirements for an interpretation of
    the single-sorted finite-product theory \(\mathbb T\). Therefore they
    define a \(\mathbb T\)-algebra structure on
    \[
        p_V:V\to A_0
    \]
    in \(\mathcal E/A_0\).

    Finally, because
    \[
        \theta_V:V^{(n)}_\theta\to V
    \]
    is a map in \(\mathcal E/Z\), one has, for every positive-arity
    operation,
    \[
        q_V\theta_V
        =
        \theta_0
        \left\langle
            q_V\pi_1,
            \dots,
            q_V\pi_n
        \right\rangle.
    \]
    For a nullary operation,
    \[
        q_V\theta_V=b_\theta!.
    \]
    Hence
    \[
        \langle p_V,q_V\rangle:
        V\to A_0\times B_0
    \]
    preserves every operation of \(\mathbb T\) in the slice
    \(\mathcal E/A_0\). It is therefore a homomorphism to the fibrewise
    \(\mathbb T\)-algebra
    \[
        A_0\times B_0\to A_0.
    \]
\end{proof}

\subsection{Semantic images and structured operations}
\label{subsec:semantic-images-structured-operations}

\begin{lemma}[Semantic images under structured operations]
\label{lem:semantic-images-structured-operations}
    Let
    \[
        \theta:\mathbb B^n\to\mathbb B
    \]
    be an operation of \(\mathbb T\), and let
    \[
        P_1,\dots,P_n
    \]
    be Mealy morphisms. For \(n>0\), the semantic quotient maps
    \[
        P_i\twoheadrightarrow\operatorname{Beh}(P_i)
    \]
    induce a regular epimorphism
    \[
        \theta(P_1,\dots,P_n)
        \twoheadrightarrow
        \operatorname{Beh}(P_1)\times_{A_0}\cdots\times_{A_0}
        \operatorname{Beh}(P_n),
    \]
    where the codomain is regarded over
    \[
        Z=A_0\times B_0
    \]
    by the \(\theta\)-induced output coordinate. Moreover,
    \[
        \operatorname{Im}
        \left(
        \theta_{\mathbf{Beh}}:
        \operatorname{Beh}(P_1)\times_{A_0}\cdots\times_{A_0}
        \operatorname{Beh}(P_n)
        \to
        \mathbf{Beh}_0
        \right)
        =
        \operatorname{Beh}(\theta(P_1,\dots,P_n)).
    \]
    For \(n=0\), the image of
    \[
        A_0\to\mathbf{Beh}_0
    \]
    induced by the constant residual behaviour is
    \[
        \operatorname{Beh}(\theta()).
    \]
\end{lemma}

\begin{proof}
    \leavevmode
    \begin{enumerate}
        \item \textbf{Form the product of semantic quotients.}
        For \(n>0\), the semantic quotients
        \[
            P_i\twoheadrightarrow\operatorname{Beh}(P_i)
        \]
        are regular epimorphisms in the carrier slice \(\mathcal E/Z\), hence
        in particular their underlying maps over \(A_0\) are regular
        epimorphisms. Since regular epimorphisms are stable under pullback and
        finite products in slices of a Grothendieck topos, their fibre product
        over \(A_0\) is a regular epimorphism:
        \[
            P_1\times_{A_0}\cdots\times_{A_0}P_n
            \twoheadrightarrow
            \operatorname{Beh}(P_1)\times_{A_0}\cdots\times_{A_0}
            \operatorname{Beh}(P_n).
        \]
        The source is the carrier of
        \[
            \theta(P_1,\dots,P_n),
        \]
        and the target is regarded as an object of \(\mathcal E/Z\) by the
        \(\theta\)-induced coordinate
        \[
            (H_1,\dots,H_n)
            \longmapsto
            \Bigl(
                p_{\mathbf{Beh}}(H_1),
                \theta_0(q_{\mathbf{Beh}}H_1,\dots,q_{\mathbf{Beh}}H_n)
            \Bigr).
        \]
        With this coordinate, the displayed map is a regular epimorphism in
        \(\mathcal E/Z\).

        \item \textbf{Compare the semantic composites.}
        By Lemma~\ref{lem:structured-operations-semantic},
        \[
            \beta_{\theta(P_i)_i}
            =
            \theta_{\mathbf{Beh}}(\beta_{P_i})_i.
        \]
        Hence the semantic map of
        \[
            \theta(P_1,\dots,P_n)
        \]
        is the composite
        \[
        \begin{tikzcd}[column sep=large, row sep=large]
            \theta(P_i)_i
                \arrow[r, two heads]
                \arrow[rr, bend left=18, "\beta_{\theta(P_i)_i}"]
            &
            \prod_{A_0}\operatorname{Beh}(P_i)
                \arrow[r, "\theta_{\mathbf{Beh}}"]
            &
            \mathbf{Beh}_0 .
        \end{tikzcd}
        \]

        \item \textbf{Compare regular images.}
        Precomposition with a regular epimorphism does not change the regular
        image. Therefore the image of
        \[
            \theta_{\mathbf{Beh}}:
            \prod_{A_0}\operatorname{Beh}(P_i)
            \to
            \mathbf{Beh}_0
        \]
        is exactly the semantic image of
        \[
            \theta(P_1,\dots,P_n).
        \]
        Hence
        \[
            \operatorname{Im}(\theta_{\mathbf{Beh}})
            =
            \operatorname{Beh}(\theta(P_1,\dots,P_n)).
        \]

        \item \textbf{Handle nullary operations.}
        For \(n=0\), the semantic map of \(\theta()\) is the map
        \[
            A_0\to\mathbf{Beh}_0
        \]
        assigning to \(v\) the constant residual behaviour
        \[
            \mathbb A/v\to 1\xrightarrow{\theta}\mathbb B.
        \]
        Its regular image is, by definition,
        \[
            \operatorname{Beh}(\theta()).
        \]
    \end{enumerate}
\end{proof}

\subsection{\texorpdfstring{\(\mathbb T\)}{T}-coequational closure}
\label{subsec:T-coequational-closure}

\begin{theorem}[\(\mathbb T\)-coequational closure]
\label{thm:T-coequational-closure}
    Let
    \[
        V\hookrightarrow\mathbf{Beh}^{\mathrm{can}}
    \]
    be a local \(\mathbb T\)-coequational theory. Then
    \[
        \mathsf{Mod}(V)
    \]
    is closed under strict homomorphic preimages, regular-epimorphic quotients,
    small coproducts, and the operations induced by \(\mathbb T\).
\end{theorem}

\begin{proof}
    \leavevmode
    \begin{enumerate}
        \item \textbf{Use local closure.}
        Closure under strict homomorphic preimages, regular-epimorphic
        quotients, and small coproducts is
        Theorem~\ref{thm:local-coequational-closure}.

        \item \textbf{Use semantic compatibility.}
        Let
        \[
            P_1,\dots,P_n\in\mathsf{Mod}(V).
        \]
        Choose strict factorizations
        \[
            \bar\beta_{P_i}:P_i\to V
        \]
        with
        \[
            j_V\bar\beta_{P_i}=\beta_{P_i}.
        \]
        These induce a map from the carrier of
        \[
            \theta(P_1,\dots,P_n)
        \]
        to \(V^{(n)}_\theta\):
        \[
        \begin{tikzcd}[column sep=large, row sep=large]
            P_1\times_{A_0}\cdots\times_{A_0}P_n
                \arrow[r, "{(\bar\beta_{P_i})_i}"]
                \arrow[dr, "{(\beta_{P_i})_i}"']
            &
            V^{(n)}_\theta
                \arrow[d, hook]
            \\
            &
            \mathbf{Beh}^{(n)}_{\theta,0}.
        \end{tikzcd}
        \]

        \item \textbf{Use algebraic closure of \(V\).}
        Since \(V\) is closed under \(\theta\), there is a map
        \[
            \theta_V:V^{(n)}_\theta\to V
        \]
        with
        \[
            j_V\theta_V
            =
            \theta_{\mathbf{Beh}}\big|_{V^{(n)}_\theta}.
        \]
        By Lemma~\ref{lem:structured-operations-semantic},
        \[
            \beta_{\theta(P_i)_i}
            =
            \theta_{\mathbf{Beh}}(\beta_{P_i})_i.
        \]
        Hence
        \[
            \beta_{\theta(P_i)_i}
        \]
        factors through \(V\), so
        \[
            \theta(P_1,\dots,P_n)\in\mathsf{Mod}(V).
        \]
        The factorization is
        \[
        \begin{tikzcd}[column sep=large, row sep=large]
            \theta(P_i)_i
                \arrow[r, "{(\bar\beta_{P_i})_i}"]
                \arrow[rrr, bend left=18, "\beta_{\theta(P_i)_i}"]
            &
            V^{(n)}_\theta
                \arrow[r, "\theta_V"]
            &
            V
                \arrow[r, hook, "j_V"]
            &
            \mathbf{Beh}^{\mathrm{can}} .
        \end{tikzcd}
        \]

        \item \textbf{Handle nullary operations.}
        The nullary case is the same argument using
        \[
            V^{(0)}_\theta=A_0
        \]
        and the closure map
        \[
            A_0\to V.
        \]
    \end{enumerate}
\end{proof}

\subsection{The structured Birkhoff theorem}
\label{subsec:T-coequational-birkhoff}

\begin{theorem}[\(\mathbb T\)-coequational Birkhoff theorem]
\label{thm:T-coequational-birkhoff}
    Let \(\mathbb B\) be a \(\mathbb T\)-structured output category, and let
    \[
        \mathcal C
        \subseteq
        \mathsf{Mly}^{\mathrm{str}}_{\mathbb A,\mathbb B}
    \]
    be a replete full subcategory. Then \(\mathcal C\) is of the form
    \[
        \mathcal C=\mathsf{Mod}(V)
    \]
    for a local \(\mathbb T\)-coequational theory
    \[
        V\hookrightarrow
        \mathbf{Beh}^{\mathrm{can}}_{\mathbb A,\mathbb B}
    \]
    if and only if \(\mathcal C\) is closed under:

    \begin{enumerate}
        \item strict homomorphic preimages;
        \item regular-epimorphic quotients;
        \item small coproducts;
        \item the operations induced by \(\mathbb T\).
    \end{enumerate}
\end{theorem}

\begin{proof}
    \leavevmode
    \begin{enumerate}
        \item \textbf{Prove necessity.}

        Suppose
        \[
            \mathcal C=\mathsf{Mod}(V)
        \]
        for a local \(\mathbb T\)-coequational theory \(V\).

        Closure under strict homomorphic preimages, regular-epimorphic
        quotients, and small coproducts follows from
        Theorem~\ref{thm:local-coequational-closure}. Closure under the
        operations induced by \(\mathbb T\) follows from
        Theorem~\ref{thm:T-coequational-closure}.

        \item \textbf{Construct the underlying local theory.}

        Conversely, suppose that \(\mathcal C\) has the four displayed
        closure properties. By the local coequational Birkhoff theorem,
        \[
            V_{\mathcal C}
            :=
            \bigvee_{P\in\mathcal C}
            \operatorname{Beh}(P)
            \hookrightarrow
            \mathbf{Beh}^{\mathrm{can}}_{\mathbb A,\mathbb B}
        \]
        is a local coequational theory and
        \[
            \mathcal C
            =
            \mathsf{Mod}(V_{\mathcal C}).
        \]

        Moreover,
        Lemma~\ref{lem:representing-theory-belongs-to-class} gives
        \[
            V_{\mathcal C}\in\mathcal C.
        \]

        It remains to prove that \(V_{\mathcal C}\) is closed under every
        pointwise operation induced by \(\mathbb T\).

        \item \textbf{Treat positive-arity operations.}

        Let
        \[
            \theta:\mathbb B^n\to\mathbb B
        \]
        be an operation with \(n>0\).

        Since
        \[
            V_{\mathcal C}\in\mathcal C
        \]
        and \(\mathcal C\) is closed under the operation induced by
        \(\theta\),
        \[
            \theta(
                V_{\mathcal C},
                \dots,
                V_{\mathcal C}
            )
            \in
            \mathcal C.
        \]
        Using
        \[
            \mathcal C
            =
            \mathsf{Mod}(V_{\mathcal C}),
        \]
        this object models \(V_{\mathcal C}\). Hence its semantic map admits
        a unique factorization
        \[
            \bar\beta_\theta:
            \theta(
                V_{\mathcal C},
                \dots,
                V_{\mathcal C}
            )
            \longrightarrow
            V_{\mathcal C}
        \]
        satisfying
        \[
            j_{V_{\mathcal C}}\bar\beta_\theta
            =
            \beta_{
                \theta(
                    V_{\mathcal C},
                    \dots,
                    V_{\mathcal C}
                )
            }.
        \]
        By Lemma~\ref{lem:model-factorizations-strict},
        \(\bar\beta_\theta\) is a strict semantic homomorphism.

        The carrier of
        \[
            \theta(
                V_{\mathcal C},
                \dots,
                V_{\mathcal C}
            )
        \]
        is
        \[
            (V_{\mathcal C})^{(n)}_\theta.
        \]
        By Lemma~\ref{lem:structured-operations-semantic},
        \[
        \begin{aligned}
            \beta_{
                \theta(
                    V_{\mathcal C},
                    \dots,
                    V_{\mathcal C}
                )
            }
            &=
            \theta_{\mathbf{Beh}}
            \bigl(
                \beta_{V_{\mathcal C}},
                \dots,
                \beta_{V_{\mathcal C}}
            \bigr)
            \\
            &=
            \theta_{\mathbf{Beh}}
            j_{V_{\mathcal C}}^{[n]}.
        \end{aligned}
        \]
        The second equality uses
        Lemma~\ref{lem:semantic-map-of-a-restricted-canonical-subobject}:
        \[
            \beta_{V_{\mathcal C}}
            =
            j_{V_{\mathcal C}}.
        \]

        Therefore
        \[
            j_{V_{\mathcal C}}\bar\beta_\theta
            =
            \theta_{\mathbf{Beh}}
            j_{V_{\mathcal C}}^{[n]}.
        \]
        Define
        \[
            \theta_{V_{\mathcal C}}
            :=
            \bar\beta_\theta:
            (V_{\mathcal C})^{(n)}_\theta
            \longrightarrow
            V_{\mathcal C}.
        \]
        Then
        \[
            j_{V_{\mathcal C}}\theta_{V_{\mathcal C}}
            =
            \theta_{\mathbf{Beh}}
            j_{V_{\mathcal C}}^{[n]},
        \]
        which is exactly the required closure factorization.

        The construction is displayed by
        \[
        \begin{tikzcd}[column sep=large, row sep=large]
            (V_{\mathcal C})^{(n)}_\theta
                \arrow[rr,
                    "{\theta_{\mathbf{Beh}}
                    j_{V_{\mathcal C}}^{[n]}}"]
                \arrow[dr, "\theta_{V_{\mathcal C}}"']
            &&
            \mathbf{Beh}_0
            \\
            &
            V_{\mathcal C}
                \arrow[ur, hook, "j_{V_{\mathcal C}}"']
            &
        \end{tikzcd}
        \]

        \item \textbf{Treat nullary operations.}

        Let
        \[
            \theta:0\to1
        \]
        be a nullary operation, interpreted as an internal functor
        \[
            \theta:1\to\mathbb B.
        \]
        Closure of \(\mathcal C\) under this nullary operation gives
        \[
            \theta()\in\mathcal C.
        \]
        Therefore
        \[
            \theta()\models V_{\mathcal C}.
        \]
        Its semantic map admits a unique strict model factorization
        \[
            \theta_{V_{\mathcal C}}:
            A_0\longrightarrow V_{\mathcal C}
        \]
        satisfying
        \[
            j_{V_{\mathcal C}}\theta_{V_{\mathcal C}}
            =
            \beta_{\theta()}.
        \]

        By the nullary case of
        Lemma~\ref{lem:structured-operations-semantic}, the semantic map of
        \(\theta()\) is the constant residual-behaviour map
        \[
            \theta_{\mathbf{Beh}}:
            A_0\to\mathbf{Beh}_0.
        \]
        Hence
        \[
            j_{V_{\mathcal C}}\theta_{V_{\mathcal C}}
            =
            \theta_{\mathbf{Beh}}.
        \]
        This is precisely the required nullary closure factorization.

        \item \textbf{Conclude.}

        Thus \(V_{\mathcal C}\) is closed under every pointwise operation
        induced by \(\mathbb T\). Since it is already a local coequational
        theory, it is a local \(\mathbb T\)-coequational theory.

        The underlying local Birkhoff theorem gives
        \[
            \mathcal C
            =
            \mathsf{Mod}(V_{\mathcal C}),
        \]
        completing the proof.
    \end{enumerate}
\end{proof}

\begin{corollary}[Uniqueness of the defining local
\(\mathbb T\)-coequational theory]
\label{cor:uniqueness-defining-local-T-coequational-theory}
    Let \(\mathbb B\) be a \(\mathbb T\)-structured output category, and let
    \[
        V\hookrightarrow
        \mathbf{Beh}^{\mathrm{can}}_{\mathbb A,\mathbb B}
    \]
    be a local \(\mathbb T\)-coequational theory. Then
    \[
        V
        =
        \bigvee_{P\in\mathsf{Mod}(V)}
        \operatorname{Beh}(P)
    \]
    as a subobject of the canonical behaviour machine.

    Consequently, if a replete full subcategory
    \[
        \mathcal C
        \subseteq
        \mathsf{Mly}^{\mathrm{str}}_{\mathbb A,\mathbb B}
    \]
    is definable by a local \(\mathbb T\)-coequational theory, then that theory
    is unique and is
    \[
        V_{\mathcal C}
        =
        \bigvee_{P\in\mathcal C}
        \operatorname{Beh}(P).
    \]
\end{corollary}

\begin{proof}
    Forgetting the additional \(\mathbb T\)-closure, a local
    \(\mathbb T\)-coequational theory is a local coequational theory.
    Therefore
    Corollary~\ref{cor:uniqueness-defining-local-coequational-theory}
    gives
    \[
        V
        =
        \bigvee_{P\in\mathsf{Mod}(V)}
        \operatorname{Beh}(P).
    \]

    Now suppose
    \[
        \mathcal C=\mathsf{Mod}(V)
    \]
    for a local \(\mathbb T\)-coequational theory \(V\). Then the underlying
    local coequational subobject is necessarily
    \[
        V
        =
        \bigvee_{P\in\mathcal C}
        \operatorname{Beh}(P).
    \]

    It remains only to observe that the \(\mathbb T\)-structure on this
    subobject is not additional independent data. For every operation
    \[
        \theta:n\to1
    \]
    of \(\mathbb T\), the induced map
    \[
        \theta_V:
        V^{(n)}_\theta\to V
    \]
    is uniquely determined by the factorization equation
    \[
        j_V\theta_V
        =
        \theta_{\mathbf{Beh}}j_V^{[n]},
    \]
    because \(j_V\) is monic.

    Hence both the underlying subobject and all its restricted
    \(\mathbb T\)-operations are uniquely determined by \(\mathcal C\).
    Therefore the defining local \(\mathbb T\)-coequational theory is unique.
\end{proof}

\section{Global input reindexing}
\label{sec:global-input-reindexing}

\subsection{Reindexing residual behaviours along input functors}
\label{subsec:input-reindexing-behaviours}

This section explains how local coequational theories vary with the input
category. This is the categorical source of closure under inverse morphic
preimage in the classical one-object setting, a closure operation appearing in
Eilenberg-type correspondences for varieties of languages
\cite{Eilenberg1974,AdamekMiliusMyersUrbat2019TOCL}.

Fix \(\mathbb B\), and, when desired, a finite-product theory \(\mathbb T\) for
which \(\mathbb B\) is a \(\mathbb T\)-structured output category. Let
\[
    \mathsf{Inp}
\]
be a chosen category of internal input categories and internal functors. For
every
\[
    \mathbb A\in\mathsf{Inp},
\]
there is a canonical behaviour object
\[
    \mathbf{Beh}_{\mathbb A,\mathbb B,0}
\]
and a strict Mealy category
\[
    \mathsf{Mly}^{\mathrm{str}}_{\mathbb A,\mathbb B}.
\]
We write
\[
    Z_{\mathbb A}:=A_0\times B_0.
\]

Let
\[
    F:\mathbb A'\to\mathbb A
\]
be an internal functor. Let
\[
    \mathbb R_{\mathbb A'}
\]
be the vertical residual category over \(A'_0\), and let
\[
    F_0^*\mathbb R_{\mathbb A}
\]
be the pullback to \(A'_0\) of the vertical residual category over \(A_0\).

\begin{definition}[Internal residual-slice functor]
\label{def:internal-residual-slice-functor}
    Application of \(F\) to residual arrows and residual triangles defines an
    internal functor over \(A'_0\)
    \[
        F_{/}:
        \mathbb R_{\mathbb A'}
        \longrightarrow
        F_0^*\mathbb R_{\mathbb A}.
    \]

    On residual objects it is
    \[
        r':u'\to v'
        \quad\longmapsto\quad
        F_1(r'):F_0(u')\to F_0(v').
    \]
    On residual arrows represented by a triangle
    \[
        w'\xrightarrow{a'}u'\xrightarrow{r'}v',
    \]
    it sends
    \[
        r'\circ a'\to r'
    \]
    to
    \[
        F_1(r')\circ F_1(a')
        \longrightarrow
        F_1(r').
    \]
\end{definition}

Functoriality follows from
\[
    F_1(1_{v'})=1_{F_0(v')}
\]
and
\[
    F_1(r'\circ a')
    =
    F_1(r')\circ F_1(a').
\]

Precomposition with \(F_{/}\) defines an internal functor over \(A'_0\)
\[
    (F_{/})^*:
    F_0^*
    \mathbf{Beh}^{\mathrm{vert}}_{\mathbb A,\mathbb B}
    \longrightarrow
    \mathbf{Beh}^{\mathrm{vert}}_{\mathbb A',\mathbb B}.
\]
Its object map is
\[
    F^\sharp:
    A'_0\times_{F_0,A_0,p_{\mathbf{Beh}}}
    \mathbf{Beh}_{\mathbb A,\mathbb B,0}
    \longrightarrow
    \mathbf{Beh}_{\mathbb A',\mathbb B,0}.
\]
In generalized-element notation,
\[
    F^\sharp(v',H)
    =
    H\circ F_{/v'},
\]
where
\[
    F_{/v'}:
    \mathbb A'/v'
    \to
    \mathbb A/F_0(v')
\]
is the fibre of \(F_{/}\) at \(v'\).

The domain of \(F^\sharp\) is regarded over
\[
    Z_{\mathbb A'}=A'_0\times B_0
\]
by
\[
    (v',H)\longmapsto(v',q_{\mathbf{Beh}}H).
\]
Since
\[
    F_{/v'}(1_{v'})=1_{F_0(v')},
\]
one has
\[
\begin{aligned}
    q_{\mathbf{Beh}}
    \bigl(
        F^\sharp(v',H)
    \bigr)
    &=
    (H\circ F_{/v'})(1_{v'})
    \\
    &=
    H(1_{F_0(v')})
    \\
    &=
    q_{\mathbf{Beh}}(H).
\end{aligned}
\]
Thus \(F^\sharp\) is a map over \(Z_{\mathbb A'}\).

\begin{lemma}[Reindexing preserves canonical structure]
\label{lem:global-reindexing-preserves-structure}
    With the input-pullback Mealy structure of
    Definition~\ref{def:input-pullback-mealy} on
    \[
        A'_0\times_{F_0,A_0,p_{\mathbf{Beh}}}
        \mathbf{Beh}_{\mathbb A,\mathbb B,0},
    \]
    the map \(F^\sharp\) is a strict semantic homomorphism
    \[
        F^\sharp:
        F^*
        \mathbf{Beh}^{\mathrm{can}}_{\mathbb A,\mathbb B}
        \longrightarrow
        \mathbf{Beh}^{\mathrm{can}}_{\mathbb A',\mathbb B}.
    \]

    Explicitly, for every
    \[
        a':u'\to v'
    \]
    and every behaviour
    \[
        H:\mathbb A/F_0(v')\to\mathbb B,
    \]
    one has
    \[
        \partial_{a'}
        \bigl(
            F^\sharp(v',H)
        \bigr)
        =
        F^\sharp
        \bigl(
            u',
            \partial_{F_1(a')}H
        \bigr),
    \]
    and
    \[
        \omega_{\mathbf{Beh}_{\mathbb A',\mathbb B}}
        \bigl(
            a',
            F^\sharp(v',H)
        \bigr)
        =
        \omega_{\mathbf{Beh}_{\mathbb A,\mathbb B}}
        \bigl(
            F_1(a'),
            H
        \bigr).
    \]

    If \(\mathbb B\) is \(\mathbb T\)-structured, then \(F^\sharp\) also
    preserves all pointwise \(\mathbb T\)-operations.
\end{lemma}

\begin{proof}
    For
    \[
        a':u'\to v',
    \]
    the residual-slice functors satisfy the strict equality
    \[
        F_{/v'}\circ a'_*
        =
        (F_1a')_*\circ F_{/u'}.
    \]
    Indeed, both sides send a residual object
    \[
        r':w'\to u'
    \]
    to
    \[
        F_1(a'\circ r')
        =
        F_1(a')\circ F_1(r'),
    \]
    and act identically on residual triangles.

    Therefore
    \[
    \begin{aligned}
        \partial_{a'}
        \bigl(
            F^\sharp(v',H)
        \bigr)
        &=
        (H\circ F_{/v'})\circ a'_*
        \\
        &=
        H\circ(F_1a')_*\circ F_{/u'}
        \\
        &=
        \bigl(
            \partial_{F_1a'}H
        \bigr)
        \circ F_{/u'}
        \\
        &=
        F^\sharp
        \bigl(
            u',
            \partial_{F_1a'}H
        \bigr).
    \end{aligned}
    \]
    This is the transition-preservation equation.

    The canonical output of \(F^\sharp(v',H)\) along \(a'\) is the arrow
    assigned by
    \[
        H\circ F_{/v'}
    \]
    to the residual morphism
    \[
        a'\to1_{v'}
    \]
    in \(\mathbb A'/v'\). The functor \(F_{/v'}\) sends this morphism to
    \[
        F_1(a')\to1_{F_0(v')}
    \]
    in \(\mathbb A/F_0(v')\). Hence
    \[
    \begin{aligned}
        \omega_{\mathbf{Beh}_{\mathbb A',\mathbb B}}
        \bigl(
            a',
            F^\sharp(v',H)
        \bigr)
        &=
        (H\circ F_{/v'})
        (a'\to1_{v'})
        \\
        &=
        H
        \bigl(
            F_1(a')\to1_{F_0(v')}
        \bigr)
        \\
        &=
        \omega_{\mathbf{Beh}_{\mathbb A,\mathbb B}}
        \bigl(
            F_1(a'),
            H
        \bigr).
    \end{aligned}
    \]
    This is the output-preservation equation.

    Suppose now that \(\mathbb B\) is \(\mathbb T\)-structured and
    \[
        \theta:\mathbb B^n\to\mathbb B
    \]
    is an operation. Since \(F^\sharp\) is defined by precomposition and
    \(\theta_{\mathbf{Beh}}\) is defined by postcomposition,
    \[
    \begin{aligned}
        F^\sharp
        \bigl(
            v',
            \theta_{\mathbf{Beh}}(H_i)_i
        \bigr)
        &=
        \bigl(
            \theta\circ\langle H_i\rangle_i
        \bigr)
        \circ F_{/v'}
        \\
        &=
        \theta
        \circ
        \left\langle
            H_i\circ F_{/v'}
        \right\rangle_i
        \\
        &=
        \theta_{\mathbf{Beh}}
        \bigl(
            F^\sharp(v',H_i)
        \bigr)_i.
    \end{aligned}
    \]
    Thus \(F^\sharp\) preserves the pointwise structured operations.
\end{proof}

\subsection{Input pullback of Mealy morphisms}
\label{subsec:input-pullback-mealy}

\begin{definition}[Input pullback]
\label{def:input-pullback-mealy}
    Let
    \[
        P:\mathbb A\rightsquigarrow\mathbb B
    \]
    be a Mealy morphism. Define
    \[
        F^*P:\mathbb A'\rightsquigarrow\mathbb B
    \]
    as follows. The carrier is
    \[
        F^*P
        :=
        A'_0\times_{F_0,A_0,p_P}P.
    \]
    A state is a pair
    \[
        (v',x)
    \]
    with
    \[
        F_0(v')=p_P(x).
    \]
    The input projection is
    \[
        p_{F^*P}(v',x)=v',
    \]
    and the output projection is
    \[
        q_{F^*P}(v',x)=q_P(x).
    \]
    For an arrow
    \[
        a':u'\to v',
    \]
    define
    \[
        \delta_{F^*P}(a',(v',x))
        =
        (u',\delta_P(F_1a',x)),
    \]
    and
    \[
        \omega_{F^*P}(a',(v',x))
        =
        \omega_P(F_1a',x).
    \]
\end{definition}

The carrier is the pullback
\[
\begin{tikzcd}[column sep=large, row sep=large]
    F^*P
        \arrow[r, "\pi_P"]
        \arrow[d]
        \arrow[dr, phantom, "\lrcorner", very near start]
    &
    P
        \arrow[d, "p_P"]
    \\
    A'_0
        \arrow[r, "F_0"']
    &
    A_0 .
\end{tikzcd}
\]
The transition is induced by the square
\[
\begin{tikzcd}[column sep=huge, row sep=large]
    A'_1\times_{A'_0}F^*P
        \arrow[r, "{F_1\times \pi_P}"]
        \arrow[d, "\delta_{F^*P}"']
    &
    A_1\times_{A_0}P
        \arrow[d, "\delta_P"]
    \\
    F^*P
        \arrow[r, "\pi_P"']
    &
    P ,
\end{tikzcd}
\]
together with the input-object component \(u'\).

For a strict semantic homomorphism
\[
    f:P\to Q
\]
define
\[
    F^*f:F^*P\to F^*Q
\]
by
\[
    F^*f(v',x)=(v',fx).
\]
This assignment is functorial.

\begin{proposition}[Input pullback and semantics]
\label{prop:input-pullback-mealy}
    The data above define a Mealy morphism
    \[
        F^*P:\mathbb A'\rightsquigarrow\mathbb B.
    \]
    Moreover,
    \[
        F^*:\mathsf{Mly}^{\mathrm{str}}_{\mathbb A,\mathbb B}
        \to
        \mathsf{Mly}^{\mathrm{str}}_{\mathbb A',\mathbb B}
    \]
    is a functor, and the semantic map satisfies
    \[
        \beta_{F^*P}
        =
        F^\sharp
        \circ
        \langle p_{F^*P},\,\beta_P\pi_P\rangle.
    \]
    In generalized elements,
    \[
        \beta_{F^*P}(v',x)
        =
        F^\sharp(v',\beta_P(x)).
    \]
\end{proposition}
\begin{proof}
    \leavevmode
    \begin{enumerate}
        \item \textbf{Check typing.}

        If
        \[
            a':u'\to v',
        \]
        then
        \[
            F_1(a'):F_0(u')\to F_0(v').
        \]
        Hence \(\delta_P(F_1a',x)\) is defined whenever
        \[
            p_P(x)=F_0(v'),
        \]
        and satisfies
        \[
            p_P(\delta_P(F_1a',x))=F_0(u').
        \]
        Therefore
        \[
            (u',\delta_P(F_1a',x))
        \]
        lies in
        \[
            A'_0\times_{F_0,A_0,p_P}P.
        \]

        The output
        \[
            \omega_P(F_1a',x)
        \]
        has source
        \[
            q_P(\delta_P(F_1a',x))
        \]
        and target
        \[
            q_P(x),
        \]
        so it has the required typing for \(F^*P\).

        \item \textbf{Verify the Mealy laws.}

        The transition unit law follows from preservation of identities by
        \(F\):
        \[
        \begin{aligned}
            \delta_{F^*P}
            \bigl(
                1_{v'},
                (v',x)
            \bigr)
            &=
            \bigl(
                v',
                \delta_P(F_1(1_{v'}),x)
            \bigr)
            \\
            &=
            \bigl(
                v',
                \delta_P(1_{F_0(v')},x)
            \bigr)
            \\
            &=
            (v',x).
        \end{aligned}
        \]

        Let
        \[
            a':u'\to v',
            \qquad
            b':v'\to w'.
        \]
        Then
        \[
        \begin{aligned}
            \delta_{F^*P}
            \bigl(
                b'\circ a',
                (w',x)
            \bigr)
            &=
            \bigl(
                u',
                \delta_P(F_1(b'\circ a'),x)
            \bigr)
            \\
            &=
            \bigl(
                u',
                \delta_P(F_1b'\circ F_1a',x)
            \bigr)
            \\
            &=
            \bigl(
                u',
                \delta_P
                \bigl(
                    F_1a',
                    \delta_P(F_1b',x)
                \bigr)
            \bigr)
            \\
            &=
            \delta_{F^*P}
            \left(
                a',
                \delta_{F^*P}(b',(w',x))
            \right).
        \end{aligned}
        \]

        The output unit law is
        \[
        \begin{aligned}
            \omega_{F^*P}
            \bigl(
                1_{v'},
                (v',x)
            \bigr)
            &=
            \omega_P(1_{F_0(v')},x)
            \\
            &=
            1_{q_P(x)}.
        \end{aligned}
        \]

        The output multiplication law is
        \[
        \begin{aligned}
            \omega_{F^*P}
            \bigl(
                b'\circ a',
                (w',x)
            \bigr)
            &=
            \omega_P(F_1b'\circ F_1a',x)
            \\
            &=
            \omega_P(F_1b',x)
            \circ
            \omega_P
            \bigl(
                F_1a',
                \delta_P(F_1b',x)
            \bigr)
            \\
            &=
            \omega_{F^*P}(b',(w',x))
            \circ
            \omega_{F^*P}
            \left(
                a',
                \delta_{F^*P}(b',(w',x))
            \right).
        \end{aligned}
        \]

        Hence the displayed data define a Mealy morphism
        \[
            F^*P:\mathbb A'\rightsquigarrow\mathbb B.
        \]

        \item \textbf{Check the action on strict semantic homomorphisms.}

        Let
        \[
            f:P\to Q
        \]
        be strict. Define
        \[
            F^*f(v',x)=(v',fx).
        \]
        This map lies over \(A'_0\) and \(B_0\).

        Transition preservation follows from
        \[
        \begin{aligned}
            F^*f
            \bigl(
                \delta_{F^*P}(a',(v',x))
            \bigr)
            &=
            F^*f
            \bigl(
                u',
                \delta_P(F_1a',x)
            \bigr)
            \\
            &=
            \bigl(
                u',
                f\delta_P(F_1a',x)
            \bigr)
            \\
            &=
            \bigl(
                u',
                \delta_Q(F_1a',fx)
            \bigr)
            \\
            &=
            \delta_{F^*Q}(a',(v',fx)).
        \end{aligned}
        \]

        Output preservation follows from
        \[
            \omega_P(F_1a',x)
            =
            \omega_Q(F_1a',fx).
        \]

        Identities and composition are inherited from \(\mathcal E\), so
        \[
            F^*:
            \mathsf{Mly}^{\mathrm{str}}_{\mathbb A,\mathbb B}
            \to
            \mathsf{Mly}^{\mathrm{str}}_{\mathbb A',\mathbb B}
        \]
        is a functor.

        \item \textbf{Identify residual semantics.}

        Let
        \[
            (v',x)\in F^*P.
        \]
        For a residual object
        \[
            r':u'\to v'
        \]
        of \(\mathbb A'/v'\), execution in \(F^*P\) gives
        \[
            \delta_{F^*P}(r',(v',x))
            =
            \bigl(
                u',
                \delta_P(F_1r',x)
            \bigr).
        \]
        Its output anchor is therefore
        \[
            q_P(\delta_P(F_1r',x))
            =
            \beta_P(x)(F_1r').
        \]

        For a residual triangle
        \[
            w'\xrightarrow{a'}u'\xrightarrow{r'}v',
        \]
        the corresponding residual output is
        \[
            \omega_P
            \bigl(
                F_1a',
                \delta_P(F_1r',x)
            \bigr),
        \]
        which is the arrow assigned by \(\beta_P(x)\) to the image under
        \(F_{/v'}\) of the residual triangle.

        Thus the residual behaviour of \((v',x)\) is the residual behaviour
        of \(x\) precomposed with
        \[
            F_{/v'}:
            \mathbb A'/v'\to\mathbb A/F_0(v').
        \]
        Hence
        \[
            \beta_{F^*P}(v',x)
            =
            F^\sharp(v',\beta_P(x)).
        \]

        \item \textbf{Display the correctly typed semantic factorization.}

        The semantic identity is the commutativity of
        \[
        \begin{tikzcd}[column sep=large, row sep=large]
            F^*P
                \arrow[r,
                    "{\langle
                        p_{F^*P},
                        \beta_P\pi_P
                    \rangle}"]
                \arrow[dr, "\beta_{F^*P}"']
            &
            A'_0
            \times_{F_0,A_0,p_{\mathbf{Beh}}}
            \mathbf{Beh}_{\mathbb A,\mathbb B,0}
                \arrow[d, "F^\sharp"]
            \\
            &
            \mathbf{Beh}_{\mathbb A',\mathbb B,0}.
        \end{tikzcd}
        \]
        Therefore
        \[
            \beta_{F^*P}
            =
            F^\sharp
            \circ
            \left\langle
                p_{F^*P},
                \beta_P\pi_P
            \right\rangle.
        \]
    \end{enumerate}
\end{proof}

\begin{proposition}[Coherence of input pullback and residual reindexing]
\label{prop:input-reindexing-coherence}
    Let
    \[
        1_{\mathbb A}:\mathbb A\to\mathbb A
    \]
    be the identity internal functor. There are canonical natural
    isomorphisms
    \[
        (1_{\mathbb A})^*
        \cong
        1_{
            \mathsf{Mly}^{\mathrm{str}}_{\mathbb A,\mathbb B}
        },
    \]
    and, after the canonical identification of its pulled-back domain,
    \[
        (1_{\mathbb A})^\sharp
        =
        1_{\mathbf{Beh}_{\mathbb A,\mathbb B,0}}.
    \]

    Let
    \[
        \mathbb A''
        \xrightarrow{G}
        \mathbb A'
        \xrightarrow{F}
        \mathbb A
    \]
    be composable internal functors. There is a canonical natural
    isomorphism
    \[
        (F\circ G)^*
        \cong
        G^*F^*.
    \]

    On residual behaviours, after the canonical pullback identification
    \[
    \begin{aligned}
        &
        A''_0
        \times_{F_0G_0,A_0,p_{\mathbf{Beh}}}
        \mathbf{Beh}_{\mathbb A,\mathbb B,0}
        \\
        &\qquad\cong
        A''_0
        \times_{G_0,A'_0}
        \left(
            A'_0
            \times_{F_0,A_0,p_{\mathbf{Beh}}}
            \mathbf{Beh}_{\mathbb A,\mathbb B,0}
        \right),
    \end{aligned}
    \]
    one has
    \[
        (F\circ G)^\sharp
        =
        G^\sharp
        \circ
        \bigl(
            1_{A''_0}\times_{A'_0}F^\sharp
        \bigr).
    \]

    These identity and composition isomorphisms satisfy the usual unit and
    associativity coherence conditions. They are compatible with residual
    semantics: under the canonical isomorphism
    \[
        (F\circ G)^*P\cong G^*F^*P,
    \]
    the two induced semantic maps to
    \[
        \mathbf{Beh}_{\mathbb A'',\mathbb B,0}
    \]
    agree.
\end{proposition}

\begin{proof}
    \leavevmode
    \begin{enumerate}
        \item \textbf{Identity coherence for input pullback.}

        The carrier of
        \[
            (1_{\mathbb A})^*P
        \]
        is
        \[
            A_0\times_{1_{A_0},A_0,p_P}P.
        \]
        The second pullback projection gives the canonical isomorphism
        \[
            A_0\times_{A_0}P
            \overset{\cong}{\longrightarrow}
            P.
        \]
        Under this isomorphism, the pulled-back transition and output maps are
        exactly those of \(P\), because
        \[
            (1_{\mathbb A})_1=1_{A_1}.
        \]
        These isomorphisms are natural in strict semantic homomorphisms, giving
        \[
            (1_{\mathbb A})^*
            \cong
            1.
        \]

        \item \textbf{Composition coherence for input pullback.}

        The carrier of \((F\circ G)^*P\) is
        \[
            A''_0
            \times_{F_0G_0,A_0,p_P}
            P.
        \]
        The carrier of \(G^*F^*P\) is
        \[
            A''_0
            \times_{G_0,A'_0}
            \left(
                A'_0
                \times_{F_0,A_0,p_P}
                P
            \right).
        \]
        The associativity universal property of pullbacks gives a canonical
        isomorphism between these two carriers.

        On an input arrow
        \[
            a'':u''\to v'',
        \]
        both transition structures apply the transition of \(P\) to
        \[
            F_1G_1(a'')
            =
            (F\circ G)_1(a'').
        \]
        Their outputs are likewise both
        \[
            \omega_P(F_1G_1(a''),x).
        \]
        Hence the carrier isomorphism is a strict Mealy isomorphism, natural
        in \(P\):
        \[
            (F\circ G)^*
            \cong
            G^*F^*.
        \]

        \item \textbf{Identity coherence for residual reindexing.}

        For every \(v\in A_0\), the residual-slice functor induced by the
        identity is
        \[
            (1_{\mathbb A})_{/v}
            =
            1_{\mathbb A/v}.
        \]
        Therefore
        \[
            (1_{\mathbb A})^\sharp(v,H)
            =
            H\circ1_{\mathbb A/v}
            =
            H.
        \]
        After identifying
        \[
            A_0\times_{A_0}\mathbf{Beh}_{\mathbb A,\mathbb B,0}
            \cong
            \mathbf{Beh}_{\mathbb A,\mathbb B,0},
        \]
        one has
        \[
            (1_{\mathbb A})^\sharp
            =
            1_{\mathbf{Beh}_{\mathbb A,\mathbb B,0}}.
        \]

        \item \textbf{Composition coherence for residual reindexing.}

        Let
        \[
            v''\in A''_0.
        \]
        The residual-slice functor of the composite satisfies
        \[
            (F\circ G)_{/v''}
            =
            F_{/G_0(v'')}
            \circ
            G_{/v''}.
        \]
        This equality holds on residual objects because
        \[
            (F\circ G)_1(r'')
            =
            F_1(G_1(r'')),
        \]
        and it holds on residual arrows by functoriality of \(F\) and \(G\).

        Hence, for
        \[
            H:
            \mathbb A/F_0G_0(v'')
            \to
            \mathbb B,
        \]
        \[
        \begin{aligned}
            (F\circ G)^\sharp(v'',H)
            &=
            H\circ(F\circ G)_{/v''}
            \\
            &=
            H\circ F_{/G_0(v'')}
            \circ G_{/v''}
            \\
            &=
            G^\sharp
            \bigl(
                v'',
                F^\sharp(G_0(v''),H)
            \bigr).
        \end{aligned}
        \]
        This is the asserted composition equation.

        \item \textbf{Verify coherence and semantic compatibility.}

        The unit and associativity coherence conditions for \(F^*\) are the
        canonical coherence conditions of iterated pullbacks. The
        corresponding conditions for \(F^\sharp\) follow from strict
        associativity and unitality of functor composition.

        Finally, Proposition~\ref{prop:input-pullback-mealy} identifies
        residual semantics by precomposition with the relevant residual-slice
        functor. Since
        \[
            (F\circ G)_{/v''}
            =
            F_{/G_0(v'')}
            \circ
            G_{/v''},
        \]
        the semantic map obtained directly from \((F\circ G)^*P\) agrees with
        the semantic map obtained through \(G^*F^*P\) under the canonical
        pullback isomorphism.
    \end{enumerate}
\end{proof}

\begin{lemma}[Input pullback and semantic images]
\label{lem:input-pullback-semantic-images}
    Let
    \[
        F:\mathbb A'\to\mathbb A
    \]
    be an internal functor and let
    \[
        P:\mathbb A\rightsquigarrow\mathbb B
    \]
    be a Mealy morphism.

    Write the semantic image factorization of \(P\) as
    \[
        P
        \xrightarrow{e_P}
        \operatorname{Beh}(P)
        \xrightarrow{m_P}
        \mathbf{Beh}_{\mathbb A,\mathbb B,0}.
    \]
    Define
    \[
        F^\sharp_P:
        A'_0
        \times_{F_0,A_0,p_{\operatorname{Beh}(P)}}
        \operatorname{Beh}(P)
        \to
        \mathbf{Beh}_{\mathbb A',\mathbb B,0}
    \]
    to be the restriction
    \[
        F^\sharp_P
        :=
        F^\sharp
        \circ
        (1_{A'_0}\times m_P).
    \]

    Then
    \[
        \operatorname{Im}(F^\sharp_P)
        =
        \operatorname{Beh}(F^*P)
    \]
    as subobjects of
    \[
        \mathbf{Beh}_{\mathbb A',\mathbb B,0}
    \]
    in the slice
    \[
        \mathcal E/(A'_0\times B_0).
    \]
\end{lemma}

\begin{proof}
    Put
    \[
        Z_{\mathbb A}:=A_0\times B_0,
        \qquad
        Z_{\mathbb A'}:=A'_0\times B_0,
    \]
    and let
    \[
        \overline F
        :=
        F_0\times1_{B_0}:
        Z_{\mathbb A'}\to Z_{\mathbb A}.
    \]

    The semantic map of \(P\) has regular image factorization
    \[
    \begin{tikzcd}[column sep=large, row sep=large]
        P
            \arrow[rr, "\beta_P"]
            \arrow[dr, two heads, "e_P"']
        &&
        \mathbf{Beh}_{\mathbb A,\mathbb B,0}
        \\
        &
        \operatorname{Beh}(P)
            \arrow[ur, hook, "m_P"']
        &
    \end{tikzcd}
    \]
    in \(\mathcal E/Z_{\mathbb A}\).

    Pulling this factorization back along
    \[
        \overline F:
        Z_{\mathbb A'}\to Z_{\mathbb A}
    \]
    gives a regular epimorphism
    \[
        \overline F^*e_P:
        \overline F^*P
        \twoheadrightarrow
        \overline F^*\operatorname{Beh}(P)
    \]
    in \(\mathcal E/Z_{\mathbb A'}\).

    There are canonical pullback isomorphisms
    \[
    \begin{aligned}
        \overline F^*P
        &=
        Z_{\mathbb A'}
        \times_{Z_{\mathbb A}}
        P
        \\
        &\cong
        A'_0
        \times_{F_0,A_0,p_P}
        P
        \\
        &=
        F^*P,
    \end{aligned}
    \]
    and
    \[
    \begin{aligned}
        \overline F^*\operatorname{Beh}(P)
        &=
        Z_{\mathbb A'}
        \times_{Z_{\mathbb A}}
        \operatorname{Beh}(P)
        \\
        &\cong
        A'_0
        \times_{F_0,A_0,p_{\operatorname{Beh}(P)}}
        \operatorname{Beh}(P).
    \end{aligned}
    \]
    The \(B_0\)-coordinate in each pullback is forced by the output projection
    of the corresponding state, which gives these canonical
    identifications.

    Under these identifications, write
    \[
        \widetilde e_P:
        F^*P
        \twoheadrightarrow
        A'_0
        \times_{F_0,A_0,p_{\operatorname{Beh}(P)}}
        \operatorname{Beh}(P)
    \]
    for the pullback of \(e_P\). In generalized-element notation,
    \[
        \widetilde e_P(v',x)
        =
        (v',e_Px).
    \]
    Since regular epimorphisms are stable under pullback,
    \(\widetilde e_P\) is a regular epimorphism in
    \[
        \mathcal E/Z_{\mathbb A'}.
    \]

    By
    Proposition~\ref{prop:input-pullback-mealy},
    \[
        \beta_{F^*P}
        =
        F^\sharp
        \circ
        \left\langle
            p_{F^*P},
            \beta_P\pi_P
        \right\rangle.
    \]
    Since
    \[
        \beta_P=m_Pe_P,
    \]
    the map into the pullback behaviour object factors as
    \[
    \begin{aligned}
        \left\langle
            p_{F^*P},
            \beta_P\pi_P
        \right\rangle
        &=
        (1_{A'_0}\times m_P)
        \widetilde e_P.
    \end{aligned}
    \]
    Consequently,
    \[
    \begin{aligned}
        \beta_{F^*P}
        &=
        F^\sharp
        (1_{A'_0}\times m_P)
        \widetilde e_P
        \\
        &=
        F^\sharp_P
        \widetilde e_P.
    \end{aligned}
    \]

    Thus the semantic map of \(F^*P\) is the composite
    \[
    \begin{tikzcd}[column sep=large, row sep=large]
        F^*P
            \arrow[r, two heads, "\widetilde e_P"]
            \arrow[rr, bend left=18, "\beta_{F^*P}"]
        &
        A'_0
        \times_{F_0,A_0,p_{\operatorname{Beh}(P)}}
        \operatorname{Beh}(P)
            \arrow[r, "F^\sharp_P"]
        &
        \mathbf{Beh}_{\mathbb A',\mathbb B,0}.
    \end{tikzcd}
    \]

    Precomposition with a regular epimorphism does not change a regular
    image. Therefore
    \[
        \operatorname{Im}(\beta_{F^*P})
        =
        \operatorname{Im}(F^\sharp_P).
    \]
    By definition,
    \[
        \operatorname{Im}(\beta_{F^*P})
        =
        \operatorname{Beh}(F^*P).
    \]
    Hence
    \[
        \operatorname{Im}(F^\sharp_P)
        =
        \operatorname{Beh}(F^*P)
    \]
    as subobjects of
    \[
        \mathbf{Beh}_{\mathbb A',\mathbb B,0}
    \]
    in \(\mathcal E/Z_{\mathbb A'}\).
\end{proof}

\section{Global coequational theories}
\label{sec:global-coequational-theories}

\subsection{Global theories and global models}
\label{subsec:global-theories-and-models}

\begin{definition}[Global coequational theory]
\label{def:global-mealy-coequational-theory}
    A \textbf{global coequational theory} consists of a local coequational
    theory
    \[
        V_{\mathbb A}\hookrightarrow
        \mathbf{Beh}^{\mathrm{can}}_{\mathbb A,\mathbb B}
    \]
    for each \(\mathbb A\in\mathsf{Inp}\), such that for every internal functor
    \[
        F:\mathbb A'\to\mathbb A
    \]
    the reindexing map
    \[
        F^\sharp:
        A'_0\times_{F_0,A_0,p_{\mathbf{Beh}}}
        \mathbf{Beh}_{\mathbb A,\mathbb B,0}
        \to
        \mathbf{Beh}_{\mathbb A',\mathbb B,0}
    \]
    sends
    \[
        A'_0\times_{F_0,A_0,p_{V_{\mathbb A}}}V_{\mathbb A}
    \]
    into
    \[
        V_{\mathbb A'}.
    \]
    If \(\mathbb B\) is \(\mathbb T\)-structured, a \textbf{global
    \(\mathbb T\)-coequational theory} is such a family for which each
    \(V_{\mathbb A}\) is a local \(\mathbb T\)-coequational theory.
\end{definition}

The global reindexing condition is displayed by
\[
\begin{tikzcd}[column sep=large, row sep=large]
    A'_0\times_{A_0}V_{\mathbb A}
        \arrow[rr, "{F^\sharp|_{V_{\mathbb A}}}"]
        \arrow[dr, dashed]
    &&
    \mathbf{Beh}_{\mathbb A',\mathbb B,0}
    \\
    &
    V_{\mathbb A'}
        \arrow[ur, hook]
    &
\end{tikzcd}
\]
where the fibre product is taken along
\[
    F_0:A'_0\to A_0
    \qquad\text{and}\qquad
    p_{V_{\mathbb A}}:V_{\mathbb A}\to A_0.
\]
Including the carrier slices, the condition is
\[
\begin{tikzcd}[column sep=large, row sep=large]
    A'_0\times_{A_0}V_{\mathbb A}
        \arrow[r]
        \arrow[d]
        \arrow[dr, phantom, "\lrcorner", very near start]
    &
    V_{\mathbb A}
        \arrow[d, hook]
    \\
    A'_0\times_{A_0}\mathbf{Beh}_{\mathbb A,\mathbb B,0}
        \arrow[r, "F^\sharp"']
        \arrow[d]
    &
    \mathbf{Beh}_{\mathbb A',\mathbb B,0}
    \\
    Z_{\mathbb A'}
        \arrow[ur, phantom, "\text{over }Z_{\mathbb A'}" description]
    &
\end{tikzcd}
\]
together with factorization through \(V_{\mathbb A'}\).

\begin{definition}[Global models]
\label{def:global-models}
    Let
    \[
        V=(V_{\mathbb A})_{\mathbb A\in\mathsf{Inp}}
    \]
    be a global coequational theory. Define
    \[
        \mathsf{Mod}^{g}(V)_{\mathbb A}
        =
        \{P\in\mathsf{Mly}^{\mathrm{str}}_{\mathbb A,\mathbb B}
          \mid
          \beta_P\text{ factors through }V_{\mathbb A}\}.
    \]
\end{definition}

\subsection{Global closure}
\label{subsec:global-coequational-closure}

\begin{theorem}[Global coequational closure]
\label{thm:global-coequational-closure}
    Let
    \[
        V=(V_{\mathbb A})_{\mathbb A\in\mathsf{Inp}}
    \]
    be a global coequational theory. Then
    \[
        \mathsf{Mod}^{g}(V)
    \]
    is closed under strict homomorphic preimages, regular-epimorphic quotients,
    small coproducts, and input pullback along functors in \(\mathsf{Inp}\).

    If \(V\) is a global \(\mathbb T\)-coequational theory, then
    \[
        \mathsf{Mod}^{g}(V)
    \]
    is also closed under the operations induced by \(\mathbb T\).
\end{theorem}

\begin{proof}
    \leavevmode
    \begin{enumerate}
        \item \textbf{Apply local closure pointwise.}
        Closure under strict homomorphic preimages, regular-epimorphic
        quotients, and small coproducts is
        Theorem~\ref{thm:local-coequational-closure} applied to each
        \(V_{\mathbb A}\). In the structured case, closure under
        \(\mathbb T\)-operations is Theorem~\ref{thm:T-coequational-closure}
        applied pointwise.

        \item \textbf{Check input pullback.}
        Let
        \[
            F:\mathbb A'\to\mathbb A
        \]
        and let
        \[
            P\in\mathsf{Mod}^{g}(V)_{\mathbb A}.
        \]
        Choose a factorization
        \[
            \bar\beta_P:P\to V_{\mathbb A}
        \]
        of \(\beta_P\). Proposition~\ref{prop:input-pullback-mealy}
        gives
        \[
            \beta_{F^*P}
            =
            F^\sharp
            \circ
            \langle p_{F^*P},\,\beta_P\pi_P\rangle.
        \]
        Since \(V\) is global, \(F^\sharp\) sends the pullback of
        \(V_{\mathbb A}\) into \(V_{\mathbb A'}\). Therefore
        \[
            \beta_{F^*P}
        \]
        factors through \(V_{\mathbb A'}\), so
        \[
            F^*P\in\mathsf{Mod}^{g}(V)_{\mathbb A'}.
        \]
        The factorization is
        \[
        \begin{tikzcd}[column sep=large, row sep=large]
            F^*P
                \arrow[r, "{\langle p_{F^*P},\,\bar\beta_P\pi_P\rangle}"]
                \arrow[rrr, bend left=18, "\beta_{F^*P}"]
            &
            A'_0\times_{A_0}V_{\mathbb A}
                \arrow[r, dashed]
            &
            V_{\mathbb A'}
                \arrow[r, hook]
            &
            \mathbf{Beh}_{\mathbb A',\mathbb B,0}.
        \end{tikzcd}
        \]
    \end{enumerate}
\end{proof}

\section{The global coequational Birkhoff theorem}
\label{sec:global-coequational-birkhoff}

\subsection{The theorem}
\label{subsec:global-coequational-birkhoff-theorem}

\begin{theorem}[Global coequational Birkhoff theorem]
\label{thm:global-coequational-birkhoff}
    Let
    \[
        \mathcal C=(\mathcal C_{\mathbb A})_{\mathbb A\in\mathsf{Inp}}
    \]
    be a replete global family of full subcategories
    \[
        \mathcal C_{\mathbb A}
        \subseteq
        \mathsf{Mly}^{\mathrm{str}}_{\mathbb A,\mathbb B}.
    \]
    Then \(\mathcal C\) is of the form
    \[
        \mathcal C_{\mathbb A}
        =
        \mathsf{Mod}^{g}(V)_{\mathbb A}
    \]
    for a global coequational theory
    \[
        V=(V_{\mathbb A})_{\mathbb A\in\mathsf{Inp}}
    \]
    if and only if each \(\mathcal C_{\mathbb A}\) is closed under:
    \begin{enumerate}
        \item strict homomorphic preimages;
        \item regular-epimorphic quotients;
        \item small coproducts;
    \end{enumerate}
    and the family \(\mathcal C\) is closed under input pullback along functors
    in \(\mathsf{Inp}\).

    If \(\mathbb B\) is \(\mathbb T\)-structured, then global
    \(\mathbb T\)-coequational theories correspond to such families that are
    also closed under the operations induced by \(\mathbb T\).
\end{theorem}

\begin{proof}
    \leavevmode
    \begin{enumerate}
        \item \textbf{Prove necessity.}

        Suppose
        \[
            \mathcal C=\mathsf{Mod}^{g}(V)
        \]
        for a global coequational theory \(V\). The required closure
        properties are
        Theorem~\ref{thm:global-coequational-closure}.

        \item \textbf{Construct the local representing theories.}

        Conversely, suppose that each \(\mathcal C_{\mathbb A}\) is closed
        under strict homomorphic preimages, regular-epimorphic quotients, and
        small coproducts, and that the family \(\mathcal C\) is closed under
        input pullback.

        For every
        \[
            \mathbb A\in\mathsf{Inp},
        \]
        define
        \[
            \mathcal S_{\mathcal C,\mathbb A}
            :=
            \{
                S\hookrightarrow
                \mathbf{Beh}^{\mathrm{can}}_{\mathbb A,\mathbb B}
                \mid
                S=\operatorname{Beh}(P)
                \text{ for some }P\in\mathcal C_{\mathbb A}
            \}.
        \]
        Choose a small representative set for the distinct subobjects in
        \(\mathcal S_{\mathcal C,\mathbb A}\), and put
        \[
            V_{\mathcal C,\mathbb A}
            :=
            \bigvee_{S\in\mathcal S_{\mathcal C,\mathbb A}}S
            \hookrightarrow
            \mathbf{Beh}^{\mathrm{can}}_{\mathbb A,\mathbb B}.
        \]

        By Theorem~\ref{thm:local-coequational-birkhoff},
        \(V_{\mathcal C,\mathbb A}\) is a local coequational theory and
        \[
            \mathcal C_{\mathbb A}
            =
            \mathsf{Mod}(V_{\mathcal C,\mathbb A}).
        \]
        By
        Lemma~\ref{lem:representing-theory-belongs-to-class},
        one also has
        \[
            V_{\mathcal C,\mathbb A}
            \in
            \mathcal C_{\mathbb A}.
        \]

        \item \textbf{Use input-pullback closure on the representing
        theory.}

        Let
        \[
            F:\mathbb A'\to\mathbb A
        \]
        be an internal functor in \(\mathsf{Inp}\).

        Since
        \[
            V_{\mathcal C,\mathbb A}
            \in
            \mathcal C_{\mathbb A}
        \]
        and the family \(\mathcal C\) is closed under input pullback,
        \[
            F^*V_{\mathcal C,\mathbb A}
            \in
            \mathcal C_{\mathbb A'}.
        \]
        Using
        \[
            \mathcal C_{\mathbb A'}
            =
            \mathsf{Mod}(V_{\mathcal C,\mathbb A'}),
        \]
        the object \(F^*V_{\mathcal C,\mathbb A}\) models
        \(V_{\mathcal C,\mathbb A'}\). Hence there is a unique factorization
        \[
            \varphi_F:
            F^*V_{\mathcal C,\mathbb A}
            \longrightarrow
            V_{\mathcal C,\mathbb A'}
        \]
        satisfying
        \[
            j_{V_{\mathcal C,\mathbb A'}}
            \varphi_F
            =
            \beta_{F^*V_{\mathcal C,\mathbb A}}.
        \]
        By Lemma~\ref{lem:model-factorizations-strict},
        \(\varphi_F\) is a strict semantic homomorphism.

        \item \textbf{Identify the model factorization with the required
        restriction of \(F^\sharp\).}

        The carrier of the input pullback is
        \[
            F^*V_{\mathcal C,\mathbb A}
            =
            A'_0
            \times_{F_0,A_0,p_{V_{\mathcal C,\mathbb A}}}
            V_{\mathcal C,\mathbb A}.
        \]
        Write
        \[
            \pi_V:
            F^*V_{\mathcal C,\mathbb A}
            \to
            V_{\mathcal C,\mathbb A}
        \]
        for its second pullback projection.

        Proposition~\ref{prop:input-pullback-mealy} gives
        \[
        \begin{aligned}
            \beta_{F^*V_{\mathcal C,\mathbb A}}
            &=
            F^\sharp
            \circ
            \left\langle
                p_{F^*V_{\mathcal C,\mathbb A}},
                \beta_{V_{\mathcal C,\mathbb A}}\pi_V
            \right\rangle.
        \end{aligned}
        \]
        By
        Lemma~\ref{lem:semantic-map-of-a-restricted-canonical-subobject},
        \[
            \beta_{V_{\mathcal C,\mathbb A}}
            =
            j_{V_{\mathcal C,\mathbb A}}.
        \]
        Therefore
        \[
        \begin{aligned}
            \beta_{F^*V_{\mathcal C,\mathbb A}}
            &=
            F^\sharp
            \circ
            \left\langle
                p_{F^*V_{\mathcal C,\mathbb A}},
                j_{V_{\mathcal C,\mathbb A}}\pi_V
            \right\rangle.
        \end{aligned}
        \]

        The map
        \[
            \left\langle
                p_{F^*V_{\mathcal C,\mathbb A}},
                j_{V_{\mathcal C,\mathbb A}}\pi_V
            \right\rangle
        \]
        is precisely the canonical pulled-back inclusion
        \[
        \begin{aligned}
            A'_0
            \times_{F_0,A_0,p_{V_{\mathcal C,\mathbb A}}}
            V_{\mathcal C,\mathbb A}
            &\longrightarrow
            A'_0
            \times_{F_0,A_0,p_{\mathbf{Beh}}}
            \mathbf{Beh}_{\mathbb A,\mathbb B,0}.
        \end{aligned}
        \]
        Hence the factorization equation for \(\varphi_F\) becomes
        \[
        \begin{aligned}
            j_{V_{\mathcal C,\mathbb A'}}
            \varphi_F
            &=
            F^\sharp
            \circ
            \left\langle
                p_{F^*V_{\mathcal C,\mathbb A}},
                j_{V_{\mathcal C,\mathbb A}}\pi_V
            \right\rangle.
        \end{aligned}
        \]

        This is exactly the required global reindexing factorization:
        \[
        \begin{tikzcd}[column sep=large, row sep=large]
            A'_0
            \times_{F_0,A_0,p_{V_{\mathcal C,\mathbb A}}}
            V_{\mathcal C,\mathbb A}
                \arrow[rr,
                    "{F^\sharp
                    \langle
                        p,\,
                        j_{V_{\mathcal C,\mathbb A}}\pi_V
                    \rangle}"]
                \arrow[dr, "\varphi_F"']
            &&
            \mathbf{Beh}_{\mathbb A',\mathbb B,0}
            \\
            &
            V_{\mathcal C,\mathbb A'}
                \arrow[ur, hook,
                    "j_{V_{\mathcal C,\mathbb A'}}"']
            &
        \end{tikzcd}
        \]
        Thus \(F^\sharp\) sends the pullback of
        \(V_{\mathcal C,\mathbb A}\) into
        \(V_{\mathcal C,\mathbb A'}\).

        Since \(F\) was arbitrary, the family
        \[
            V_{\mathcal C}
            :=
            (V_{\mathcal C,\mathbb A})_{\mathbb A\in\mathsf{Inp}}
        \]
        is a global coequational theory.

        \item \textbf{Handle the structured case.}

        Suppose now that \(\mathbb B\) is \(\mathbb T\)-structured and that
        every \(\mathcal C_{\mathbb A}\) is also closed under the operations
        induced by \(\mathbb T\).

        By Theorem~\ref{thm:T-coequational-birkhoff}, applied separately to
        every input category \(\mathbb A\), each
        \[
            V_{\mathcal C,\mathbb A}
        \]
        is a local \(\mathbb T\)-coequational theory. The preceding argument
        has already shown that these local theories satisfy the global
        reindexing condition. Therefore
        \[
            V_{\mathcal C}
        \]
        is a global \(\mathbb T\)-coequational theory.

        \item \textbf{Conclude.}

        For every
        \[
            \mathbb A\in\mathsf{Inp},
        \]
        the local Birkhoff theorem gives
        \[
            \mathcal C_{\mathbb A}
            =
            \mathsf{Mod}(V_{\mathcal C,\mathbb A}).
        \]
        Since \(V_{\mathcal C}\) is a global theory, this equality is precisely
        \[
            \mathcal C_{\mathbb A}
            =
            \mathsf{Mod}^{g}(V_{\mathcal C})_{\mathbb A}.
        \]
        Hence
        \[
            \mathcal C
            =
            \mathsf{Mod}^{g}(V_{\mathcal C}).
        \]
    \end{enumerate}
\end{proof}

\begin{corollary}[Uniqueness of the defining global theory]
\label{cor:uniqueness-defining-global-coequational-theory}
    Let
    \[
        V=(V_{\mathbb A})_{\mathbb A\in\mathsf{Inp}}
    \]
    be a global coequational theory. Then, for every
    \[
        \mathbb A\in\mathsf{Inp},
    \]
    one has
    \[
        V_{\mathbb A}
        =
        \bigvee_{
            P\in\mathsf{Mod}^{g}(V)_{\mathbb A}
        }
        \operatorname{Beh}(P).
    \]

    Consequently a globally coequationally definable family of model classes
    has a unique defining global coequational theory. The same statement holds
    for global \(\mathbb T\)-coequational theories.
\end{corollary}

\begin{proof}
    For each fixed input category \(\mathbb A\),
    \[
        \mathsf{Mod}^{g}(V)_{\mathbb A}
        =
        \mathsf{Mod}(V_{\mathbb A}).
    \]
    Therefore
    Corollary~\ref{cor:uniqueness-defining-local-coequational-theory} gives
    \[
        V_{\mathbb A}
        =
        \bigvee_{
            P\in\mathsf{Mod}^{g}(V)_{\mathbb A}
        }
        \operatorname{Beh}(P).
    \]
    Thus every component \(V_{\mathbb A}\) is determined uniquely by the
    corresponding model class. Hence the entire global family \(V\) is
    uniquely determined.

    The structured case follows pointwise in the same way.
\end{proof}

\section{One-object specialization}
\label{sec:coeq-one-object-specialization}

This section records the one-object form of the preceding constructions. Let
the input category be determined by a monoid object \(M\). In the purely
one-object output case, let the output category be determined by a monoid
object \(N\). We keep the standing convention
\[
    b\circ a=ba .
\]
Thus a Mealy morphism is a right \(M\)-object \(P\), written
\[
    x\cdot m:=\delta(m,x),
\]
together with an \(N\)-valued output cocycle
\[
    \omega:M\times P\to N
\]
satisfying
\[
    \omega(1,x)=1,
    \qquad
    \omega(mn,x)=\omega(m,x)\omega(n,x\cdot m).
\]

The residual category of \(M\) is the right Cayley residual category. Its
objects are elements \(r\in M\), and a residual arrow labelled \(m\) has the
form
\[
    rm\longrightarrow r .
\]
A residual behaviour is therefore an \(N\)-valued cocycle
\[
    \lambda:M\times M\to N
\]
satisfying
\[
    \lambda(r,1)=1,
    \qquad
    \lambda(r,mn)=\lambda(r,m)\lambda(rm,n).
\]
The canonical derivative is left translation in the residual-history variable:
\[
    (\partial_s\lambda)(r,m)=\lambda(sr,m),
\]
and the canonical output on input \(s\) is
\[
    \lambda(1,s).
\]

Hence a local coequational theory in the one-object monoid-valued case is a
subobject
\[
    V\hookrightarrow \mathbf{Beh}_0
\]
of residual cocycles closed under all derivatives:
\[
    \lambda\in V
    \quad\Longrightarrow\quad
    \partial_s\lambda\in V
    \qquad(s\in M).
\]
A Mealy morphism \(P\) models \(V\) precisely when every residual cocycle
generated by a state of \(P\) lies in \(V\). The residual cocycle generated by
\(x\in P\) is
\[
    \lambda_x(r,m)=\omega(m,x\cdot r).
\]
Thus
\[
    P\models V
    \quad\Longleftrightarrow\quad
    \lambda_x\in V
    \text{ for every }x\in P .
\]

The semantic image of \(P\) is the image of the state-to-behaviour map
\[
    P\to \mathbf{Beh}_0,
    \qquad
    x\longmapsto \lambda_x .
\]
In \(\mathbf{Set}\), this is simply the set
\[
    \operatorname{Beh}(P)
    =
    \{\lambda_x\mid x\in P\}
    \subseteq \mathbf{Beh}_0 .
\]
The semantic quotient
\[
    P\twoheadrightarrow \operatorname{Beh}(P)
\]
identifies precisely those states with the same residual cocycle:
\[
    x\sim y
    \quad\Longleftrightarrow\quad
    \lambda_x=\lambda_y .
\]

For a fixed one-object input monoid \(M\), the local coequational Birkhoff
theorem specializes as follows. Derivative-closed subobjects
\[
    V\hookrightarrow\mathbf{Beh}_0
\]
correspond to replete full classes of strict Mealy morphisms closed under
strict homomorphic preimages, regular-epimorphic quotients, and small
coproducts. The theory generated by such a class \(\mathcal C\) is
\[
    V_{\mathcal C}
    =
    \bigvee_{P\in\mathcal C}\operatorname{Beh}(P).
\]
In \(\mathbf{Set}\), this join is the union of the residual behaviours
generated by the objects of \(\mathcal C\).

The Boolean language case is obtained by taking
\[
    \mathcal E=\mathbf{Set},
    \qquad
    M=\Sigma^*,
    \qquad
    \mathbb B=2_{\mathrm{ch}},
\]
where \(2_{\mathrm{ch}}\) is the chaotic category on the Boolean set \(2\).
Since \(2_{\mathrm{ch}}\) has a unique arrow between any two Boolean values,
the observable data of a state is its Boolean output anchor. A residual
behaviour is therefore just a language
\[
    L:\Sigma^*\to 2 .
\]
The derivative convention becomes
\[
    (\partial_uL)(w)=L(uw).
\]
Thus a local coequational theory over the alphabet \(\Sigma\) is a collection
\[
    V\subseteq 2^{\Sigma^*}
\]
closed under left derivatives:
\[
    L\in V
    \quad\Longrightarrow\quad
    \partial_uL\in V
    \qquad(u\in\Sigma^*).
\]
A Boolean recognizer \(P\) models \(V\) precisely when all residual languages
generated by its states lie in \(V\):
\[
    \{\beta_P(x)\mid x\in P\}\subseteq V .
\]
Its semantic image is
\[
    \operatorname{Beh}(P)
    =
    \{\beta_P(x)\mid x\in P\}
    \subseteq
    2^{\Sigma^*},
\]
and the semantic quotient is the usual Nerode-style quotient identifying
states with the same residual language.

If the Boolean operations
\[
    \wedge,\vee,\neg,0,1
\]
are included as structure on \(2_{\mathrm{ch}}\), then a local Boolean
coequational theory is a Boolean subalgebra of \(2^{\Sigma^*}\) closed under
residual derivatives. Finite recognizability is not part of this
coequational closure condition by itself; it enters only through an additional
finite-state restriction or by working inside the regular languages. Thus the
derivative-closed Boolean structure is separated from the finite-state
regularity condition, as in the classical distinction between all languages and
regular languages
\cite{Eilenberg1974,HopcroftMotwaniUllman2006,Brzozowski1964}.

The global input-reindexing construction specializes to inverse morphic
preimage. A functor between one-object free-monoid input categories is a monoid
homomorphism
\[
    h:\Gamma^*\to\Sigma^* .
\]
For a language
\[
    L:\Sigma^*\to 2,
\]
the reindexed behaviour is
\[
    h^{-1}L=L\circ h,
    \qquad
    h^{-1}L(w)=L(h(w)).
\]
Therefore a global Boolean coequational theory assigns to each alphabet
\(\Sigma\) a derivative-closed collection
\[
    V_\Sigma\subseteq 2^{\Sigma^*}
\]
such that
\[
    L\in V_\Sigma
    \quad\Longrightarrow\quad
    h^{-1}L\in V_\Gamma
\]
for every homomorphism
\[
    h:\Gamma^*\to\Sigma^*.
\]
If the theory is Boolean, each \(V_\Sigma\) is also closed under the pointwise
Boolean operations. The global coequational Birkhoff theorem then recovers the
coequational inverse-image closure condition familiar from Eilenberg-type
correspondences: the associated model classes are exactly the recognizers whose
generated residual languages lie in the corresponding \(V_\Sigma\).

\section{Summary}
\label{sec:coequations-birkhoff-summary}

This chapter has reformulated residual semantics in terms of coequational
subobjects of the terminal behaviour Mealy morphism.

\[
\begin{array}{c|l}
\text{construction} & \text{role} \\
\hline
\mathbf{Beh}^{\mathrm{can}}_{\mathbb A,\mathbb B}
&
\text{terminal residual semantic Mealy morphism}
\\
V\hookrightarrow \mathbf{Beh}^{\mathrm{can}}
&
\text{local coequational theory}
\\
\beta_P:P\to\mathbf{Beh}^{\mathrm{can}}
&
\text{unique semantic map to the terminal object}
\\
\operatorname{Beh}(P)
&
\text{regular semantic image of a Mealy morphism}
\\
\mathsf{Mod}(V)
&
\text{Mealy morphisms whose semantics land in }V
\\
V_{\mathcal C}
=
\bigvee_{P\in\mathcal C}\operatorname{Beh}(P)
&
\text{coequational theory generated by a model class}
\\
\theta_{\mathbf{Beh}}
&
\text{pointwise structured operation on behaviours}
\\
F^\sharp
&
\text{input reindexing of residual behaviours}
\\
F^*P
&
\text{input pullback of a Mealy morphism}
\end{array}
\]

The local coequational Birkhoff theorem identifies local coequational theories
with replete full classes of strict Mealy morphisms closed under strict
homomorphic preimages, regular-epimorphic quotients, and small coproducts. The
structured version adds closure under pointwise algebraic operations induced by
a \(\mathbb T\)-structured output category. The global version adds stability
under input pullback, equivalently under residual behaviour reindexing.

The next chapter extracts the algebraic side of this coequational semantics.
For a local coequational theory \(V\), the restricted canonical object on \(V\)
determines a residual transition--output action of
\[
    \mathbb A^{\mathrm{op}},
\]
reflecting the target-to-source processing convention. The regular image of
this action is the transition--output quotient associated to \(V\). For a
specification \(K\), applying this construction to the derivative theory
\[
    \operatorname{Der}_0(K)
\]
produces the syntactic transition--output category.
